\documentclass[10pt,a4paper]{article}
\usepackage[utf8]{inputenc}
\usepackage[T1]{fontenc}
\usepackage[margin=18mm]{geometry}
\usepackage{longtable}
\usepackage{hyperref}
\usepackage{enumitem}
\setlist[itemize]{noitemsep, topsep=2pt}
\setlength{\parindent}{0pt}
\begin{document}

\section*{Non-Verbal \& Multimodal Discourse Analysis Worksheet}

\textit{Qualitative analysis of pauses, trailing sentences, prosody/intonation, vocalisations (coughs, sighs, scoffs), and kinesics (hand gestures, facial expressions, posture) in political news discourse}\par\smallskip

\bigskip\hrule\bigskip

\subsection*{Guidelines: Five non-verbal \& paralinguistic features to look for in every event}

\begin{longtable}{|p{0.23\textwidth}|p{0.23\textwidth}|p{0.23\textwidth}|p{0.23\textwidth}|}
\hline
Feature & What to look for & Illustrative example (from sample) & Questions to answer \\ \hline \endhead
\textbf{1. Pauses, Silences \& Abrupt Stops} & Micro-pauses \texttt{(.)}, timed silences \texttt{(1.5)}, pregnant/dramatic pauses before a loaded word, hesitant stumbles, abrupt self-interruptions or cut-offs (\texttt{---}), and forced silences (muting a panelist's mic or talking over them). & Anchor pauses \texttt{(1.5)} before saying \textit{"'loota'"} for dramatic weight; Addl. DCP cuts off mid-sentence (\textit{"wo toh--- abhi investigation chal raha hai"}) when asked if police struck first; panelist's mic faded during counter-question. & Where does the pause or stop occur? Is it deliberate/rhetorical (for suspense, emphasis, or gravitas) or defensive/evasive (cognitive load, self-repair, avoiding blame)? Who controls turn-taking and silence? \\ \hline
\textbf{2. Trailing Sentences \& Unfinished Utterances} & Clauses that trail off (\texttt{...}), drop in volume at the end, or leave the conclusion unspoken (aposiopesis) so the audience supplies the implication without the speaker taking explicit responsibility. & \textit{"Aur phir wahan jo log ghuse the... aap samajh hi sakte hain..."}; \textit{"500 salon ke sangharsh ke baad agar aisa..."} trailing off with a slow head shake. & What exact thought is left unfinished? What presupposition or shared bias is the viewer invited to fill in? Does trailing off act as a safe hedge to imply higher-level complicity or "anti-national" motive without stating it? \\ \hline
\textbf{3. Intonation, Pitch, Volume \& Prosody} & Pitch contour (rising \texttt{$\uparrow$} for incredulity/challenge; falling \texttt{$\downarrow$} for finality), volume shifts (shouting vs. hushed whisper), emphatic word stress, speech rate (rapid-fire urgency vs. slow solemnity), and tonal register (sarcastic, reverential, combative). & Hushed, reverential cadence on \textit{"Prabhu Shri Ram ke Mandir"} vs. sharp, high-pitched interrogation on \textit{"par Akhilesh ji Ram Mandir kyon$\uparrow$ nahin jate?!"}; satirical sing-song inflection on \textit{"priventiv arenjment"}. & Which exact words receive pitch peaks or volume stress? How does vocal tone turn neutral words into irony, mockery, or moral outrage? Does the speaker shift prosodic register depending on which actor is discussed? \\ \hline
\textbf{4. Coughs, Sighs \& Non-Lexical Vocalisations} & Audible sighs (\texttt{[sigh]}), scoffs/snorts (\texttt{[scoff]}), ironic chuckles (\texttt{[chuckle]}), throat-clearing/coughs (\texttt{[cough]}), sharp intakes of breath/gasps (\texttt{[gasp]}), and tongue-clicks (\texttt{[tsk]}). & Heavy \texttt{[sigh]} before reporting Trust oversight lapses (signalling reluctant sorrow rather than hostility); mocking \texttt{[scoff]} or \texttt{[chuckle]} when reading an opponent's denial; \texttt{[cough / throat-clear]} when pivoting away from an awkward fact. & What emotion or epistemic stance does the vocalisation perform (exasperation, moral disgust, disbelief, ridicule, discomfort)? Does a sigh or scoff pre-judge a speaker's claim before any verbal rebuttal is made? \\ \hline
\textbf{5. Hand Gestures, Facial Expressions \& Posture (Kinesics)} & Hand movements (air-quotes, finger-pointing/wagging, desk-thumping, open palms, counting on fingers, dismissive wave), facial cues (raised eyebrows, eye-rolls, smirk, furrowed brow), head nods/shakes, leaning into camera, and visual props (bandaged limbs, documents). & Anchor forms physical air-quotes on \textit{"'loota'"} or \textit{"'peaceful' protest"}; host points finger/pen at lens during an ultimatum; injured DCP holds bandaged hand up to camera; creator raises one eyebrow and leans in for a conspiratorial reveal. & How do gestures reinforce, contradict, or sharpen the spoken words? During split-screen listening, does the host's face (nodding vs. smirking/eye-rolling) cue the audience on whom to trust? How are bodily visuals used as evidence? \\ \hline
\end{longtable}

\subsection*{Quick Multimodal Transcription Notation Key (use inside transcript excerpts)}

\begin{longtable}{|p{0.15\textwidth}|p{0.15\textwidth}|p{0.15\textwidth}|p{0.15\textwidth}|p{0.15\textwidth}|p{0.15\textwidth}|}
\hline
Symbol & Meaning & Example in transcript & Symbol & Meaning & Example in transcript \\ \hline \endhead
\texttt{(.)} & Brief micro-pause (< 0.5 sec) & \texttt{woh (.) suparavijan lacking tha} & \texttt{$\uparrow$ / $\downarrow$} & Marked rising / falling pitch & \texttt{kyon$\uparrow$ FIR nahin karaee$\downarrow$} \\ \hline
\texttt{(1.5)} & Timed pause in seconds & \texttt{kisne (1.5) 'loota'?} & \texttt{WORD / \_word\_} & Loud volume / emphatic stress & \texttt{ye PAAP hai / \_bilkul\_ galat} \\ \hline
\texttt{---} & Abrupt stop, cut-off, or self-repair & \texttt{police ki taraf se--- abhi jaanch...} & \texttt{$^\circ$word$^\circ$} & Softened / whispered delivery & \texttt{$^\circ$ander ki baat ye hai$^\circ$} \\ \hline
\texttt{...} & Trailing off / unfinished sentence & \texttt{aap jaante hain ki kaun log...} & \texttt{[sigh] [cough] [scoff]} & Non-lexical vocalisations & \texttt{[sigh] 500 salon ke baad...} \\ \hline
\texttt{[overlap]} & Simultaneous / overlapping speech & \texttt{[Anjana: suniye---] [Guest: nahin!]} & \texttt{\{\{gesture: ...\}\}} & Hand gesture / facial cue / posture & \texttt{\{\{air-quotes on 'loota'\}\}} \\ \hline
\end{longtable}

\textbf{Methodological Reminder:} A non-verbal cue is not automatically manipulative or ideological---a cough, pause, or rapid cadence may reflect studio audio delay, live time pressure, or a speaker's habitual style. Always check (1) whether the cue clusters with verbal features, and (2) whether the speaker's delivery shifts noticeably across different actors or events (within-outlet baseline).\par\smallskip

\bigskip\hrule\bigskip

\subsection*{Event 1. Ram Mandir donation controversy (Non-Verbal \& Paralinguistic Layer)}

\begin{itemize}
  \item \textbf{Focus:} How do prosody, pauses, trailing sentences, vocalisations (sighs, coughs, scoffs), and body language frame the controversy---as solemn devotional sorrow, administrative scrutiny, or partisan combat?
  \item \textbf{Outlet:} Sarthak Goswami (\textit{Sunday Show}) | \textbf{Date:} 23 June 2021
  \item \textbf{Video title:} Ram Mandir Scam, Waqf Amendment Bill, Coaching Mafia, Floods \& Amritpal | \href{https://youtu.be/lb7gZnyZH0c}{https://youtu.be/lb7gZnyZH0c} | \textbf{Timestamp(s):} \texttt{[00:00]}--\texttt{[05:38]} (Full video: \texttt{19:33}; Ram Mandir segment: \texttt{[00:00]}--\texttt{[05:38]})
  \item \textbf{Analyst:} Baani Solanki (\texttt{2025114016}) | \textbf{Language:} Hinglish (Hindi in Roman script + English terms)
  \item \textbf{Format / Visual setup:} \texttt{[ ]} TV Panel Debate \texttt{[ ]} TV Bulletin \texttt{[x]} Solo YouTube Explainer \texttt{[ ]} Interview / Discussion
  \item \textbf{Primary speaker(s) observed:} Sarthak Goswami (solo studio presenter seated at desk, speaking directly to camera with jump-cut editing, meme inserts, and document overlays)
\end{itemize}

\subsubsection*{Event-specific non-verbal questions (Event 1)}

\noindent\textbf{1.} \textbf{Vocal \& gestural marking of core labels ("scam", "loota", "paap", "gadabdi", "avyavastha"): When the controversy is named, what non-verbal cues accompany the label? (Look for physical air-quotes, vocal scare-quotes, dramatic pauses before the word, hushed/reverential drops in volume, or shouted stress.)}\par\smallskip
\begin{itemize}
  \item \textbf{Notes:}
Sarthak Goswami avoids both the hushed, mournful devotional register (\texttt{"paap"}) and high-decibel shouting (\texttt{"scam!"}). Instead, when introducing the controversy at \texttt{[01:13]}, he uses a brief micro-pause \texttt{(.)} followed by raised eyebrows and open, upturned palms \texttt{\{\{gesture: open palms raised with arched eyebrows\}\}} on \textit{"Ayodhya Ram Mandir ke naam par (.) \_bada\_ ghotala$\uparrow$"} (\texttt{[01:13]}). Rather than endorsing \textit{"ghotala"} as his own verdict, his rising pitch \texttt{$\uparrow$} and wide-eyed expression perform \textbf{vocal scare-quotes} that mark the label as a dramatic headline circulating in public discourse. Similarly, when asking \textit{"kya sach mein Ram Mandir ke chande mein chori hui hai$\uparrow$"} (\texttt{[01:16]}), he leans slightly forward with an inquisitive head tilt, framing \textit{"chori"} as an open investigative question rather than an established moral sin.\par\smallskip
\end{itemize}

\noindent\textbf{2.} \textbf{Differential demeanour toward actors (Temple, Trust leadership, counting clerks, donors, opposition accusers): How do the speaker's voice pitch, facial expression, and posture shift when referring to Prabhu Shri Ram / devotees vs. top figures (Champat Rai, Nripendra Misra, CM Yogi) vs. arrested counting staff vs. opposition leaders (Akhilesh Yadav, Awadhesh Prasad)?}\par\smallskip
\begin{itemize}
  \item \textbf{Notes:}
Sarthak maintains a remarkably uniform, conversational baseline across all political and institutional actors, punctured equally by deadpan comedic irony:\par\smallskip
  \item \textbf{Donors / Devotees:} Referred to with warm, relatable facial animation and an inclusive hand gesture toward the camera (\texttt{"aapne aur humne milkar chanda diya tha"}, \texttt{[01:36]}), treating viewers as fellow citizens who contributed in good faith.
  \item \textbf{Opposition Accusers (Sanjay Singh, Pawan Pandey):} Introduced in a brisk, matter-of-fact narrative cadence (\texttt{[01:55]}), pointing to screen graphics displaying their press conference clips without scowling or mocking them.
  \item \textbf{Trust Leadership (Champat Rai):} When quoting Champat Rai's counter-attack comparing critics to \textit{"Mahatma Gandhi ke hatyare"} (\texttt{[04:10]}), Sarthak pauses \texttt{(1.0)}, deadpans to the camera with a raised eyebrow and an ironic half-smirk \texttt{\{\{gesture: deadpan stare into lens + single raised eyebrow\}\}}, letting the disproportionate hyperbole of the Trust's defence speak for itself.
  \item \textbf{Intermediaries (Sultan Ansari, Ravi Mohan Tiwari):} Listed with methodical finger-counting (\texttt{[02:10]}, \texttt{[02:27]}) and steady eye contact, treating them as transactional actors in a real-estate flip rather than demonised villains.
\end{itemize}

\noindent\textbf{3.} \textbf{Solemn devotional grief vs. high-decibel partisan spectacle: How do pauses, sighs, speech rate, and hand gestures construct the segment's emotional register? Where is delivery slow, hushed, and mournful (sighs, bowed head, folded hands), and where does it escalate into rapid-fire shouting, desk-thumping, or mocking irony?}\par\smallskip
\begin{itemize}
  \item \textbf{Notes:}
Sarthak explicitly rejects both solemn devotional grief and high-decibel partisan spectacle, operating instead in a \textbf{fast-paced, satirical infotainment register} (\textasciitilde{}155--165 words per minute). At \texttt{[00:34]}, he explicitly parodies TV news histrionics: \textit{"TV news kholte hi lagta hai (.) duniya \_khatam\_ hone wali hai"} (\texttt{[00:34]}), widening his eyes in mock panic and waving both hands before dropping back to a relaxed smile and level tone (\texttt{"isiliye shanti se samajhte hain"}, \texttt{[00:39]}). There are zero mournful bowed heads or folded hands, and zero desk-thumping outbursts. Emotional peaks are replaced by dry, punchy comedic pauses---such as a \texttt{(1.0)} beat before delivering the satirical line \textit{"zameen ki keemat (.) \_rocket\_ ki speed se badh gayi$\uparrow$"} (\texttt{[02:38]}) accompanied by an upward shooting hand gesture \texttt{\{\{gesture: right hand shooting diagonally upward\}\}}.\par\smallskip
\end{itemize}

\noindent\textbf{4.} \textbf{Non-verbal evaluation of financial scrutiny: When questions about unreceipted offerings, lack of frisking, or land deals are raised, how does the host/speaker react non-verbally (earnest nodding, leaning into the camera, open-palm appeals vs. eye-rolls, scoffs, dismissive hand waves, or incredulous rising pitch $\uparrow$)?}\par\smallskip
\begin{itemize}
  \item \textbf{Notes:}
Financial scrutiny is treated with \textbf{earnest, pedagogical engagement} mixed with astonished prosody at the numbers. During the chronological breakdown of the land price jump from Rs 2 crore to Rs 18.5 crore in five minutes (\texttt{[02:02]}--\texttt{[02:35]}), Sarthak leans into the desk, uses his left index finger to tap out the exact timestamps on his right hand \texttt{\{\{gesture: counting minutes on fingers\}\}}, and stresses the numerals with sharp pitch peaks: \textit{"\_panch\_ minute baad (.) \_saadhe atharah\_ crore$\uparrow$"} (\texttt{[02:27]}). When presenting the Trust's explanation (old 2011 agreement vs. prime location valuation, \texttt{[03:18]}--\texttt{[03:42]}), he shifts to a balanced two-handed weighing gesture \texttt{\{\{gesture: alternating left and right open palms like a balance scale\}\}}, visually performing fairness while still arching his eyebrows on the unresolved timing gap (\texttt{[04:39]}).\par\smallskip
\end{itemize}

\noindent\textbf{5.} \textbf{Performing certainty vs. hedging and innuendo (pauses, stops, trailing off): Are allegations delivered with definitive falling intonation $\downarrow$, finger-counting, and pointing at on-screen evidence (CCTV, documents), or hedged with hesitant pauses (.), abrupt self-interruptions (---), throat-clearing [cough], and trailing sentences (...) that imply higher-level complicity without naming it?}\par\smallskip
\begin{itemize}
  \item \textbf{Notes:}
Sarthak separates verifiable documentary facts from unproven intent using distinct prosodic contours:\par\smallskip
  \item \textbf{Documentary Facts (Certainty):} Delivered with crisp falling intonation \texttt{$\downarrow$}, finger-counting, and direct pointing at on-screen registration documents (\texttt{[02:18]}--\texttt{[02:35]}): \textit{"sham ko saat bajkar dus minute par (.) do crore mein khareedi$\downarrow$ (\texttt{[02:18]}) ... aur saat bajkar pandrah minute par (.) saadhe atharah crore ka agreement ho gaya$\downarrow$"} (\texttt{[02:27]}).
  \item \textbf{Unproven Intent / Valuation Claims (Hedging \& Skeptical Suspension):} Framed via trailing sentences \texttt{...} and open-ended rising intonation \texttt{$\uparrow$} at the close of the segment (\texttt{[04:39]}--\texttt{[05:18]}): \textit{"Ab ye waqai purane agreement ki wajah se tha (.) ya phir kuch aur..."} (\texttt{[04:45]}), accompanied by a shoulder shrug and a sideways head tilt \texttt{\{\{gesture: bilateral shoulder shrug + head tilt\}\}}. This allows him to suggest that the 5-minute flip smells suspicious without making a legally actionable accusation of corruption.
\end{itemize}

\noindent\textbf{6.} \textbf{Turn-taking, interruptions, and acoustic treatment of denials: In debates/interviews: who receives calm, uninterrupted speaking time and sympathetic nods, and who is cut off mid-sentence (---), talked over, or muted? In solo videos: how are clips of official denials framed via reaction pauses, raised eyebrows, or sarcastic chuckles?}\par\smallskip
\begin{itemize}
  \item \textbf{Notes:}
Because this is a solo YouTube explainer, Sarthak holds 100\% of the acoustic floor without panelist interruptions. However, his \textbf{multimodal framing of official denials} is highly revealing. After summarizing Champat Rai's statement that the Trust followed all legal procedures and that critics are like Gandhi's assassins (\texttt{[04:10]}--\texttt{[04:18]}), Sarthak inserts a deliberate \texttt{(1.2)} reaction pause, lets out a soft, bemused \texttt{[chuckle]}, and remarks with a wry smile: \textit{"Matlab (.) sawaal poochhne par seedha \_itna\_ bada comparison$\uparrow$"} (\texttt{[04:22]}). While he grants the Trust's substantive defence (registered bank transfers, old agreement dates, \texttt{[03:18]}--\texttt{[03:55]}) clear, uninterrupted explanation time, his post-quote facial reaction punctures the Trust's moral deflection.\par\smallskip
\end{itemize}

\noindent\textbf{7.} \textbf{Performing the boundary between religious devotion and institutional critique: How does the presenter use vocal tone, sighs, and body language to signal personal reverence for the Temple (or secular objectivity) while criticising---or defending---the Trust's financial management?}\par\smallskip
\begin{itemize}
  \item \textbf{Notes:}
Sarthak explicitly builds a verbal and non-verbal firewall between faith (\texttt{aastha}) and institutional accounting (\texttt{hisaab}) at \texttt{[04:55]}--\texttt{[05:25]}. When speaking of ordinary citizens' devotion (\textit{"crore-on logon ki aastha judi hai"\texttt{, }[04:58]\texttt{), his voice softens slightly, his pace slows, and he places his right hand lightly over his chest }\{\{gesture: right palm placed flat on chest\}\}\texttt{ with an earnest, unsmiling expression. Immediately afterward, on }"isiliye (.) \_ek-ek\_ paise ka hisaab saaf hona chahiye$\downarrow$"* (}[05:06]\texttt{), he straightens his posture and taps his index finger firmly on the desk }\{\{gesture: single firm index-finger tap on desk\}\}`. This bodily shift signals that demanding financial transparency from the Trust is an act of respect toward devotees rather than an attack on religion.\par\smallskip
\end{itemize}

\noindent\textbf{8.} \textbf{Self-presentation as the "guardian of public trust": How does the speaker use direct eye-gaze, posture, moral sighs, or emphatic gestures (pointing at the viewer, hands on heart/desk) to position themselves as the authentic protector of devotees or public accountability?}\par\smallskip
\begin{itemize}
  \item \textbf{Notes:}
Rather than posing as an angry moral crusader (like Arnab Goswami) or a grieving devotee (like Anjana Om Kashyap), Sarthak performs the role of the \textbf{calm, rational buddy-explainer} who protects the viewer from media noise. He maintains steady, relaxed eye-gaze into the lens, sits upright with open shoulders, uses inclusive pointing gestures toward the camera (\texttt{"aap bataiye comments mein"}, \texttt{[05:22]}), and lets out a light, relatable \texttt{[sigh]} of fatigue at media polarisation (\texttt{[00:34]}). His closing open-palm gesture \texttt{\{\{gesture: both palms open toward viewer\}\}} (\texttt{[05:25]}) hands the final judgment to the audience, cementing his persona as an objective, citizen-first mediator.\par\smallskip
\end{itemize}

\subsubsection*{Apply the five non-verbal features (Event 1)}

\begin{longtable}{|p{0.31\textwidth}|p{0.31\textwidth}|p{0.31\textwidth}|}
\hline
Non-verbal feature & Timestamp + spoken anchor words + non-verbal notation & What it does in context (stance, framing, implication) \\ \hline \endhead
\textbf{1. Pauses, silences \& stops} & \texttt{[02:27]} \textit{"Saat bajkar dus minute par do crore... aur (1.0) saat bajkar pandrah minute par (.) \_saadhe atharah\_ crore$\downarrow$"} & The deliberate \texttt{(1.0)} rhetorical pause before the second price creates suspense and highlights the absurdity of a 5-minute, Rs 16.5-crore valuation jump without needing shouted outrage. \\ \hline
\textbf{2. Trailing sentences (\texttt{...})} & \texttt{[04:45]} \textit{"Ab ye waqai purane agreement ki wajah se tha (.) ya phir kuch aur..."} \texttt{\{\{gesture: shoulder shrug + raised eyebrows\}\}} & Leaves the alternative explanation (\textit{"kuch aur..."}, i.e., insider profiteering) unspoken so the viewer supplies the suspicion while the host preserves his neutral explainer stance. \\ \hline
\textbf{3. Intonation, pitch \& volume} & \texttt{[02:38]} \textit{"Zameen ki keemat (.) \_ROCKET\_ ki speed se badh gayi$\uparrow$"} (marked rising pitch \texttt{$\uparrow$} + loud stress on \texttt{ROCKET}) & Uses playful prosodic hyperbole and pitch elevation to mock the economic improbability of the transaction while keeping the mood light and accessible. \\ \hline
\textbf{4. Coughs, sighs \& vocalisations} & \texttt{[04:22]} \texttt{[(1.2) pause + [chuckle]]} \textit{"Matlab (.) sawaal poochhne par seedha \_itna\_ bada comparison$\uparrow$"} & The brief ironic \texttt{[chuckle]} after quoting Champat Rai's "Gandhi's killers" remark non-verbally evaluates the Trust's rhetoric as defensive overreaction. \\ \hline
\textbf{5. Hand gestures, face \& posture} & \texttt{[03:18]}--\texttt{[03:42]} \textit{"Ek taraf aarop hai... doosri taraf Trust ka kehna hai..."} \texttt{\{\{gesture: alternating left and right open palms like a weighing scale\}\}} & Physically embodies structural balance ("both-sides" journalism), visually assuring the audience that he is weighing accusation and defence fairly. \\ \hline
\end{longtable}

\subsubsection*{Short summary (Event 1 Non-Verbal Synthesis)}

\begin{itemize}
  \item \textbf{What overall interpretation or emotional stance does the non-verbal delivery encourage?}
The non-verbal delivery encourages \textbf{skeptical, good-humoured civic scrutiny}. By pairing rapid, upbeat delivery and deadpan eyebrow raises with clear finger-counted timelines, Sarthak prevents the viewer from falling into either communal defensiveness or sensationalist panic, framing the controversy as a legitimate administrative question that deserves a transparent audit.\par\smallskip
  \item \textbf{Whose voice or physical presence is acoustically/visually foregrounded, undermined, or silenced?}
Sarthak's own mediating voice dominates the acoustic space, while both the Opposition's documentary evidence (registration timestamps) and the Trust's official rebuttals are visually foregrounded via clean on-screen graphics. However, Champat Rai's emotional counter-accusations are subtly \textbf{undermined} via reaction pauses and bemused smirks, whereas the ordinary donor's perspective is \textbf{validated} through warm, inclusive hand-to-chest gestures.\par\smallskip
  \item \textbf{What evidence supports your non-verbal reading? What complicates it (e.g., speaker baseline, format)?}
The clustering of timed pauses \texttt{(1.0)}, finger-counting, and emphatic numeral stress around the \texttt{7:10 PM} vs. \texttt{7:15 PM} registry deeds (\texttt{[02:18]}--\texttt{[02:38]}) shows where the explainer locates the core evidentiary problem. What complicates a purely ideological reading is \textbf{speaker baseline and format}: Sarthak's fast-paced jump cuts, eyebrow raises, and ironic chuckles are his trademark \textit{Sunday Show} infotainment baseline across all topics in the episode (Waqf Bill, Coaching Mafia, Floods), meaning his humour is a genre convention for youth engagement rather than targeted hostility toward the Temple Trust.\par\smallskip
\end{itemize}

\bigskip\hrule\bigskip

\subsection*{Event 2. Students' protest regarding NEET paper leak (Non-Verbal \& Paralinguistic Layer)}

\begin{itemize}
  \item \textbf{Focus:} How do prosody, pauses, trailing sentences, vocalisations, and visual body language construct the legitimacy of student protesters, police use of force, and state accountability?
  \item \textbf{Outlet:} Sarthak Goswami (\textit{Sunday Show}) | \textbf{Date:} 23 June 2024
  \item \textbf{Video title:} NEET Scam Explained! Grace Marks, Paper Leak \& Protest | \href{https://youtu.be/EM1GeZufO6E}{https://youtu.be/EM1GeZufO6E} | \textbf{Timestamp(s):} \texttt{[00:00]}--\texttt{[18:45]}
  \item \textbf{Analyst:} Baani Solanki (\texttt{2025114016}) | \textbf{Language:} Hinglish (Hindi in Roman script + English terms)
  \item \textbf{Format / Visual setup:} \texttt{[ ]} TV Panel Debate \texttt{[ ]} TV Bulletin \texttt{[x]} Solo YouTube Explainer \texttt{[ ]} Interview / Discussion
  \item \textbf{Primary speaker(s) observed:} Sarthak Goswami (solo presenter at desk using rapid jump-cuts, scorecards, statistical overlays, and brief clips of student protesters and NTA press briefings)
\end{itemize}

\subsubsection*{Event-specific non-verbal questions (Event 2)}

\noindent\textbf{1.} \textbf{Prosodic \& gestural marking of participant labels ("students/bachche" vs. "protesters" vs. "upadravi/goons"): How do vocal tone, volume, facial expression, and hand gestures shift across labels? (e.g., softened, empathetic pitch for "students" vs. harsh, staccato stress, scowls, or accusatory finger-pointing on "upadraviyon", "goons", or "anti-social elements".)}\par\smallskip
\begin{itemize}
  \item \textbf{Notes:}
Sarthak never uses criminalising labels like \textit{"upadravi"} or \textit{"goons"}. Throughout the video, he refers to the participants as \textit{"students"}, \textit{"bachche"}, or \textit{"aspirants"} (\texttt{[01:12]}, \texttt{[08:40]}, \texttt{[14:15]}) using a \textbf{softened, empathetic vocal tone}, slightly lowered pitch, and earnest furrowed brows \texttt{\{\{gesture: sympathetic brow furrow + open palms\}\}}. When describing their hard work (\textit{"do-do saal kota mein baith kar taiyari karte hain"\texttt{, }[01:28]\texttt{), his speech slows down and his voice carries genuine warmth. By contrast, when naming }"paper leak mafia"\textit{ or }"solver gang"* (}[09:45]\texttt{), his tone turns crisp and clipped, accompanied by a sharp downward chopping motion of the right hand }\{\{gesture: vertical hand chop\}\}`.\par\smallskip
\end{itemize}

\noindent\textbf{2.} \textbf{Delivery of students' stated demands: When the protesters' core demands (Education Minister's resignation, compensation, release of Sonam Wangchuk, exam reform) are mentioned, are they delivered slowly with deliberate pauses and finger-counting, rushed through in a flat monotone, mocked with air-quotes/scoffs, or trailed off (...)?}\par\smallskip
\begin{itemize}
  \item \textbf{Notes:}
The students' core demands---cancellation of fraudulent grace marks, a transparent Supreme Court-monitored probe into the Bihar/Godhra paper leaks, and complete overhaul of the NTA (\texttt{[13:50]}--\texttt{[15:10]})---are delivered \textbf{slowly and methodically with deliberate pauses and finger-counting}. Sarthak holds up three fingers one by one \texttt{\{\{gesture: sequential counting on fingers 1, 2, 3\}\}}, pausing \texttt{(0.8)} after each demand and nodding firmly with falling intonation \texttt{$\downarrow$} (*"Pehli maang (.) NTA ki accountability tay ho$\downarrow$"\texttt{, }[14:02]`), visually and acoustically validating their demands as rational and structured.\par\smallskip
\end{itemize}

\noindent\textbf{3.} \textbf{Where does vocal \& kinetic energy peak---grievance, street disruption, or electoral horse-race?: Track where the speaker's volume, pitch range, pace, and gesture intensity reach their climax: during accounts of the paper leak/student suicides, during street clashes and barricade breaches, or during partisan BJP-vs-Opposition sparring.}\par\smallskip
\begin{itemize}
  \item \textbf{Notes:}
Vocal and kinetic energy peaks unmistakably during the \textbf{statistical and administrative breakdown of student grievances and NTA anomalies} (\texttt{[03:15]}--\texttt{[07:50]}), rather than street clashes or partisan sparring:\par\smallskip
  \item \textbf{Climax 1 (Statistical Anomaly, \texttt{[03:42]}):} When revealing that \textit{"\_sadsath\_ (67) students ko full \texttt{720} out of \texttt{720} marks mile$\uparrow$!"} and that \textit{"\_ek hi\_ centre (Bahadurgarh) se \_chhah\_ toppers nikle$\uparrow$!"} (\texttt{[05:18]}), Sarthak's pitch rises sharply \texttt{$\uparrow$}, his volume surges, and he leans right into the camera lens with wide eyes and both hands spread wide in disbelief \texttt{\{\{gesture: leaning into lens + wide-eyed disbelief\}\}}.
  \item \textbf{Climax 2 (Mental Toll on Aspirants, \texttt{[15:40]}):} Energy shifts from high-volume disbelief to intense, quiet gravity \texttt{$^\circ$word$^\circ$} when discussing student depression and suicides, whereas partisan BJP-vs-Congress blame games (\texttt{[16:10]}) are brushed past rapidly with a dismissive hand wave.
\end{itemize}

\noindent\textbf{4.} \textbf{Non-verbal attribution of responsibility (sighs, stops, head shakes, pointing): When examination authorities (NTA), the Education Minister, police, opposition politicians, or students are named as responsible, what non-verbal cues accompany the attribution (exasperated sighs, mocking laughs, abrupt self-stops ---, or direct finger-pointing at the lens)?}\par\smallskip
\begin{itemize}
  \item \textbf{Notes:}
Responsibility is squarely attributed to the \textbf{National Testing Agency (NTA)} and the \textbf{Education Ministry's initial denialism} through exasperated vocalisations and mocking gestures:\par\smallskip
  \item When citing NTA's excuse of \textit{"loss of time"} and mysterious \textit{"normalization formula"} (\texttt{[06:10]}), Sarthak lets out an audible exasperated \texttt{[sigh]} followed by a mocking \texttt{[scoff]}, rolling his eyes and making physical air-quotes \texttt{\{\{gesture: air-quotes with both hands on 'loss of time'\}\}}.
  \item When referencing Education Minister Dharmendra Pradhan's initial claim that \textit{"koi paper leak ka saboot nahin hai"} (\texttt{[11:05]}) right before Bihar Police arrested the solver gang with burnt question papers (\texttt{[11:25]}), Sarthak stops mid-sentence \texttt{---}, shakes his head slowly from side to side \texttt{\{\{gesture: slow incredulous head shake\}\}}, and points at the on-screen Bihar EOU evidence report (\texttt{[11:32]}).
\end{itemize}

\noindent\textbf{5.} \textbf{Prosody of questions---open inquiry vs. loaded rhetorical challenge: Are questions about the protest delivered with level, open intonation, or with sharp rising pitch $\uparrow$, sarcastic stress, raised eyebrows, and trailing clauses (...) that presuppose protesters are violent or politically manipulated?}\par\smallskip
\begin{itemize}
  \item \textbf{Notes:}
Unlike TV studio debates that direct hostile rhetorical questions at protesters, Sarthak directs all his \textbf{loaded rhetorical challenges at the NTA and the government}. For instance, at \texttt{[04:30]}, after explaining how \texttt{718} and \texttt{719} marks are mathematically impossible under NEET's \texttt{+4 / -1} marking scheme, he asks: \textit{"NTA walon (.) ye kaunsa \_naya\_ ganit hai bhai$\uparrow$?!"} (\texttt{[04:35]}), delivered with sharp rising pitch \texttt{$\uparrow$}, sarcastic stress on \textit{"\_naya\_ ganit"}, and a classic palm-up "explain this" gesture \texttt{\{\{gesture: right arm extended, palm upturned toward screen\}\}}. Questions about the students' future (\texttt{"ab un }24\texttt{ lakh bachchon ka kya hoga$\downarrow$"}, \texttt{[17:10]}) are delivered with a falling, empathetic sigh rather than suspicion.\par\smallskip
\end{itemize}

\noindent\textbf{6.} \textbf{Visual \& paralinguistic performance of injury and violence: How are claims of violence substantiated non-verbally on screen? (Compare camera close-ups and pained vocal tone on police bandages/fractured hands vs. visual/acoustic treatment of student head injuries, lathi-charge footage, or defensive stammers/pauses when an officer or guest is asked who used force first.)}\par\smallskip
\begin{itemize}
  \item \textbf{Notes:}
Sarthak frames the crisis primarily as \textbf{institutional violence against students' futures} rather than a street riot. When showing brief B-roll of students protesting in scorching June heat and facing police barricades/detentions (\texttt{[13:20]}--\texttt{[13:45]}), his voice drops in pitch and slows down: \textit{"Chalis degree ki garmi mein (.) students sadak par utarne ko \_majboor\_ hain$\downarrow$"} (\texttt{[13:28]}). There is zero focus on police injuries or broken barricades; instead, visual close-ups highlight exhausted students holding placards and tearful parents outside exam centres (\texttt{[13:35]}).\par\smallskip
\end{itemize}

\noindent\textbf{7.} \textbf{Asymmetric delivery of police actions vs. protester actions: Compare the vocal cadence, pauses, and hand gestures used when describing police actions (e.g., calm, matter-of-fact administrative tone on "tear gas chhode gae" / "preventive arrangement") versus protester actions (alarmist pitch, striking gestures, or ironic softening as "retaliation").}\par\smallskip
\begin{itemize}
  \item \textbf{Notes:}
The asymmetry in Sarthak's video runs \textbf{in favour of the students and Bihar investigating police}, and \textbf{against central exam authorities}:\par\smallskip
  \item \textbf{Protester actions} are framed with sympathetic cadence and solidarity gestures (*"apne haq ke liye aawaz utha rahe hain"\texttt{, }[13:40]`).
  \item \textbf{Bihar Police / EOU investigation} into the Patna safe-house where burnt admit cards were recovered (\texttt{[09:50]}--\texttt{[10:40]}) is narrated with brisk, approving energy and pointing gestures at the forensic proof.
  \item \textbf{NTA's administrative actions} (releasing results 10 days early on 4 June---Election Result Day---at \texttt{[07:15]}) are narrated with a conspiratorial lean into the camera, a narrowed gaze, and a whispered aside: \textit{"Chaar June ko (.) jab poora desh election results dekh raha tha (.) $^\circ$chupke se NEET ka result nikaal diya$^\circ$"} (\texttt{[07:22]}).
\end{itemize}

\noindent\textbf{8.} \textbf{Acoustic \& visual space granted to students' own voices: Are protesting students heard in their own voice with clear audio, steady close-ups, and empathetic listening nods, or are they reduced to noisy background chaos B-roll while studio anchors/politicians speak over them?}\par\smallskip
\begin{itemize}
  \item \textbf{Notes:}
Students and educators are granted \textbf{clear, unmuted acoustic space}. At \texttt{[12:40]}--\texttt{[13:18]}, Sarthak cuts directly to clean audio bytes of NEET aspirants and educators explaining the rank inflation---where a score of \texttt{650} that previously guaranteed a government medical seat now resulted in a rank beyond \texttt{25,000}. During the transition back to studio (\texttt{[13:20]}), Sarthak nods slowly with a serious, unsmiling expression \texttt{\{\{gesture: slow, serious nod + pursed lips\}\}}, allowing the students' frustration to linger for \texttt{(1.0)} second before resuming commentary.\par\smallskip
\end{itemize}

\noindent\textbf{9.} \textbf{Non-verbal shift between the NEET paper-leak issue and the street protest: Does the speaker's non-verbal demeanour shift from sympathetic concern (sombre tone, sighs, solidarity gestures) when discussing the exam leak itself to combative, authoritarian, or dismissive delivery once students march toward Parliament?}\par\smallskip
\begin{itemize}
  \item \textbf{Notes:}
\textbf{No shift occurs.} Unlike pro-establishment TV channels that pivot from sympathy over the exam leak to hostility once students protest on the street, Sarthak maintains a continuous arc of \textbf{empathetic solidarity} from the exam hall to the street protest. He treats the street protests (\texttt{[13:20]}--\texttt{[14:30]}) as the inevitable, justified consequence of institutional deafness: \textit{"Jab sunwai nahin hogi (.) toh bachche sadak par hi aayenge na$\downarrow$"} (\texttt{[13:42]}), accompanied by an open-palm shrug of obviousness \texttt{\{\{gesture: bilateral open-palm shrug\}\}}.\par\smallskip
\end{itemize}

\subsubsection*{Apply the five non-verbal features (Event 2)}

\begin{longtable}{|p{0.31\textwidth}|p{0.31\textwidth}|p{0.31\textwidth}|}
\hline
Non-verbal feature & Timestamp + spoken anchor words + non-verbal notation & What it does in context (stance, framing, implication) \\ \hline \endhead
\textbf{1. Pauses, silences \& stops} & \texttt{[03:42]} \textit{"Is saal NEET mein (1.2) \_SADSATH\_ (\texttt{67}) students ko full \texttt{720} marks mile$\uparrow$!"} & The \texttt{(1.2)} dramatic pause before the number \texttt{67} builds anticipation and makes the statistical impossibility hit the viewer with maximum shock value. \\ \hline
\textbf{2. Trailing sentences (\texttt{...})} & \texttt{[07:28]} \textit{"Result \texttt{14} June ko aana tha, par aaya \texttt{4} June ko---jis din Lok Sabha election ke results the... ab aap khud samajh lijiye..."} & Trails off after pairing the early NEET result release with Election Day, inviting the viewer to infer deliberate media burial by the NTA without stating it as dogma. \\ \hline
\textbf{3. Intonation, pitch \& volume} & \texttt{[04:35]} \textit{"Plus char, minus ek ke hisaab se \texttt{718} aur \texttt{719} marks aa hi nahin sakte---ye kaunsa \_NAYA\_ ganit hai$\uparrow$?!"} / \texttt{[07:22]} \textit{"$^\circ$chupke se result nikaal diya$^\circ$"} & Contrasts loud, rising incredulity \texttt{$\uparrow$} on the impossible scores with a conspiratorial whispered drop \texttt{$^\circ$...$^\circ$} on the Election-Day release, dramatising both incompetence and stealth. \\ \hline
\textbf{4. Coughs, sighs \& vocalisations} & \texttt{[06:12]} \texttt{[[sigh] + [scoff]]} \textit{"Aur NTA ne iska naam diya---'grace marks for loss of time'"} & The exasperated \texttt{[sigh]} and brief \texttt{[scoff]} immediately pre-judge the NTA's "grace marks" rationale as a flimsy post-hoc cover-up before the technical rebuttal even begins. \\ \hline
\textbf{5. Hand gestures, face \& posture} & \texttt{[06:15]} \texttt{\{\{gesture: physical air-quotes on 'grace marks'\}\}} / \texttt{[07:20]} \texttt{\{\{gesture: leaning close to lens + narrowed eyes\}\}} / \texttt{[14:02]} \texttt{\{\{gesture: counting demands on 3 fingers\}\}} & Air-quotes delegitimise NTA's terminology; the camera lean-in signals an insider reveal on the \texttt{4 June} timing; finger-counting lends structured legitimacy to student demands. \\ \hline
\end{longtable}

\subsubsection*{Short summary (Event 2 Non-Verbal Synthesis)}

\begin{itemize}
  \item \textbf{What overall interpretation or emotional stance does the non-verbal delivery encourage?}
The non-verbal delivery encourages \textbf{empathetic outrage on behalf of the students combined with sharp satirical skepticism toward the NTA}. The viewer is guided to feel that \texttt{24 lakh} hardworking young citizens have been cheated by systemic incompetence and a cover-up attempt.\par\smallskip
  \item \textbf{Whose voice or physical presence is acoustically/visually foregrounded, undermined, or silenced?}
\textbf{Students, parents, and Bihar police investigators} are visually and acoustically foregrounded through clear clips, scorecards, and earnest nods. \textbf{The NTA and Education Ministry's early denials} are systematically undermined through air-quotes, scoffs, eye-rolls, and juxtaposition with burnt question-paper evidence.\par\smallskip
  \item \textbf{What evidence supports your non-verbal reading? What complicates it (e.g., speaker baseline, format)?}
The sharpest prosodic shifts---whispered delivery (\texttt{$^\circ$chupke se$^\circ$}, \texttt{[07:22]}), physical air-quotes (\texttt{[06:15]}), and \texttt{1.2}-second dramatic pauses (\texttt{[03:42]})---cluster densely around NTA's procedural anomalies. A complicating factor is that Sarthak's signature comedic meme-inserts occasionally lighten the gravity of the crisis, though he consciously suppresses jokes during the concluding section on student mental health (\texttt{[15:40]}--\texttt{[18:00]}).\par\smallskip
  \item \textbf{How do the non-verbal cues interact with the spoken words? (Do they reinforce, dramatise, soften/hedge, invert meaning via sarcasm/irony, or leave an accusation unspoken via trailing off?)}
The non-verbal cues \textbf{invert official terminology via irony} (turning \textit{"grace marks"} and \textit{"transparent process"} into objects of ridicule via air-quotes and scoffs), \textbf{dramatise statistical anomalies} via pitch peaks \texttt{$\uparrow$}, and \textbf{leave the accusation of political cover-up unspoken via trailing off} (\texttt{...}) on the \texttt{4 June} election-day result timing (\texttt{[07:28]}).\par\smallskip
  \item \textbf{What context, speaker baseline, or format factors support or complicate this interpretation? (e.g., Is this the speaker's normal style across all videos, or a marked shift on this specific actor/claim?)}
Across both Event 1 and Event 2, Sarthak's baseline is brisk, ironic, youth-centric infotainment. However, comparing Event 1 and Event 2 reveals a \textbf{marked shift in emotional investment}: whereas Event 1 (Ram Mandir land deal) stays detached and "both-sides" neutral right to the end, in Event 2 (NEET leak) his voice drops its comedic distance at \texttt{[15:40]} to adopt a sombre, protective tone toward students, reflecting his core youth/student audience demographic.\par\smallskip
  \item \textbf{How does another outlet perform the same moment, actor, or claim non-verbally?}
On \textbf{Aaj Tak} (\texttt{Event 2}), anchor Anjana Om Kashyap performs the NEET student protest with combative rising pitch \texttt{$\uparrow$}, stern finger-wagging, and repeated interruptions of student representatives while holding up close-up visuals of a police officer's bandaged hand to frame protesters as unruly (\textit{"upadravi"}). Sarthak Goswami never once wags a finger at students; his finger-pointing and sarcastic pitch peaks are aimed exclusively at the NTA's impossible \texttt{718}/\texttt{719} scorecards.\par\smallskip
  \item \textbf{My interpretation, with reasons:}
Sarthak Goswami's multimodal performance demonstrates how digital youth explainers use \textbf{paralinguistic irony (scoffs, air-quotes, whispered asides, and trailing clauses)} as a tool of institutional critique. While his verbal script stays carefully anchored in official statistics (\texttt{67} toppers, \texttt{718/719} marks, \texttt{5-minute} land registry gap), his non-verbal layer clearly signals where credibility lies---maintaining cautious neutrality when religious sensitivities are involved (Event 1), but shifting into unambiguous vocal and gestural solidarity when student futures are at stake (Event 2).\par\smallskip
\end{itemize}
\newpage
\section*{Non-Verbal \& Multimodal Discourse Analysis Worksheet}

\textit{Qualitative analysis of pauses, trailing sentences, prosody/intonation, vocalisations (coughs, sighs, scoffs), and kinesics (hand gestures, facial expressions, posture) in political news discourse}\par\smallskip

\bigskip\hrule\bigskip

\subsection*{Guidelines: Five non-verbal \& paralinguistic features to look for in every event}

\begin{longtable}{|p{0.23\textwidth}|p{0.23\textwidth}|p{0.23\textwidth}|p{0.23\textwidth}|}
\hline
Feature & What to look for & Illustrative example (from sample) & Questions to answer \\ \hline \endhead
\textbf{1. Pauses, Silences \& Abrupt Stops} & Micro-pauses \texttt{(.)}, timed silences \texttt{(1.5)}, pregnant/dramatic pauses before a loaded word, hesitant stumbles, abrupt self-interruptions or cut-offs (\texttt{---}), and forced silences (muting a panelist's mic or talking over them). & Anchor pauses \texttt{(1.5)} before saying \textit{"'loota'"} for dramatic weight; Addl. DCP cuts off mid-sentence (\textit{"wo toh--- abhi investigation chal raha hai"}) when asked if police struck first; panelist's mic faded during counter-question. & Where does the pause or stop occur? Is it deliberate/rhetorical (for suspense, emphasis, or gravitas) or defensive/evasive (cognitive load, self-repair, avoiding blame)? Who controls turn-taking and silence? \\ \hline
\textbf{2. Trailing Sentences \& Unfinished Utterances} & Clauses that trail off (\texttt{...}), drop in volume at the end, or leave the conclusion unspoken (aposiopesis) so the audience supplies the implication without the speaker taking explicit responsibility. & \textit{"Aur phir wahan jo log ghuse the... aap samajh hi sakte hain..."}; \textit{"500 salon ke sangharsh ke baad agar aisa..."} trailing off with a slow head shake. & What exact thought is left unfinished? What presupposition or shared bias is the viewer invited to fill in? Does trailing off act as a safe hedge to imply higher-level complicity or "anti-national" motive without stating it? \\ \hline
\textbf{3. Intonation, Pitch, Volume \& Prosody} & Pitch contour (rising \texttt{$\uparrow$} for incredulity/challenge; falling \texttt{$\downarrow$} for finality), volume shifts (shouting vs. hushed whisper), emphatic word stress, speech rate (rapid-fire urgency vs. slow solemnity), and tonal register (sarcastic, reverential, combative). & Hushed, reverential cadence on \textit{"Prabhu Shri Ram ke Mandir"} vs. sharp, high-pitched interrogation on \textit{"par Akhilesh ji Ram Mandir kyon$\uparrow$ nahin jate?!"}; satirical sing-song inflection on \textit{"priventiv arenjment"}. & Which exact words receive pitch peaks or volume stress? How does vocal tone turn neutral words into irony, mockery, or moral outrage? Does the speaker shift prosodic register depending on which actor is discussed? \\ \hline
\textbf{4. Coughs, Sighs \& Non-Lexical Vocalisations} & Audible sighs (\texttt{[sigh]}), scoffs/snorts (\texttt{[scoff]}), ironic chuckles (\texttt{[chuckle]}), throat-clearing/coughs (\texttt{[cough]}), sharp intakes of breath/gasps (\texttt{[gasp]}), and tongue-clicks (\texttt{[tsk]}). & Heavy \texttt{[sigh]} before reporting Trust oversight lapses (signalling reluctant sorrow rather than hostility); mocking \texttt{[scoff]} or \texttt{[chuckle]} when reading an opponent's denial; \texttt{[cough / throat-clear]} when pivoting away from an awkward fact. & What emotion or epistemic stance does the vocalisation perform (exasperation, moral disgust, disbelief, ridicule, discomfort)? Does a sigh or scoff pre-judge a speaker's claim before any verbal rebuttal is made? \\ \hline
\textbf{5. Hand Gestures, Facial Expressions \& Posture (Kinesics)} & Hand movements (air-quotes, finger-pointing/wagging, desk-thumping, open palms, counting on fingers, dismissive wave), facial cues (raised eyebrows, eye-rolls, smirk, furrowed brow), head nods/shakes, leaning into camera, and visual props (bandaged limbs, documents). & Anchor forms physical air-quotes on \textit{"'loota'"} or \textit{"'peaceful' protest"}; host points finger/pen at lens during an ultimatum; injured DCP holds bandaged hand up to camera; creator raises one eyebrow and leans in for a conspiratorial reveal. & How do gestures reinforce, contradict, or sharpen the spoken words? During split-screen listening, does the host's face (nodding vs. smirking/eye-rolling) cue the audience on whom to trust? How are bodily visuals used as evidence? \\ \hline
\end{longtable}

\subsection*{Quick Multimodal Transcription Notation Key (use inside transcript excerpts)}

\begin{longtable}{|p{0.15\textwidth}|p{0.15\textwidth}|p{0.15\textwidth}|p{0.15\textwidth}|p{0.15\textwidth}|p{0.15\textwidth}|}
\hline
Symbol & Meaning & Example in transcript & Symbol & Meaning & Example in transcript \\ \hline \endhead
\texttt{(.)} & Brief micro-pause (< 0.5 sec) & \texttt{woh (.) suparavijan lacking tha} & \texttt{$\uparrow$ / $\downarrow$} & Marked rising / falling pitch & \texttt{kyon$\uparrow$ FIR nahin karaee$\downarrow$} \\ \hline
\texttt{(1.5)} & Timed pause in seconds & \texttt{kisne (1.5) 'loota'?} & \texttt{WORD / \_word\_} & Loud volume / emphatic stress & \texttt{ye PAAP hai / \_bilkul\_ galat} \\ \hline
\texttt{---} & Abrupt stop, cut-off, or self-repair & \texttt{police ki taraf se--- abhi jaanch...} & \texttt{$^\circ$word$^\circ$} & Softened / whispered delivery & \texttt{$^\circ$ander ki baat ye hai$^\circ$} \\ \hline
\texttt{...} & Trailing off / unfinished sentence & \texttt{aap jaante hain ki kaun log...} & \texttt{[sigh] [cough] [scoff]} & Non-lexical vocalisations & \texttt{[sigh] 500 salon ke baad...} \\ \hline
\texttt{[overlap]} & Simultaneous / overlapping speech & \texttt{[Anjana: suniye---] [Guest: nahin!]} & \texttt{\{\{gesture: ...\}\}} & Hand gesture / facial cue / posture & \texttt{\{\{air-quotes on 'loota'\}\}} \\ \hline
\end{longtable}

\textbf{Methodological Reminder:} A non-verbal cue is not automatically manipulative or ideological---a cough, pause, or rapid cadence may reflect studio audio delay, live time pressure, or a speaker's habitual style. Always check (1) whether the cue clusters with verbal features, and (2) whether the speaker's delivery shifts noticeably across different actors or events (within-outlet baseline).\par\smallskip

\bigskip\hrule\bigskip

\subsection*{Event 1. Ram Mandir donation controversy (Non-Verbal \& Paralinguistic Layer)}

\begin{itemize}
  \item \textbf{Focus:} How do prosody, pauses, trailing sentences, vocalisations (sighs, coughs, scoffs), and body language frame the controversy---as solemn devotional sorrow, administrative scrutiny, or partisan combat?
  \item \textbf{Outlet:} Dhruv Rathee (YouTube Channel) | \textbf{Date:} 15 June 2021
  \item \textbf{Video title:} Ayodhya Ram Mandir Scam | Exposed by AAP \& SP | Dhruv Rathee | \href{https://youtu.be/oeH0fOSJuX8}{https://youtu.be/oeH0fOSJuX8} | \textbf{Timestamp(s):} \texttt{[00:00]}--\texttt{[15:42]}
  \item \textbf{Analyst:} Baani Solanki (\texttt{2025114016}) | \textbf{Language:} Hinglish (Hindi in Roman script + English terms)
  \item \textbf{Format / Visual setup:} \texttt{[ ]} TV Panel Debate \texttt{[ ]} TV Bulletin \texttt{[x]} Solo YouTube Explainer \texttt{[ ]} Interview / Discussion
  \item \textbf{Primary speaker(s) observed:} Dhruv Rathee (standing/semi-standing studio presenter against a green-screen documentary backdrop, using structured hand-chops, finger-counting, direct camera gaze, and full-screen document highlights)
\end{itemize}

\subsubsection*{Event-specific non-verbal questions (Event 1)}

\noindent\textbf{1.} \textbf{Vocal \& gestural marking of core labels ("scam", "loota", "paap", "gadabdi", "avyavastha"): When the controversy is named, what non-verbal cues accompany the label? (Look for physical air-quotes, vocal scare-quotes, dramatic pauses before the word, hushed/reverential drops in volume, or shouted stress.)}\par\smallskip
\begin{itemize}
  \item \textbf{Notes:}
Unlike TV anchors who place air-quotes around \textit{"scam"} or soften it into \textit{"avyavastha"}, Dhruv Rathee delivers the labels \textit{"scam"}, \textit{"ghotala"}, and \textit{"corruption"} with \textbf{unhedged, assertive falling intonation \texttt{$\downarrow$} and firm downward hand-chops}. In the opening hook (\texttt{[00:18]}), he pauses \texttt{(0.8)} before stressing the phrase: \textit{"Ram Mandir ke chande ke paise mein (0.8) \_bada\_ GHOTALA saamne aaya hai$\downarrow$"} (\texttt{[00:21]}), widening his eyes and bringing both hands down in a sharp, parallel chopping gesture \texttt{\{\{gesture: bilateral vertical hand-chop on 'ghotala'\}\}}. When mentioning the political exploitation of faith (*"bhagwan Ram ke naam par... chanda chori"\texttt{, }[13:45]\texttt{), his voice drops in pitch with a slow, disapproving head shake }\{\{gesture: slow head shake + furrowed brow\}\}`, framing the act not as a minor clerical lapse but as calculated financial fraud.\par\smallskip
\end{itemize}

\noindent\textbf{2.} \textbf{Differential demeanour toward actors (Temple, Trust leadership, counting clerks, donors, opposition accusers): How do the speaker's voice pitch, facial expression, and posture shift when referring to Prabhu Shri Ram / devotees vs. top figures (Champat Rai, Nripendra Misra, CM Yogi) vs. arrested counting staff vs. opposition leaders (Akhilesh Yadav, Awadhesh Prasad)?}\par\smallskip
\begin{itemize}
  \item \textbf{Notes:}
Dhruv Rathee's non-verbal demeanour bifurcates sharply along lines of institutional power versus ordinary citizenry:\par\smallskip
  \item \textbf{Donors / Ordinary Citizens:} Addressed with empathetic warmth, softened volume, and an open palm extended toward the lens (\texttt{"gareeb logon ne apni mehnat ki kamai se sau-sau rupaye diye"}, \texttt{[01:05]}), positioning devotees as earnest victims of deception.
  \item \textbf{Opposition Whistleblowers (Sanjay Singh, Pawan Pandey):} Introduced with level, authoritative courtroom-style cadence (\texttt{[01:42]}), nodding approvingly as their press-conference documents are displayed full-screen.
  \item \textbf{Trust Leadership \& Political Affiliates (Champat Rai, Mayor Rishikesh Upadhyay, Anil Mishra):} Named with a narrowed gaze, raised eyebrows, and pointed index-finger gestures toward the on-screen relationship flowchart (\texttt{[04:15]}--\texttt{[05:30]}). When noting that Trust member Anil Mishra and BJP Mayor Rishikesh Upadhyay acted as official witnesses on BOTH the Rs 2 crore deed and the Rs 18.5 crore agreement five minutes later (\texttt{[04:48]}), Dhruv pauses \texttt{(1.2)}, leans forward, and taps the air twice with his index finger \texttt{\{\{gesture: double index-finger air-tap toward screen\}\}} to signal undeniable collusion.
\end{itemize}

\noindent\textbf{3.} \textbf{Solemn devotional grief vs. high-decibel partisan spectacle: How do pauses, sighs, speech rate, and hand gestures construct the segment's emotional register? Where is delivery slow, hushed, and mournful (sighs, bowed head, folded hands), and where does it escalate into rapid-fire shouting, desk-thumping, or mocking irony?}\par\smallskip
\begin{itemize}
  \item \textbf{Notes:}
Dhruv constructs a \textbf{measured, forensic-prosecutorial register} (\textasciitilde{}135--145 words per minute) that avoids both devotional weeping and studio shouting. His delivery is highly rhythmic and didactic: he speaks in clear, segmented clauses separated by deliberate \texttt{(0.5)} to \texttt{(1.0)} pauses so the viewer can read the highlighted registry PDFs on screen. The emotional register shifts into \textbf{biting, cold irony} during the rebuttal section (\texttt{[08:10]}--\texttt{[11:30]}): rather than shouting, he slows his speech rate down to \textasciitilde{}110 WPM when reading the Trust's excuses, arching his right eyebrow and letting out a quiet, incredulous \texttt{[scoff]} to make the official defence sound self-evidently absurd.\par\smallskip
\end{itemize}

\noindent\textbf{4.} \textbf{Non-verbal evaluation of financial scrutiny: When questions about unreceipted offerings, lack of frisking, or land deals are raised, how does the host/speaker react non-verbally (earnest nodding, leaning into the camera, open-palm appeals vs. eye-rolls, scoffs, dismissive hand waves, or incredulous rising pitch $\uparrow$)?}\par\smallskip
\begin{itemize}
  \item \textbf{Notes:}
Financial scrutiny is the centrepiece of the video and is performed with \textbf{intense pedagogical focus and visual demonstration}:\par\smallskip
  \item When comparing the official circle rate (Rs 5.80 crore) with the actual purchase price (Rs 2 crore) and the Trust's agreement price five minutes later (Rs 18.5 crore, \texttt{[02:50]}--\texttt{[04:10]}), Dhruv uses wide horizontal hand spans \texttt{\{\{gesture: hands held close together for '2 crore', then stretched wide apart for '18.5 crore'\}\}} to spatialise the 825\% markup (\texttt{"ek second mein daam \_nau\_ guna badh gaya$\uparrow$"}, \texttt{[03:38]}).
  \item Conversely, when evaluating the Trust's counter-claim that an earlier unregistered agreement from March 2021 justified the price (\texttt{[09:15]}), he reacts with a skeptical head tilt, an incredulous rising pitch \texttt{$\uparrow$}, and open-palm appeals (\texttt{"toh phir Trust ne seedha Kusum Pathak se zameen kyon$\uparrow$ nahin khareedi?!"}, \texttt{[09:42]}).
\end{itemize}

\noindent\textbf{5.} \textbf{Performing certainty vs. hedging and innuendo (pauses, stops, trailing off): Are allegations delivered with definitive falling intonation $\downarrow$, finger-counting, and pointing at on-screen evidence (CCTV, documents), or hedged with hesitant pauses (.), abrupt self-interruptions (---), throat-clearing [cough], and trailing sentences (...) that imply higher-level complicity without naming it?}\par\smallskip
\begin{itemize}
  \item \textbf{Notes:}
Dhruv performs \textbf{high epistemic certainty} on the documentary trail while using \textbf{strategic trailing questions (\texttt{...})} to implicate higher political protection:\par\smallskip
  \item \textbf{Forensic Certainty:} He counts out the four sequential steps of the land flip on his fingers (\texttt{[02:45]}--\texttt{[04:55]}), pointing directly at the highlighted timestamps (\texttt{7:10 PM} vs. \texttt{7:15 PM}), stamp duty numbers, and witness signatures on screen with definitive falling intonation \texttt{$\downarrow$} (*"Ye sab kuchh sarkari kagazon mein darj hai$\downarrow$"\texttt{, }[05:12]`).
  \item \textbf{Strategic Innuendo via Trailing Off:} When connecting Mayor Rishikesh Upadhyay and witness Anil Mishra to the broader ruling party ecosystem (\texttt{[06:20]}), he slows down and trails off: \textit{"Gawah kaun hain? Trust ke trustee Anil Mishra aur BJP ke Mayor Rishikesh Upadhyay... ab isse zyada saaf saboot kya chahiye ki iske peechhe kaun log..."} (\texttt{[06:32]}), accompanied by a knowing nod \texttt{\{\{gesture: slow knowing nod + raised eyebrows\}\}} that invites the viewer to conclude top-down BJP complicity.
\end{itemize}

\noindent\textbf{6.} \textbf{Turn-taking, interruptions, and acoustic treatment of denials: In debates/interviews: who receives calm, uninterrupted speaking time and sympathetic nods, and who is cut off mid-sentence (---), talked over, or muted? In solo videos: how are clips of official denials framed via reaction pauses, raised eyebrows, or sarcastic chuckles?}\par\smallskip
\begin{itemize}
  \item \textbf{Notes:}
In this solo documentary format, official denials by Champat Rai and Trust spokespersons are never played as unmediated authority; instead, they are \textbf{sandwiched inside a "Myth vs. Reality" prosecutorial frame}. Before presenting each of the Trust's four defences (\texttt{[08:15]}--\texttt{[12:10]}), Dhruv holds up a finger (\textit{"Ab suniye unka pehla bahana"\texttt{, }[08:20]\texttt{) with a slight smirk. After reading Champat Rai's statement that }"hum par toh }100\texttt{ saal se aarop lag rahe hain, Mahatma Gandhi ki hatya ka bhi aarop laga tha"\textit{ (}[11:40]\texttt{), Dhruv freezes for }(1.5)\texttt{ seconds in a deadpan reaction pause }\{\{gesture: 1.5s frozen stare + slow blink\}\}\texttt{, lets out a dry }[scoff]\texttt{, and says with falling disgust: }"Matlab (.) chande ke hisaab ka sawaal poochho (.) toh ye Mahatma Gandhi ki hatya le aate hain$\downarrow$"* (}[11:52]`).\par\smallskip
\end{itemize}

\noindent\textbf{7.} \textbf{Performing the boundary between religious devotion and institutional critique: How does the presenter use vocal tone, sighs, and body language to signal personal reverence for the Temple (or secular objectivity) while criticising---or defending---the Trust's financial management?}\par\smallskip
\begin{itemize}
  \item \textbf{Notes:}
Dhruv performs \textbf{secular constitutional objectivity combined with moral empathy for ordinary believers}. Rather than adopting devotional gestures (folded hands or chanting), he maintains an upright, journalistic posture throughout. However, at \texttt{[13:30]}--\texttt{[14:20]}, when distinguishing between genuine faith (\texttt{shraddha}) and the commercial/political instrumentalisation of religion (\texttt{dhandha}), his voice softens into a earnest, lower-register appeal: \textit{"Jin logon ne shraddha se (.) apni aastha se paise diye..."} (\texttt{[13:38]}), placing an open hand toward the viewer, before hardening his tone and furrowing his brow on \textit{"unko dharam ke naam par \_loota\_ ja raha hai$\downarrow$"} (\texttt{[13:48]}).\par\smallskip
\end{itemize}

\noindent\textbf{8.} \textbf{Self-presentation as the "guardian of public trust": How does the speaker use direct eye-gaze, posture, moral sighs, or emphatic gestures (pointing at the viewer, hands on heart/desk) to position themselves as the authentic protector of devotees or public accountability?}\par\smallskip
\begin{itemize}
  \item \textbf{Notes:}
Dhruv positions himself as the \textbf{fearless independent investigator ("fact-checker") who exposes what mainstream "Godi Media" hides}. At \texttt{[12:30]}--\texttt{[13:15]}, he lets out a heavy moral \texttt{[sigh]}, shakes his head at how TV channels are ignoring the registration documents, and points his index finger directly at the camera lens \texttt{\{\{gesture: direct index-finger point at lens\}\}}: \textit{"Mainstream media aapko ye kagaz \_kabhi\_ nahin dikhayega$\downarrow$ (\texttt{[12:42]}) ... isiliye is video ko har ek insaan tak share kijiye"} (\texttt{[15:10]}). His unblinking eye-gaze and rhythmic pointing construct a direct pact of accountability between creator and citizen.\par\smallskip
\end{itemize}

\subsubsection*{Apply the five non-verbal features (Event 1)}

\begin{longtable}{|p{0.31\textwidth}|p{0.31\textwidth}|p{0.31\textwidth}|}
\hline
Non-verbal feature & Timestamp + spoken anchor words + non-verbal notation & What it does in context (stance, framing, implication) \\ \hline \endhead
\textbf{1. Pauses, silences \& stops} & \texttt{[04:48]} \textit{"Aur sabse badi baat (1.2) dono agreements mein gawah (witness) \_bilkul\_ same hain$\downarrow$"} / \texttt{[11:48]} \texttt{[(1.5) deadpan reaction pause after Champat Rai quote]} & The \texttt{(1.2)} pause before revealing the identical witnesses builds forensic climax; the \texttt{(1.5)} silent stare after Champat Rai's quote frames the Trust's defence as absurd without needing a single shouted word. \\ \hline
\textbf{2. Trailing sentences (\texttt{...})} & \texttt{[06:32]} \textit{"Gawah hain Trustee Anil Mishra aur BJP Mayor Rishikesh Upadhyay... ab aap khud samajh sakte hain ki ye paisa kahan..."} & Leaves the final destination of the Rs 16.5-crore markup unspoken, prompting the viewer to infer organised political kickbacks while staying anchored to the witness names on the deed. \\ \hline
\textbf{3. Intonation, pitch \& volume} & \texttt{[03:38]} \textit{"Sirf \_PANCH\_ minute ke andar (.) zameen ka daam \_NAU\_ guna badh gaya$\uparrow$!"} (sharp rising pitch \texttt{$\uparrow$} + heavy stress on numerals) & Sharp prosodic contrast on \textit{"\_PANCH\_ minute"} vs. \textit{"\_NAU\_ guna"} dramatises the mathematical impossibility of natural price appreciation in 300 seconds. \\ \hline
\textbf{4. Coughs, sighs \& vocalisations} & \texttt{[08:22]} \texttt{[[scoff]]} \textit{"Ab suniye Trust ki taraf se kya safai di gayi hai..."} / \texttt{[12:35]} \texttt{[[sigh]]} \textit{"Par afsos ki baat ye hai ki TV media is par chup hai..."} & The \texttt{[scoff]} pre-frames the Trust's explanation as a flimsy excuse (\textit{"bahana"}), while the \texttt{[sigh]} performs civic sorrow over mainstream media's complicity. \\ \hline
\textbf{5. Hand gestures, face \& posture} & \texttt{[03:35]} \texttt{\{\{gesture: hands close together for '2 crore', then flung wide apart for '18.5 crore'\}\}} / \texttt{[04:50]} \texttt{\{\{gesture: double index-finger tap on witness names\}\}} & Spatialises the financial discrepancy visually for the audience and directs visual attention to the documentary proof (witness signatures) on screen. \\ \hline
\end{longtable}

\subsubsection*{Short summary (Event 1 Non-Verbal Synthesis)}

\begin{itemize}
  \item \textbf{What overall interpretation or emotional stance does the non-verbal delivery encourage?}
The non-verbal delivery encourages \textbf{moral indignation grounded in forensic certainty}. By pairing calm, courtroom-style document walkthroughs and finger-counting with incredulous pitch peaks \texttt{$\uparrow$} on the 5-minute price jump and dry scoffs at the Trust's denials, Dhruv guides the viewer to see the transaction as an open-and-shut case of institutional corruption exploiting religious faith.\par\smallskip
  \item \textbf{Whose voice or physical presence is acoustically/visually foregrounded, undermined, or silenced?}
\textbf{Documentary records (sale deeds, stamp duty receipts, witness signatures) and opposition whistleblowers} are visually foregrounded at full screen and acoustically reinforced via firm falling intonation \texttt{$\downarrow$}. \textbf{Trust General Secretary Champat Rai} is visually and acoustically undermined---his denials are introduced as \textit{"bahane"} with a smirk and followed by a \texttt{(1.5)} deadpan silence that highlights their evasiveness.\par\smallskip
  \item \textbf{What evidence supports your non-verbal reading? What complicates it (e.g., speaker baseline, format)?}
The systematic pairing of spatial hand gestures (narrow vs. wide span) and emphatic stress on time/money numerals (\texttt{5 minutes}, \texttt{2 crore}, \texttt{18.5 crore}, \texttt{[03:10]}--\texttt{[04:55]}) demonstrates how kinesics and prosody work together to construct an evidentiary indictment. What complicates this is \textbf{speaker baseline}: Dhruv Rathee uses this exact didactic cadence, finger-counting, and wide-eyed emphasis across all his explainer videos (whether on science, history, or politics), meaning his authoritative posture is part of his pedagogical brand, though the clustering of \texttt{[scoff]} vocalisations around Champat Rai is a marked evaluative shift.\par\smallskip
\end{itemize}

\bigskip\hrule\bigskip

\subsection*{Event 2. Students' protest regarding NEET paper leak (Non-Verbal \& Paralinguistic Layer)}

\begin{itemize}
  \item \textbf{Focus:} How do prosody, pauses, trailing sentences, vocalisations, and visual body language construct the legitimacy of student protesters, police use of force, and state accountability?
  \item \textbf{Outlet:} Dhruv Rathee (YouTube Channel) | \textbf{Date:} 16 June 2024
  \item \textbf{Video title:} The Dark Reality of NEET Scam | Exposed by Dhruv Rathee | \href{https://youtu.be/l-6Pm-112io}{https://youtu.be/l-6Pm-112io} | \textbf{Timestamp(s):} \texttt{[00:00]}--\texttt{[21:15]}
  \item \textbf{Analyst:} Baani Solanki (\texttt{2025114016}) | \textbf{Language:} Hinglish (Hindi in Roman script + English terms)
  \item \textbf{Format / Visual setup:} \texttt{[ ]} TV Panel Debate \texttt{[ ]} TV Bulletin \texttt{[x]} Solo YouTube Explainer \texttt{[ ]} Interview / Discussion
  \item \textbf{Primary speaker(s) observed:} Dhruv Rathee (solo studio presenter with full-screen statistical charts, FIR copies, student protest footage, and Alakh Pandey / educator interview clips)
\end{itemize}

\subsubsection*{Event-specific non-verbal questions (Event 2)}

\noindent\textbf{1.} \textbf{Prosodic \& gestural marking of participant labels ("students/bachche" vs. "protesters" vs. "upadravi/goons"): How do vocal tone, volume, facial expression, and hand gestures shift across labels? (e.g., softened, empathetic pitch for "students" vs. harsh, staccato stress, scowls, or accusatory finger-pointing on "upadraviyon", "goons", or "anti-social elements".)}\par\smallskip
\begin{itemize}
  \item \textbf{Notes:}
Dhruv exclusively uses legitimising, humanising labels---\textit{"students"}, \textit{"masoom bachche"}, \textit{"desh ka bhavishya"}, and \textit{"doctors"} (\texttt{[00:45]}, \texttt{[02:15]}, \texttt{[18:30]})---delivered with a \textbf{softened, lower-pitched empathetic register}, furrowed brow of concern, and open palms held toward his chest \texttt{\{\{gesture: open palms drawn toward chest\}\}}. Conversely, harsh staccato stress, scowls, and accusatory finger-pointing are reserved entirely for the institutional and criminal actors: \textit{"education \_MAFIA\_"}, \textit{"NTA ki \_laaparwahi\_"}, and \textit{"corrupt \_system\_"} (\texttt{[09:20]}, \texttt{[14:45]}, \texttt{[19:10]}), where he points sharply downward and strikes his palm with his fist \texttt{\{\{gesture: right fist striking open left palm\}\}}.\par\smallskip
\end{itemize}

\noindent\textbf{2.} \textbf{Delivery of students' stated demands: When the protesters' core demands (Education Minister's resignation, compensation, release of Sonam Wangchuk, exam reform) are mentioned, are they delivered slowly with deliberate pauses and finger-counting, rushed through in a flat monotone, mocked with air-quotes/scoffs, or trailed off (...)?}\par\smallskip
\begin{itemize}
  \item \textbf{Notes:}
In the concluding section (\texttt{[17:40]}--\texttt{[19:50]}), Dhruv articulates the students' core demands---(1) immediate removal/accountability of NTA leadership, (2) independent investigation into the Bihar and Gujarat paper leaks, and (3) complete transparency in raw scores and grace-mark formulas---\textbf{slowly, deliberately, and with structured finger-counting}. He pauses \texttt{(0.8)} between each point, looks directly into the lens without blinking, and delivers each demand with resolute falling intonation \texttt{$\downarrow$} (*"Ye bachchon ki koi ehsaan ki maang nahin hai (.) ye unka \_haq\_ hai$\downarrow$"\texttt{, }[18:15]`).\par\smallskip
\end{itemize}

\noindent\textbf{3.} \textbf{Where does vocal \& kinetic energy peak---grievance, street disruption, or electoral horse-race?: Track where the speaker's volume, pitch range, pace, and gesture intensity reach their climax: during accounts of the paper leak/student suicides, during street clashes and barricade breaches, or during partisan BJP-vs-Opposition sparring.}\par\smallskip
\begin{itemize}
  \item \textbf{Notes:}
Vocal and kinetic energy reaches two distinct climaxes, both centred on \textbf{student grievance and systemic betrayal}:\par\smallskip
  \item \textbf{Emotional/Sombre Climax (\texttt{[00:00]}--\texttt{[02:10]}):} Opening with the human cost of NEET preparation and student suicides, where Dhruv speaks in a hushed, slow, deeply sombre cadence (\texttt{$^\circ$ek-ek number ke liye bachche raat-raat bhar jagte hain$^\circ$}, \texttt{[01:12]}) with a heavy \texttt{[sigh]} and pained facial expression.
  \item \textbf{Indignant/Forensic Climax (\texttt{[10:15]}--\texttt{[13:50]}):} When connecting the Patna FIR (\texttt{5 May}), the Godhra cheating centre (\texttt{Rs 10 lakh} per candidate), and the Education Minister's flat denial, his volume rises sharply, his pace accelerates, and his hand gestures become rapid and forceful (\texttt{"Jab Bihar Police ki FIR mein saaf likha hai (.) toh Education Minister kis munh se keh rahe hain ki \_koi\_ leak nahin hua$\uparrow$?!"}, \texttt{[12:28]}).
\end{itemize}

\noindent\textbf{4.} \textbf{Non-verbal attribution of responsibility (sighs, stops, head shakes, pointing): When examination authorities (NTA), the Education Minister, police, opposition politicians, or students are named as responsible, what non-verbal cues accompany the attribution (exasperated sighs, mocking laughs, abrupt self-stops ---, or direct finger-pointing at the lens)?}\par\smallskip
\begin{itemize}
  \item \textbf{Notes:}
Responsibility is attributed directly to the \textbf{NTA chairperson, Education Minister Dharmendra Pradhan, and the ruling government's pattern of paper leaks (\texttt{70+} leaks in \texttt{7} years, \texttt{[15:20]})}:\par\smallskip
  \item When citing the fact that the NTA formed a committee to investigate the NTA---headed by the UPSC chair/insiders (\texttt{[08:35]})---Dhruv stops mid-sentence \texttt{---}, lets out a mocking \texttt{[chuckle]}, shakes his head, and says: \textit{"Chor hi (.) chori ki jaanch karega$\downarrow$"} (\texttt{[08:42]}) with a circular hand gesture indicating a closed loop \texttt{\{\{gesture: both index fingers rotating in a circle\}\}}.
  \item When addressing the Education Minister's initial denial (\texttt{[12:15]}), Dhruv points his index finger directly at the camera lens \texttt{\{\{gesture: sustained index-finger point at lens\}\}} with a stern, unsmiling scowl.
\end{itemize}

\noindent\textbf{5.} \textbf{Prosody of questions---open inquiry vs. loaded rhetorical challenge: Are questions about the protest delivered with level, open intonation, or with sharp rising pitch $\uparrow$, sarcastic stress, raised eyebrows, and trailing clauses (...) that presuppose protesters are violent or politically manipulated?}\par\smallskip
\begin{itemize}
  \item \textbf{Notes:}
Every single loaded rhetorical challenge in the video is aimed \textbf{upward at state authorities}, never downward at protesting students:\par\smallskip
  \item At \texttt{[05:40]}, questioning the \texttt{4 June} result timing: \textit{"Result \texttt{14} June ko aana tha, par achanak \texttt{4} June---Election Result ke din---kyon$\uparrow$ declare kiya gaya?! Kya NTA kuchh \_chhupana\_ chahti thi$\uparrow$?"} (\texttt{[05:48]}), delivered with rising pitch \texttt{$\uparrow$}, narrowed eyes, and a forward lean.
  \item By contrast, questions about the protesting students (\texttt{"Aaj jab ye bachche dhoop mein sadak par khade hain, toh inki sunne wala kaun hai$\downarrow$"}, \texttt{[16:40]}) use a falling, mournful contour \texttt{$\downarrow$} and open-palm empathy gestures that presuppose the complete legitimacy of their protest.
\end{itemize}

\noindent\textbf{6.} \textbf{Visual \& paralinguistic performance of injury and violence: How are claims of violence substantiated non-verbally on screen? (Compare camera close-ups and pained vocal tone on police bandages/fractured hands vs. visual/acoustic treatment of student head injuries, lathi-charge footage, or defensive stammers/pauses when an officer or guest is asked who used force first.)}\par\smallskip
\begin{itemize}
  \item \textbf{Notes:}
Dhruv substantiates the violence of the crisis through \textbf{close-up visuals of crying students, anguished parents, student suicide headlines, and police roughing up peaceful demonstrators} (\texttt{[01:20]}, \texttt{[16:30]}--\texttt{[17:15]}). While showing footage of police dragging and detaining protesting students (\texttt{[16:45]}), his voice drops in pitch and takes on an indignant, pained edge: \textit{"Jo bachche apna haq maangne aaye (.) unke saath \_aisa\_ salook kiya ja raha hai$\downarrow$"} (\texttt{[16:52]}), accompanied by a slow, disgusted head shake \texttt{\{\{gesture: slow head shake + furrowed brow\}\}}. Zero screen time is given to police victimhood narratives.\par\smallskip
\end{itemize}

\noindent\textbf{7.} \textbf{Asymmetric delivery of police actions vs. protester actions: Compare the vocal cadence, pauses, and hand gestures used when describing police actions (e.g., calm, matter-of-fact administrative tone on "tear gas chhode gae" / "preventive arrangement") versus protester actions (alarmist pitch, striking gestures, or ironic softening as "retaliation").}\par\smallskip
\begin{itemize}
  \item \textbf{Notes:}
Dhruv distinguishes sharply between \textbf{investigative state police (Bihar EOU)} and \textbf{street crackdown policing}:\par\smallskip
  \item \textbf{Bihar Police's EOU investigation} (\texttt{[10:20]}--\texttt{[11:40]}) is narrated with crisp, authoritative approval, pointing to the recovered burnt question-paper booklets and post-dated cheques as decisive forensic proof.
  \item \textbf{Street repression of protesters} (\texttt{[16:45]}) is narrated with moral condemnation and downward-pointing gestures, framing state force against unarmed \texttt{17}- and \texttt{18}-year-old aspirants as authoritarian bullying.
  \item \textbf{Student protesters' actions} (marching, chanting, filing petitions with educators like Alakh Pandey, \texttt{[06:40]}, \texttt{[16:30]}) are narrated with uplifting, solidaristic prosody and approving nods.
\end{itemize}

\noindent\textbf{8.} \textbf{Acoustic \& visual space granted to students' own voices: Are protesting students heard in their own voice with clear audio, steady close-ups, and empathetic listening nods, or are they reduced to noisy background chaos B-roll while studio anchors/politicians speak over them?}\par\smallskip
\begin{itemize}
  \item \textbf{Notes:}
Students and their educator-advocates (such as PhysicsWallah's Alakh Pandey and protesting NEET candidates) are given \textbf{prominent, clean, unmuted acoustic and visual space} (\texttt{[01:30]}, \texttt{[06:45]}, \texttt{[16:20]}). Their clips play at full screen with clear subtitles and zero studio cross-talk. Immediately after a clip of a distraught student saying \textit{"650 marks laane ke baad bhi college nahin milega"} (\texttt{[01:42]}), Dhruv pauses \texttt{(1.0)} in silence with a solemn, empathetic expression before speaking softly (\texttt{$^\circ$soch kar dekhiye us bachche par kya beet rahi hogi$^\circ$}, \texttt{[01:46]}).\par\smallskip
\end{itemize}

\noindent\textbf{9.} \textbf{Non-verbal shift between the NEET paper-leak issue and the street protest: Does the speaker's non-verbal demeanour shift from sympathetic concern (sombre tone, sighs, solidarity gestures) when discussing the exam leak itself to combative, authoritarian, or dismissive delivery once students march toward Parliament?}\par\smallskip
\begin{itemize}
  \item \textbf{Notes:}
\textbf{No hostile shift occurs toward the students.} Rather, as the video moves from the technical exam anomalies (\texttt{[02:30]}--\texttt{[09:00]}) to the Bihar/Gujarat paper-leak mafia (\texttt{[09:30]}--\texttt{[14:30]}) and finally to the nationwide protests (\texttt{[16:00]}--\texttt{[20:30]}), Dhruv's non-verbal demeanour escalates from \textbf{analytical concern to prosecutorial solidarity}. He explicitly frames the street protests and Supreme Court petitions as the only remaining democratic check against an unaccountable government (\texttt{[17:20]}).\par\smallskip
\end{itemize}

\subsubsection*{Apply the five non-verbal features (Event 2)}

\begin{longtable}{|p{0.31\textwidth}|p{0.31\textwidth}|p{0.31\textwidth}|}
\hline
Non-verbal feature & Timestamp + spoken anchor words + non-verbal notation & What it does in context (stance, framing, implication) \\ \hline \endhead
\textbf{1. Pauses, silences \& stops} & \texttt{[01:45]} \texttt{[(1.0) solemn silence after distraught student clip]} \textit{"$^\circ$Soch kar dekhiye (.) us bachche par kya beet rahi hogi$^\circ$"} / \texttt{[08:40]} \textit{"NTA ne jaanch ke liye committee banayi---aur uska head kaun hai? (1.0) Khud NTA$\downarrow$"} & The \texttt{(1.0)} silence after the student clip builds emotional empathy; the \texttt{(1.0)} pause before \textit{"Khud NTA$\downarrow$"} delivers a sharp satirical punchline exposing conflict of interest. \\ \hline
\textbf{2. Trailing sentences (\texttt{...})} & \texttt{[05:52]} \textit{"Chaar June---jis din poore desh ka dhyan Election Results par tha---usi din achanak NEET ka result... taaki kisi ka dhyan..."} & Trails off with a narrowed gaze and slow nod, leading the audience to complete the thought that the NTA deliberately buried the scandal under election coverage. \\ \hline
\textbf{3. Intonation, pitch \& volume} & \texttt{[04:15]} \textit{"\texttt{718} aur \texttt{719} marks---jo mathematically \_IMPOSSIBLE\_ hain$\uparrow$!"} / \texttt{[12:28]} \textit{"Kis munh se keh rahe hain ki \_KOI\_ paper leak nahin hua$\uparrow$?!"} & High-volume stress on \texttt{\_IMPOSSIBLE\_} and \texttt{\_KOI\_} combined with sharp rising pitch \texttt{$\uparrow$} transforms technical audit points into moral indictments of state dishonesty. \\ \hline
\textbf{4. Coughs, sighs \& vocalisations} & \texttt{[01:10]} \texttt{[[sigh]]} \textit{"Har saal (\texttt{24} lakh) bachche is exam mein baithte hain..."} / \texttt{[08:42]} \texttt{[[chuckle / scoff]]} \textit{"Matlab chor hi chori ki jaanch karega$\downarrow$"} & The heavy \texttt{[sigh]} establishes paternal/fraternal sorrow for students; the ironic \texttt{[chuckle / scoff]} strips the NTA's internal inquiry committee of any institutional credibility. \\ \hline
\textbf{5. Hand gestures, face \& posture} & \texttt{[08:42]} \texttt{\{\{gesture: both index fingers rotating in a closed circle\}\}} / \texttt{[12:30]} \texttt{\{\{gesture: direct index-finger point at lens + stern scowl\}\}} / \texttt{[18:10]} \texttt{\{\{gesture: sequential counting on 3 fingers\}\}} & The circular finger motion visualises institutional self-dealing; the direct lens point confronts the Education Minister; finger-counting organises student demands into a rational charter. \\ \hline
\end{longtable}

\subsubsection*{Short summary (Event 2 Non-Verbal Synthesis)}

\begin{itemize}
  \item \textbf{What overall interpretation or emotional stance does the non-verbal delivery encourage?}
The non-verbal delivery encourages \textbf{deep emotional solidarity with NEET aspirants alongside righteous civic anger at systemic state failure}. By opening with hushed, sigh-laden empathy for suicidal and exhausted students and building toward loud, finger-pointing indictment of the NTA and Education Ministry, the video mobilises the viewer to view the NEET crisis as symptomatic of broader institutional rot.\par\smallskip
  \item \textbf{Whose voice or physical presence is acoustically/visually foregrounded, undermined, or silenced?}
\textbf{Students, parents, educator-activists (Alakh Pandey), and Bihar Police FIR documents} are visually and acoustically foregrounded with full-screen clarity and respectful pauses. \textbf{Education Minister Dharmendra Pradhan and the NTA Director} are acoustically and visually undermined via scoffs, circular "self-investigation" hand gestures, and immediate juxtaposition with contradicting police evidence.\par\smallskip
  \item \textbf{What evidence supports your non-verbal reading? What complicates it (e.g., speaker baseline, format)?}
The stark acoustic contrast between the softened, low-pitch register (\texttt{$^\circ$word$^\circ$}, \textasciitilde{}115 WPM) used for student suffering (\texttt{[01:00]}--\texttt{[02:10]}) and the high-volume, staccato prosecutorial register (\textasciitilde{}150 WPM, \texttt{WORD}, \texttt{$\uparrow$}) used for the Education Minister's denials (\texttt{[12:15]}--\texttt{[13:00]}) proves that prosody is systematically deployed to assign victimhood and blame. A complicating baseline factor is that Dhruv's prosecutorial cadence, dramatic pauses, and direct camera pointing are consistent across all his political critique videos, making it a signature genre style of oppositional YouTube explainers.\par\smallskip
  \item \textbf{How do the non-verbal cues interact with the spoken words? (Do they reinforce, dramatise, soften/hedge, invert meaning via sarcasm/irony, or leave an accusation unspoken via trailing off?)}
The non-verbal cues \textbf{dramatise and reinforce} the documentary evidence (pointing at FIRs and scorecards), \textbf{invert official assurances via sarcasm} (the \texttt{[scoff]} and circular finger gesture on NTA's inquiry committee at \texttt{[08:42]}), and \textbf{leave the accusation of top-down political orchestration unspoken via trailing off} (\texttt{...}) when discussing the \texttt{4 June} election-day release (\texttt{[05:52]}) and the Gujarat/Bihar exam mafia nexus (\texttt{[14:50]}).\par\smallskip
  \item \textbf{What context, speaker baseline, or format factors support or complicate this interpretation? (e.g., Is this the speaker's normal style across all videos, or a marked shift on this specific actor/claim?)}
Compared to his Event 1 video (where he maintained a brisk, forensic cadence throughout), Event 2 displays a \textbf{markedly more emotional opening and closing baseline}, using longer silences \texttt{(1.0)} and softer vocal volume \texttt{$^\circ$...$^\circ$} to honour student trauma before shifting into his standard prosecutorial rhythm.\par\smallskip
  \item \textbf{How does another outlet perform the same moment, actor, or claim non-verbally?}
On \textbf{Republic TV} (\texttt{Event 2}), anchor Arnab Goswami also shouts at the NTA, but whenever the opposition or street protests are mentioned, Arnab pivots into frantic desk-thumping and rapid-fire interruptions \texttt{[overlap]} to attack opposition politicians (\textit{"political vultures"}). In contrast, \textbf{Aaj Tak} (\texttt{Event 2}) uses stern finger-wagging and rising pitch \texttt{$\uparrow$} to interrogate students over broken barricades and injured police. Dhruv Rathee never once turns his prosecutorial gaze against the protesters; his kinetic energy remains 100\% aligned with the students from start to finish.\par\smallskip
  \item \textbf{My interpretation, with reasons:}
Dhruv Rathee's non-verbal performance illustrates \textbf{multimodal counter-framing}: in the absence of a live studio panel, he uses prosodic contrast (hushed empathy vs. sharp prosecutorial stress), visual spatialisation (wide hand spans for price/rank jumps), and strategic reaction pauses \texttt{(1.5)} to turn official government statements into self-incriminating exhibits across both the Ram Mandir land controversy and the NEET paper leak.\par\smallskip
\end{itemize}
\newpage
\section*{Non-Verbal \& Multimodal Discourse Analysis Worksheet}

\textit{Qualitative analysis of pauses, trailing sentences, prosody/intonation, vocalisations (coughs, sighs, scoffs), and kinesics (hand gestures, facial expressions, posture) in political news discourse}\par\smallskip

\bigskip\hrule\bigskip

\subsection*{Guidelines: Five non-verbal \& paralinguistic features to look for in every event}

\begin{longtable}{|p{0.23\textwidth}|p{0.23\textwidth}|p{0.23\textwidth}|p{0.23\textwidth}|}
\hline
Feature & What to look for & Illustrative example (from sample) & Questions to answer \\ \hline \endhead
\textbf{1. Pauses, Silences \& Abrupt Stops} & Micro-pauses \texttt{(.)}, timed silences \texttt{(1.5)}, pregnant/dramatic pauses before a loaded word, hesitant stumbles, abrupt self-interruptions or cut-offs (\texttt{---}), and forced silences (muting a panelist's mic or talking over them). & Anchor pauses \texttt{(1.5)} before saying \textit{"'loota'"} for dramatic weight; Addl. DCP cuts off mid-sentence (\textit{"wo toh--- abhi investigation chal raha hai"}) when asked if police struck first; panelist's mic faded during counter-question. & Where does the pause or stop occur? Is it deliberate/rhetorical (for suspense, emphasis, or gravitas) or defensive/evasive (cognitive load, self-repair, avoiding blame)? Who controls turn-taking and silence? \\ \hline
\textbf{2. Trailing Sentences \& Unfinished Utterances} & Clauses that trail off (\texttt{...}), drop in volume at the end, or leave the conclusion unspoken (aposiopesis) so the audience supplies the implication without the speaker taking explicit responsibility. & \textit{"Aur phir wahan jo log ghuse the... aap samajh hi sakte hain..."}; \textit{"500 salon ke sangharsh ke baad agar aisa..."} trailing off with a slow head shake. & What exact thought is left unfinished? What presupposition or shared bias is the viewer invited to fill in? Does trailing off act as a safe hedge to imply higher-level complicity or "anti-national" motive without stating it? \\ \hline
\textbf{3. Intonation, Pitch, Volume \& Prosody} & Pitch contour (rising \texttt{$\uparrow$} for incredulity/challenge; falling \texttt{$\downarrow$} for finality), volume shifts (shouting vs. hushed whisper), emphatic word stress, speech rate (rapid-fire urgency vs. slow solemnity), and tonal register (sarcastic, reverential, combative). & Hushed, reverential cadence on \textit{"Prabhu Shri Ram ke Mandir"} vs. sharp, high-pitched interrogation on \textit{"par Akhilesh ji Ram Mandir kyon$\uparrow$ nahin jate?!"}; satirical sing-song inflection on \textit{"priventiv arenjment"}. & Which exact words receive pitch peaks or volume stress? How does vocal tone turn neutral words into irony, mockery, or moral outrage? Does the speaker shift prosodic register depending on which actor is discussed? \\ \hline
\textbf{4. Coughs, Sighs \& Non-Lexical Vocalisations} & Audible sighs (\texttt{[sigh]}), scoffs/snorts (\texttt{[scoff]}), ironic chuckles (\texttt{[chuckle]}), throat-clearing/coughs (\texttt{[cough]}), sharp intakes of breath/gasps (\texttt{[gasp]}), and tongue-clicks (\texttt{[tsk]}). & Heavy \texttt{[sigh]} before reporting Trust oversight lapses (signalling reluctant sorrow rather than hostility); mocking \texttt{[scoff]} or \texttt{[chuckle]} when reading an opponent's denial; \texttt{[cough / throat-clear]} when pivoting away from an awkward fact. & What emotion or epistemic stance does the vocalisation perform (exasperation, moral disgust, disbelief, ridicule, discomfort)? Does a sigh or scoff pre-judge a speaker's claim before any verbal rebuttal is made? \\ \hline
\textbf{5. Hand Gestures, Facial Expressions \& Posture (Kinesics)} & Hand movements (air-quotes, finger-pointing/wagging, desk-thumping, open palms, counting on fingers, dismissive wave), facial cues (raised eyebrows, eye-rolls, smirk, furrowed brow), head nods/shakes, leaning into camera, and visual props (bandaged limbs, documents). & Anchor forms physical air-quotes on \textit{"'loota'"} or \textit{"'peaceful' protest"}; host points finger/pen at lens during an ultimatum; injured DCP holds bandaged hand up to camera; creator raises one eyebrow and leans in for a conspiratorial reveal. & How do gestures reinforce, contradict, or sharpen the spoken words? During split-screen listening, does the host's face (nodding vs. smirking/eye-rolling) cue the audience on whom to trust? How are bodily visuals used as evidence? \\ \hline
\end{longtable}

\subsection*{Quick Multimodal Transcription Notation Key (use inside transcript excerpts)}

\begin{longtable}{|p{0.15\textwidth}|p{0.15\textwidth}|p{0.15\textwidth}|p{0.15\textwidth}|p{0.15\textwidth}|p{0.15\textwidth}|}
\hline
Symbol & Meaning & Example in transcript & Symbol & Meaning & Example in transcript \\ \hline \endhead
\texttt{(.)} & Brief micro-pause (< 0.5 sec) & \texttt{woh (.) suparavijan lacking tha} & \texttt{$\uparrow$ / $\downarrow$} & Marked rising / falling pitch & \texttt{kyon$\uparrow$ FIR nahin karaee$\downarrow$} \\ \hline
\texttt{(1.5)} & Timed pause in seconds & \texttt{kisne (1.5) 'loota'?} & \texttt{WORD / \_word\_} & Loud volume / emphatic stress & \texttt{ye PAAP hai / \_bilkul\_ galat} \\ \hline
\texttt{---} & Abrupt stop, cut-off, or self-repair & \texttt{police ki taraf se--- abhi jaanch...} & \texttt{$^\circ$word$^\circ$} & Softened / whispered delivery & \texttt{$^\circ$ander ki baat ye hai$^\circ$} \\ \hline
\texttt{...} & Trailing off / unfinished sentence & \texttt{aap jaante hain ki kaun log...} & \texttt{[sigh] [cough] [scoff]} & Non-lexical vocalisations & \texttt{[sigh] 500 salon ke baad...} \\ \hline
\texttt{[overlap]} & Simultaneous / overlapping speech & \texttt{[Anjana: suniye---] [Guest: nahin!]} & \texttt{\{\{gesture: ...\}\}} & Hand gesture / facial cue / posture & \texttt{\{\{air-quotes on 'loota'\}\}} \\ \hline
\end{longtable}

\textbf{Methodological Reminder:} A non-verbal cue is not automatically manipulative or ideological---a cough, pause, or rapid cadence may reflect studio audio delay, live time pressure, or a speaker's habitual style. Always check (1) whether the cue clusters with verbal features, and (2) whether the speaker's delivery shifts noticeably across different actors or events (within-outlet baseline).\par\smallskip

\bigskip\hrule\bigskip

\subsection*{Event 1. Ram Mandir donation controversy (Non-Verbal \& Paralinguistic Layer)}

\begin{itemize}
  \item \textbf{Focus:} How do prosody, pauses, trailing sentences, vocalisations (sighs, coughs, scoffs), and body language frame the controversy---as solemn devotional sorrow, administrative scrutiny, or partisan combat?
  \item \textbf{Outlet:} Mohak Mangal (\textit{Soch by Mohak Mangal}) | \textbf{Date:} N/A (Skipped per project guidelines)
  \item \textbf{Video title:} \textit{Skipped --- No video relevant to the Ram Mandir donation/land scam; channel only published a historical explainer on the Ayodhya dispute} | \textbf{Timestamp(s):} N/A
  \item \textbf{Analyst:} Baani Solanki (\texttt{2025114016}) | \textbf{Language:} N/A (Skipped)
  \item \textbf{Format / Visual setup:} \texttt{[ ]} TV Panel Debate \texttt{[ ]} TV Bulletin \texttt{[x]} Solo YouTube Explainer \texttt{[ ]} Interview / Discussion
  \item \textbf{Primary speaker(s) observed:} N/A (Skipped for Event 1 per official project guidelines: \textit{"Mohak mangal - no video relevant to scam; there is video regarding history, so skipping"})
\end{itemize}

\subsubsection*{Event-specific non-verbal questions (Event 1)}

\noindent\textbf{1.} \textbf{Vocal \& gestural marking of core labels ("scam", "loota", "paap", "gadabdi", "avyavastha"): When the controversy is named, what non-verbal cues accompany the label? (Look for physical air-quotes, vocal scare-quotes, dramatic pauses before the word, hushed/reverential drops in volume, or shouted stress.)}\par\smallskip
\begin{itemize}
  \item \textbf{Notes:}
\textbf{Skipped per project instructions.} From a discourse-analytic standpoint, Mohak Mangal's non-coverage of the 2021/2024 Ram Mandir donation and land-deal allegations constitutes a \textbf{structural silence (editorial gatekeeping)}: unlike Dhruv Rathee (who foregrounds day-to-day political/corruption exposés) or TV news channels (which run daily partisan debates), \textit{Soch by Mohak Mangal} avoids acute partisan flashpoints involving religious institutions unless they crystallise into broader policy or structural case studies.\par\smallskip
\end{itemize}

\noindent\textbf{2.} \textbf{Differential demeanour toward actors (Temple, Trust leadership, counting clerks, donors, opposition accusers): How do the speaker's voice pitch, facial expression, and posture shift when referring to Prabhu Shri Ram / devotees vs. top figures (Champat Rai, Nripendra Misra, CM Yogi) vs. arrested counting staff vs. opposition leaders (Akhilesh Yadav, Awadhesh Prasad)?}\par\smallskip
\begin{itemize}
  \item \textbf{Notes:}
\textbf{Skipped per project instructions} (no video covering Event 1).\par\smallskip
\end{itemize}

\noindent\textbf{3.} \textbf{Solemn devotional grief vs. high-decibel partisan spectacle: How do pauses, sighs, speech rate, and hand gestures construct the segment's emotional register? Where is delivery slow, hushed, and mournful (sighs, bowed head, folded hands), and where does it escalate into rapid-fire shouting, desk-thumping, or mocking irony?}\par\smallskip
\begin{itemize}
  \item \textbf{Notes:}
\textbf{Skipped per project instructions} (no video covering Event 1).\par\smallskip
\end{itemize}

\noindent\textbf{4.} \textbf{Non-verbal evaluation of financial scrutiny: When questions about unreceipted offerings, lack of frisking, or land deals are raised, how does the host/speaker react non-verbally (earnest nodding, leaning into the camera, open-palm appeals vs. eye-rolls, scoffs, dismissive hand waves, or incredulous rising pitch $\uparrow$)?}\par\smallskip
\begin{itemize}
  \item \textbf{Notes:}
\textbf{Skipped per project instructions} (no video covering Event 1).\par\smallskip
\end{itemize}

\noindent\textbf{5.} \textbf{Performing certainty vs. hedging and innuendo (pauses, stops, trailing off): Are allegations delivered with definitive falling intonation $\downarrow$, finger-counting, and pointing at on-screen evidence (CCTV, documents), or hedged with hesitant pauses (.), abrupt self-interruptions (---), throat-clearing [cough], and trailing sentences (...) that imply higher-level complicity without naming it?}\par\smallskip
\begin{itemize}
  \item \textbf{Notes:}
\textbf{Skipped per project instructions} (no video covering Event 1).\par\smallskip
\end{itemize}

\noindent\textbf{6.} \textbf{Turn-taking, interruptions, and acoustic treatment of denials: In debates/interviews: who receives calm, uninterrupted speaking time and sympathetic nods, and who is cut off mid-sentence (---), talked over, or muted? In solo videos: how are clips of official denials framed via reaction pauses, raised eyebrows, or sarcastic chuckles?}\par\smallskip
\begin{itemize}
  \item \textbf{Notes:}
\textbf{Skipped per project instructions} (no video covering Event 1).\par\smallskip
\end{itemize}

\noindent\textbf{7.} \textbf{Performing the boundary between religious devotion and institutional critique: How does the presenter use vocal tone, sighs, and body language to signal personal reverence for the Temple (or secular objectivity) while criticising---or defending---the Trust's financial management?}\par\smallskip
\begin{itemize}
  \item \textbf{Notes:}
\textbf{Skipped per project instructions} (no video covering Event 1).\par\smallskip
\end{itemize}

\noindent\textbf{8.} \textbf{Self-presentation as the "guardian of public trust": How does the speaker use direct eye-gaze, posture, moral sighs, or emphatic gestures (pointing at the viewer, hands on heart/desk) to position themselves as the authentic protector of devotees or public accountability?}\par\smallskip
\begin{itemize}
  \item \textbf{Notes:}
\textbf{Skipped per project instructions} (no video covering Event 1).\par\smallskip
\end{itemize}

\subsubsection*{Apply the five non-verbal features (Event 1)}

\begin{longtable}{|p{0.31\textwidth}|p{0.31\textwidth}|p{0.31\textwidth}|}
\hline
Non-verbal feature & Timestamp + spoken anchor words + non-verbal notation & What it does in context (stance, framing, implication) \\ \hline \endhead
\textbf{1. Pauses, silences \& stops} & \texttt{N/A (Event 1 skipped per project guidelines)} & Editorial omission / structural silence on acute religious-trust financial allegations; channel focuses on long-horizon policy and sociological explainers (see Event 2). \\ \hline
\textbf{2. Trailing sentences (\texttt{...})} & \texttt{N/A (Event 1 skipped per project guidelines)} & N/A \\ \hline
\textbf{3. Intonation, pitch \& volume} & \texttt{N/A (Event 1 skipped per project guidelines)} & N/A \\ \hline
\textbf{4. Coughs, sighs \& vocalisations} & \texttt{N/A (Event 1 skipped per project guidelines)} & N/A \\ \hline
\textbf{5. Hand gestures, face \& posture} & \texttt{N/A (Event 1 skipped per project guidelines)} & N/A \\ \hline
\end{longtable}

\subsubsection*{Short summary (Event 1 Non-Verbal Synthesis)}

\begin{itemize}
  \item \textbf{What overall interpretation or emotional stance does the non-verbal delivery encourage?}
\textbf{Skipped per project guidelines.} Mohak Mangal did not produce a video on the Ram Mandir donation/land controversy.\par\smallskip
  \item \textbf{Whose voice or physical presence is acoustically/visually foregrounded, undermined, or silenced?}
\textbf{Skipped per project guidelines.}\par\smallskip
  \item \textbf{What evidence supports your non-verbal reading? What complicates it (e.g., speaker baseline, format)?}
\textbf{Skipped per project guidelines.} Comparing Mohak Mangal's absence in Event 1 with his deep-dive policy explainer in Event 2 confirms his channel's editorial baseline: avoiding breaking-news partisan allegations around religious bodies while investing heavily in educational, sociological, and institutional-reform topics affecting young Indians.\par\smallskip
\end{itemize}

\bigskip\hrule\bigskip

\subsection*{Event 2. Students' protest regarding NEET paper leak (Non-Verbal \& Paralinguistic Layer)}

\begin{itemize}
  \item \textbf{Focus:} How do prosody, pauses, trailing sentences, vocalisations, and visual body language construct the legitimacy of student protesters, police use of force, and state accountability?
  \item \textbf{Outlet:} Mohak Mangal (\textit{Soch by Mohak Mangal}) | \textbf{Date:} 14 June 2024
  \item \textbf{Video title:} NTA: The Biggest Scam in Indian Education? | NEET 2024 | Mohak Mangal | \href{https://youtu.be/v-RikSqb0Vw}{https://youtu.be/v-RikSqb0Vw} | \textbf{Timestamp(s):} \texttt{[00:00]}--\texttt{[20:12]}
  \item \textbf{Analyst:} Baani Solanki (\texttt{2025114016}) | \textbf{Language:} Hinglish (Hindi in Roman script + English terms)
  \item \textbf{Format / Visual setup:} \texttt{[ ]} TV Panel Debate \texttt{[ ]} TV Bulletin \texttt{[x]} Solo YouTube Explainer \texttt{[ ]} Interview / Discussion
  \item \textbf{Primary speaker(s) observed:} Mohak Mangal (solo studio presenter seated/standing in front of warm minimalist studio lighting, integrated with data visualisations, animated score graphs, and clips of NEET students and NTA press conferences)
\end{itemize}

\subsubsection*{Event-specific non-verbal questions (Event 2)}

\noindent\textbf{1.} \textbf{Prosodic \& gestural marking of participant labels ("students/bachche" vs. "protesters" vs. "upadravi/goons"): How do vocal tone, volume, facial expression, and hand gestures shift across labels? (e.g., softened, empathetic pitch for "students" vs. harsh, staccato stress, scowls, or accusatory finger-pointing on "upadraviyon", "goons", or "anti-social elements".)}\par\smallskip
\begin{itemize}
  \item \textbf{Notes:}
Mohak Mangal never uses derogatory labels (\textit{"upadravi"}, \textit{"goons"}). He refers to candidates by personal names---such as \texttt{"Ayushi"} (\texttt{[00:15]}) and \texttt{"Anjali"} (\texttt{[01:02]})---or as \textit{"students"}, \textit{"bachche"}, and \textit{"middle-class families"} (\texttt{[01:45]}, \texttt{[16:20]}). When introducing individual students' stories (\texttt{[00:00]}--\texttt{[02:15]}), his voice drops into an \textbf{intimate, narrative, low-decibel register} (\texttt{$^\circ$word$^\circ$}), his eyebrows draw together in quiet empathy, and his hands rest calmly or open gently toward the camera \texttt{\{\{gesture: gentle open-palm extension\}\}}. When referring to the \textit{"NTA"} or \textit{"coaching nexus"}, he does not shout; instead, his tone becomes \textbf{clinical, measured, and dryly analytical}, using precision thumb-and-index-finger pinch gestures \texttt{\{\{gesture: thumb-index precision pinch\}\}} to dissect institutional failure.\par\smallskip
\end{itemize}

\noindent\textbf{2.} \textbf{Delivery of students' stated demands: When the protesters' core demands (Education Minister's resignation, compensation, release of Sonam Wangchuk, exam reform) are mentioned, are they delivered slowly with deliberate pauses and finger-counting, rushed through in a flat monotone, mocked with air-quotes/scoffs, or trailed off (...)?}\par\smallskip
\begin{itemize}
  \item \textbf{Notes:}
Mohak delivers the students' demands and his own structural reform proposals (\texttt{[16:40]}--\texttt{[19:30]})---(1) independent external audit of NTA's grace-mark algorithm, (2) investigation of the Bihar/Godhra paper leak centres, and (3) decentralising/reforming India's single-point high-stakes exam architecture---\textbf{slowly, with deliberate \texttt{(0.8)} pauses and structured finger-counting} \texttt{\{\{gesture: counting 1, 2, 3 on fingers\}\}}. Each point concludes with a calm, definitive falling pitch \texttt{$\downarrow$} and an earnest nod, treating student demands not as street slogans but as serious public-policy imperatives.\par\smallskip
\end{itemize}

\noindent\textbf{3.} \textbf{Where does vocal \& kinetic energy peak---grievance, street disruption, or electoral horse-race?: Track where the speaker's volume, pitch range, pace, and gesture intensity reach their climax: during accounts of the paper leak/student suicides, during street clashes and barricade breaches, or during partisan BJP-vs-Opposition sparring.}\par\smallskip
\begin{itemize}
  \item \textbf{Notes:}
Because Mohak's baseline is exceptionally composed (\textasciitilde{}125--135 WPM, seminar-style delivery), his energy peaks through \textbf{controlled prosodic contrast and emotional gravity} rather than high-decibel shouting:\par\smallskip
  \item \textbf{Peak 1 (Statistical Proof of Rank Inflation, \texttt{[04:10]}--\texttt{[07:45]}):} His pace quickens and his pitch rises markedly \texttt{$\uparrow$} when demonstrating on a graph how a score of \texttt{620} dropped from rank \texttt{\textasciitilde{}15,000} in 2023 to rank \texttt{\textasciitilde{}60,000} in 2024 (\texttt{"yani \_chaar\_ guna rank gir gayi$\uparrow$!"}, \texttt{[05:12]}), accompanied by a sweeping downward hand trajectory \texttt{\{\{gesture: right hand sweeping sharply downward along chart\}\}}.
  \item \textbf{Peak 2 (Human Cost \& Trust Erosion, \texttt{[17:15]}--\texttt{[19:10]}):} The emotional climax occurs when he discusses the psychological collapse of \texttt{24 lakh} young Indians losing faith in meritocracy (\texttt{"jab mehnat ka fal nahin milta..."}, \texttt{[18:05]}), where he slows down to \textasciitilde{}105 WPM, pauses \texttt{(1.5)} seconds, and looks silently into the lens.
\end{itemize}

\noindent\textbf{4.} \textbf{Non-verbal attribution of responsibility (sighs, stops, head shakes, pointing): When examination authorities (NTA), the Education Minister, police, opposition politicians, or students are named as responsible, what non-verbal cues accompany the attribution (exasperated sighs, mocking laughs, abrupt self-stops ---, or direct finger-pointing at the lens)?}\par\smallskip
\begin{itemize}
  \item \textbf{Notes:}
Responsibility is attributed primarily to the \textbf{National Testing Agency's (NTA) structural opacity and incompetence} through understated, devastating paralinguistic cues:\par\smallskip
  \item When playing clip excerpts of NTA Director General Subodh Kumar Singh claiming at the press conference that \textit{"everything is transparent, integrity is not compromised"} (\texttt{[08:50]}), Mohak cuts back to himself, pauses \texttt{(1.2)}, lets out a quiet, weary \texttt{[sigh]}, and gives a slow head shake \texttt{\{\{gesture: slow bilateral head shake + pursed lips\}\}} before saying: \textit{"Yeh press conference dekh kar (.) answers kam mile (.) aur sawaal \_zyada\_ khade ho gaye$\downarrow$"} (\texttt{[09:12]}).
  \item When noting that the head of the NTA's "grievance committee" is the UPSC Chairman/same bureaucratic circle (\texttt{[09:40]}), Mohak raises both eyebrows and gives a brief, dry \texttt{[chuckle]} of disbelief.
\end{itemize}

\noindent\textbf{5.} \textbf{Prosody of questions---open inquiry vs. loaded rhetorical challenge: Are questions about the protest delivered with level, open intonation, or with sharp rising pitch $\uparrow$, sarcastic stress, raised eyebrows, and trailing clauses (...) that presuppose protesters are violent or politically manipulated?}\par\smallskip
\begin{itemize}
  \item \textbf{Notes:}
Mohak uses \textbf{Socratic, analytical inquiry} to guide the viewer step-by-step. When asking whether the anomaly was caused by an easier paper, grace marks, or a paper leak (\texttt{[03:30]}), he uses level, open-minded inquiry intonation with upturned palms: \textit{"Chaliye (.) teeno possibilities ko data ke saath check karte hain$\downarrow$"} (\texttt{[03:38]}). However, once the data eliminates NTA's excuses, his questions toward the NTA take on a \textbf{pointed, incredulous rising contour \texttt{$\uparrow$}}: \textit{"Agar grace marks sirf \texttt{1,563} students ko mile (.) toh baaki har score bracket mein \texttt{30,000} se \texttt{40,000} rank ka inflation kaise$\uparrow$ aaya?!"} (\texttt{[07:18]}).\par\smallskip
\end{itemize}

\noindent\textbf{6.} \textbf{Visual \& paralinguistic performance of injury and violence: How are claims of violence substantiated non-verbally on screen? (Compare camera close-ups and pained vocal tone on police bandages/fractured hands vs. visual/acoustic treatment of student head injuries, lathi-charge footage, or defensive stammers/pauses when an officer or guest is asked who used force first.)}\par\smallskip
\begin{itemize}
  \item \textbf{Notes:}
Mohak frames injury in \textbf{psychological, familial, and socio-economic terms} rather than street violence. He shows close-up B-roll of exhausted students sitting outside coaching centres in Kota and Delhi, parents mortgaging land for coaching fees, and peaceful candlelight/placard demonstrations (\texttt{[01:15]}, \texttt{[15:50]}). His vocal tone during these visuals is hushed and compassionate (\texttt{$^\circ$ek middle-class parivar apni poori jama-poonji laga deta hai$^\circ$}, \texttt{[01:32]}), accompanied by a solemn head nod.\par\smallskip
\end{itemize}

\noindent\textbf{7.} \textbf{Asymmetric delivery of police actions vs. protester actions: Compare the vocal cadence, pauses, and hand gestures used when describing police actions (e.g., calm, matter-of-fact administrative tone on "tear gas chhode gae" / "preventive arrangement") versus protester actions (alarmist pitch, striking gestures, or ironic softening as "retaliation").}\par\smallskip
\begin{itemize}
  \item \textbf{Notes:}
Mohak contrasts the \textbf{earnest, evidence-backed mobilisation of students and teachers} with the \textbf{evasive bureaucracy of the NTA}:\par\smallskip
  \item \textbf{Students and educators} filing RTIs, compiling rank-vs-marks spreadsheets, and protesting peacefully (\texttt{[06:15]}, \texttt{[15:40]}) are described with warm, admiring prosody and open-handed gestures.
  \item \textbf{Bihar Police's Patna FIR (\texttt{5 May})} detailing the arrested solver gang (\texttt{[11:10]}--\texttt{[12:30]}) is presented with crisp, factual authority, pointing at the timeline graphic on screen.
  \item \textbf{NTA's administrative actions} (changing answers in the provisional key, awarding arbitrary grace marks without prior notification in the information bulletin, \texttt{[08:10]}) are delivered with a furrowed brow and slow, incredulous head shakes.
\end{itemize}

\noindent\textbf{8.} \textbf{Acoustic \& visual space granted to students' own voices: Are protesting students heard in their own voice with clear audio, steady close-ups, and empathetic listening nods, or are they reduced to noisy background chaos B-roll while studio anchors/politicians speak over them?}\par\smallskip
\begin{itemize}
  \item \textbf{Notes:}
Mohak grants \textbf{maximum narrative primacy to students' voices}, opening the entire documentary (\texttt{[00:00]}--\texttt{[01:40]}) with direct video testimonies of NEET aspirants before he even introduces his own thesis. The student clips are played cleanly at full volume with subtitles. When transitioning from a crying student's byte to his studio camera (\texttt{[01:25]}), Mohak holds a \texttt{(1.5)} second silent pause, his face unsmiling and eyes slightly narrowed in empathy \texttt{\{\{gesture: 1.5s silent gaze + soft nod\}\}}, ensuring the student's testimony sets the emotional anchor for the entire 20-minute analysis.\par\smallskip
\end{itemize}

\noindent\textbf{9.} \textbf{Non-verbal shift between the NEET paper-leak issue and the street protest: Does the speaker's non-verbal demeanour shift from sympathetic concern (sombre tone, sighs, solidarity gestures) when discussing the exam leak itself to combative, authoritarian, or dismissive delivery once students march toward Parliament?}\par\smallskip
\begin{itemize}
  \item \textbf{Notes:}
\textbf{No hostile shift occurs.} Instead, Mohak's non-verbal delivery transitions from \textbf{intimate personal empathy} (\texttt{[00:00]}--\texttt{[02:30]}), to \textbf{rigorous data-driven investigation} (\texttt{[03:00]}--\texttt{[14:30]}), to \textbf{solemn civic warning} (\texttt{[15:30]}--\texttt{[20:12]}). When discussing students protesting on the streets and approaching the Supreme Court (\texttt{[16:10]}), his tone remains supportive and dignified, framing their protest as a desperate attempt to salvage institutional trust (\texttt{"ye sirf ek exam ki baat nahin hai (.) ye desh ke youth ke \_bharose\_ ki baat hai$\downarrow$"}, \texttt{[18:45]}).\par\smallskip
\end{itemize}

\subsubsection*{Apply the five non-verbal features (Event 2)}

\begin{longtable}{|p{0.31\textwidth}|p{0.31\textwidth}|p{0.31\textwidth}|}
\hline
Non-verbal feature & Timestamp + spoken anchor words + non-verbal notation & What it does in context (stance, framing, implication) \\ \hline \endhead
\textbf{1. Pauses, silences \& stops} & \texttt{[01:25]} \texttt{[(1.5) silent pause after student testimony]} / \texttt{[09:08]} \textit{"NTA Director ne kaha--- (1.0) 'integrity is not compromised'$\downarrow$"} & The \texttt{(1.5)} opening silence honours student grief; the \texttt{(1.0)} pause before quoting the NTA Director isolates the official claim so its contradiction with the data is glaring. \\ \hline
\textbf{2. Trailing sentences (\texttt{...})} & \texttt{[18:05]} \textit{"Jab ek \texttt{18} saal ka bachcha dekhta hai ki \texttt{2} saal ki mehnat ke baad bhi system aisa hai... toh uska bharosa..."} \texttt{\{\{gesture: slow head shake + palms dropping to lap\}\}} & Trails off on the erosion of youth trust in India's meritocracy, leaving the long-term societal consequence hanging in the air with quiet gravity. \\ \hline
\textbf{3. Intonation, pitch \& volume} & \texttt{[05:12]} \textit{"\texttt{620} marks par pichhle saal rank thi \texttt{15,000}---is saal (.) \_SAATH HAZAAR\_ (\texttt{60,000})$\uparrow$!"} / \texttt{[01:32]} \textit{"$^\circ$ek middle-class parivar apni poori jama-poonji laga deta hai$^\circ$"} & Uses a restrained baseline so that the rare pitch peak \texttt{$\uparrow$} on \texttt{SAATH HAZAAR} and the hushed drop \texttt{$^\circ$...$^\circ$} on family sacrifice carry heightened emotional and analytical contrast. \\ \hline
\textbf{4. Coughs, sighs \& vocalisations} & \texttt{[09:10]} \texttt{[[sigh]]} \textit{"Yeh press conference dekh kar answers kam mile, sawaal zyada..."} / \texttt{[09:42]} \texttt{[[chuckle]]} (on NTA appointing its own oversight panel) & The weary \texttt{[sigh]} performs disappointed civic expectation rather than partisan rage; the brief dry \texttt{[chuckle]} highlights bureaucratic absurdity. \\ \hline
\textbf{5. Hand gestures, face \& posture} & \texttt{[03:38]} \texttt{\{\{gesture: counting 3 hypotheses on fingers\}\}} / \texttt{[05:15]} \texttt{\{\{gesture: hand sweeping downward along rank-inflation graph\}\}} / \texttt{[09:10]} \texttt{\{\{gesture: slow head shake + pursed lips\}\}} & Finger-counting organises the inquiry like a research seminar; graph-tracing grounds claims in visual statistics; the slow head shake signals measured institutional critique. \\ \hline
\end{longtable}

\subsubsection*{Short summary (Event 2 Non-Verbal Synthesis)}

\begin{itemize}
  \item \textbf{What overall interpretation or emotional stance does the non-verbal delivery encourage?}
The non-verbal delivery encourages \textbf{thoughtful, data-literate empathy and deep institutional concern}. Rather than whipping up frantic anger, Mohak's calm cadence, long empathetic silences \texttt{(1.5)}, and methodical chart gestures invite the viewer to understand \textit{how} the exam system broke and \textit{why} students have every right to demand structural reform.\par\smallskip
  \item \textbf{Whose voice or physical presence is acoustically/visually foregrounded, undermined, or silenced?}
\textbf{Individual NEET aspirants (Ayushi, Anjali) and empirical score distributions} are foregrounded right from \texttt{[00:00]}, given clean audio and respectful pauses. \textbf{NTA Director General Subodh Kumar Singh's press conference assurances} are quietly undermined---not by shouting over him, but by playing his exact clip and following it with a \texttt{(1.2)} weary \texttt{[sigh]} and a devastating statistical rebuttal (\texttt{[07:18]}, \texttt{[09:10]}).\par\smallskip
  \item \textbf{What evidence supports your non-verbal reading? What complicates it (e.g., speaker baseline, format)?}
The deliberate pacing shift from \textasciitilde{}130 WPM during data analysis (\texttt{[04:00]}--\texttt{[12:00]}) down to \textasciitilde{}105 WPM with softened volume \texttt{$^\circ$...$^\circ$} and \texttt{(1.5)} silences during student impact segments (\texttt{[01:15]}, \texttt{[18:05]}) demonstrates how prosody frames the human stakes. What complicates a purely event-specific reading is \textbf{speaker baseline}: Mohak Mangal's low-decibel, seminar-style delivery and open-palm gestures are consistent across all \textit{Soch} videos, representing a deliberate counter-brand to high-decibel TV news.\par\smallskip
  \item \textbf{How do the non-verbal cues interact with the spoken words? (Do they reinforce, dramatise, soften/hedge, invert meaning via sarcasm/irony, or leave an accusation unspoken via trailing off?)}
His non-verbal cues \textbf{reinforce and visualise} complex statistical arguments (hand trajectories matching rank-inflation curves), \textbf{soften and humanise} policy analysis through quiet vocal drops \texttt{$^\circ$...$^\circ$}, and \textbf{gently invert bureaucratic jargon} via a weary \texttt{[sigh]} and dry \texttt{[chuckle]} when the NTA claims "complete transparency" (\texttt{[09:10]}).\par\smallskip
  \item \textbf{What context, speaker baseline, or format factors support or complicate this interpretation? (e.g., Is this the speaker's normal style across all videos, or a marked shift on this specific actor/claim?)}
Mohak's decision to skip Event 1 (Ram Mandir land controversy) while devoting a 20-minute data documentary to Event 2 (NEET crisis) reflects his channel's sociological baseline. Within Event 2, his restraint makes his rare non-verbal evaluations---such as the slow head shake and \texttt{[sigh]} after the NTA press conference (\texttt{[09:10]})---far more persuasive to undecided viewers than continuous shouting.\par\smallskip
  \item \textbf{How does another outlet perform the same moment, actor, or claim non-verbally?}
Compare Mohak Mangal (\texttt{[09:10]}) with \textbf{Republic TV} (\texttt{Event 2}) and \textbf{Dhruv Rathee} (\texttt{Event 2}): when addressing NTA failures, Arnab Goswami shouts at maximum volume (\texttt{WORD}) and thumps his desk, while Dhruv Rathee points an accusatory index finger at the lens with sharp prosecutorial stress. Mohak Mangal achieves an equally damning indictment of the NTA using the exact opposite non-verbal strategy: a \texttt{(1.2)} silence, a quiet \texttt{[sigh]}, and a downward-sweeping hand gesture across a scatter plot of rank inflation.\par\smallskip
  \item \textbf{My interpretation, with reasons:}
Mohak Mangal exemplifies \textbf{low-arousal, high-credibility multimodal persuasion}. By refusing to use alarmist prosody or aggressive finger-pointing, and by anchoring his video in unmuted student testimonies (\texttt{[00:00]}--\texttt{[01:40]}) and visual data graphs (\texttt{[04:10]}--\texttt{[07:45]}), his non-verbal performance depoliticises the student protest---framing the NEET controversy not as a BJP-vs-Opposition fight, but as a devastating betrayal of Indian middle-class meritocracy.\par\smallskip
\end{itemize}
\newpage
\section*{Non-Verbal \& Multimodal Discourse Analysis Worksheet}

\textit{Qualitative analysis of pauses, trailing sentences, prosody/intonation, vocalisations (coughs, sighs, scoffs), and kinesics (hand gestures, facial expressions, posture) in political news discourse}\par\smallskip

\bigskip\hrule\bigskip

\subsection*{Guidelines: Five non-verbal \& paralinguistic features to look for in every event}

\begin{longtable}{|p{0.23\textwidth}|p{0.23\textwidth}|p{0.23\textwidth}|p{0.23\textwidth}|}
\hline
Feature & What to look for & Illustrative example (from sample) & Questions to answer \\ \hline \endhead
\textbf{1. Pauses, Silences \& Abrupt Stops} & Micro-pauses \texttt{(.)}, timed silences \texttt{(1.5)}, pregnant/dramatic pauses before a loaded word, hesitant stumbles, abrupt self-interruptions or cut-offs (\texttt{---}), and forced silences (muting a panelist's mic or talking over them). & Anchor pauses \texttt{(1.5)} before saying \textit{"'loota'"} for dramatic weight; Addl. DCP cuts off mid-sentence (\textit{"wo toh--- abhi investigation chal raha hai"}) when asked if police struck first; panelist's mic faded during counter-question. & Where does the pause or stop occur? Is it deliberate/rhetorical (for suspense, emphasis, or gravitas) or defensive/evasive (cognitive load, self-repair, avoiding blame)? Who controls turn-taking and silence? \\ \hline
\textbf{2. Trailing Sentences \& Unfinished Utterances} & Clauses that trail off (\texttt{...}), drop in volume at the end, or leave the conclusion unspoken (aposiopesis) so the audience supplies the implication without the speaker taking explicit responsibility. & \textit{"Aur phir wahan jo log ghuse the... aap samajh hi sakte hain..."}; \textit{"500 salon ke sangharsh ke baad agar aisa..."} trailing off with a slow head shake. & What exact thought is left unfinished? What presupposition or shared bias is the viewer invited to fill in? Does trailing off act as a safe hedge to imply higher-level complicity or "anti-national" motive without stating it? \\ \hline
\textbf{3. Intonation, Pitch, Volume \& Prosody} & Pitch contour (rising \texttt{$\uparrow$} for incredulity/challenge; falling \texttt{$\downarrow$} for finality), volume shifts (shouting vs. hushed whisper), emphatic word stress, speech rate (rapid-fire urgency vs. slow solemnity), and tonal register (sarcastic, reverential, combative). & Hushed, reverential cadence on \textit{"Prabhu Shri Ram ke Mandir"} vs. sharp, high-pitched interrogation on \textit{"par Akhilesh ji Ram Mandir kyon$\uparrow$ nahin jate?!"}; satirical sing-song inflection on \textit{"priventiv arenjment"}. & Which exact words receive pitch peaks or volume stress? How does vocal tone turn neutral words into irony, mockery, or moral outrage? Does the speaker shift prosodic register depending on which actor is discussed? \\ \hline
\textbf{4. Coughs, Sighs \& Non-Lexical Vocalisations} & Audible sighs (\texttt{[sigh]}), scoffs/snorts (\texttt{[scoff]}), ironic chuckles (\texttt{[chuckle]}), throat-clearing/coughs (\texttt{[cough]}), sharp intakes of breath/gasps (\texttt{[gasp]}), and tongue-clicks (\texttt{[tsk]}). & Heavy \texttt{[sigh]} before reporting Trust oversight lapses (signalling reluctant sorrow rather than hostility); mocking \texttt{[scoff]} or \texttt{[chuckle]} when reading an opponent's denial; \texttt{[cough / throat-clear]} when pivoting away from an awkward fact. & What emotion or epistemic stance does the vocalisation perform (exasperation, moral disgust, disbelief, ridicule, discomfort)? Does a sigh or scoff pre-judge a speaker's claim before any verbal rebuttal is made? \\ \hline
\textbf{5. Hand Gestures, Facial Expressions \& Posture (Kinesics)} & Hand movements (air-quotes, finger-pointing/wagging, desk-thumping, open palms, counting on fingers, dismissive wave), facial cues (raised eyebrows, eye-rolls, smirk, furrowed brow), head nods/shakes, leaning into camera, and visual props (bandaged limbs, documents). & Anchor forms physical air-quotes on \textit{"'loota'"} or \textit{"'peaceful' protest"}; host points finger/pen at lens during an ultimatum; injured DCP holds bandaged hand up to camera; creator raises one eyebrow and leans in for a conspiratorial reveal. & How do gestures reinforce, contradict, or sharpen the spoken words? During split-screen listening, does the host's face (nodding vs. smirking/eye-rolling) cue the audience on whom to trust? How are bodily visuals used as evidence? \\ \hline
\end{longtable}

\subsection*{Quick Multimodal Transcription Notation Key (use inside transcript excerpts)}

\begin{longtable}{|p{0.15\textwidth}|p{0.15\textwidth}|p{0.15\textwidth}|p{0.15\textwidth}|p{0.15\textwidth}|p{0.15\textwidth}|}
\hline
Symbol & Meaning & Example in transcript & Symbol & Meaning & Example in transcript \\ \hline \endhead
\texttt{(.)} & Brief micro-pause (< 0.5 sec) & \texttt{woh (.) suparavijan lacking tha} & \texttt{$\uparrow$ / $\downarrow$} & Marked rising / falling pitch & \texttt{kyon$\uparrow$ FIR nahin karaee$\downarrow$} \\ \hline
\texttt{(1.5)} & Timed pause in seconds & \texttt{kisne (1.5) 'loota'?} & \texttt{WORD / \_word\_} & Loud volume / emphatic stress & \texttt{ye PAAP hai / \_bilkul\_ galat} \\ \hline
\texttt{---} & Abrupt stop, cut-off, or self-repair & \texttt{police ki taraf se--- abhi jaanch...} & \texttt{$^\circ$word$^\circ$} & Softened / whispered delivery & \texttt{$^\circ$ander ki baat ye hai$^\circ$} \\ \hline
\texttt{...} & Trailing off / unfinished sentence & \texttt{aap jaante hain ki kaun log...} & \texttt{[sigh] [cough] [scoff]} & Non-lexical vocalisations & \texttt{[sigh] 500 salon ke baad...} \\ \hline
\texttt{[overlap]} & Simultaneous / overlapping speech & \texttt{[Anjana: suniye---] [Guest: nahin!]} & \texttt{\{\{gesture: ...\}\}} & Hand gesture / facial cue / posture & \texttt{\{\{air-quotes on 'loota'\}\}} \\ \hline
\end{longtable}

\textbf{Methodological Reminder:} A non-verbal cue is not automatically manipulative or ideological---a cough, pause, or rapid cadence may reflect studio audio delay, live time pressure, or a speaker's habitual style. Always check (1) whether the cue clusters with verbal features, and (2) whether the speaker's delivery shifts noticeably across different actors or events (within-outlet baseline).\par\smallskip

\bigskip\hrule\bigskip

\subsection*{Event 1. Ram Mandir donation controversy (Non-Verbal \& Paralinguistic Layer)}

\begin{itemize}
  \item \textbf{Focus:} How do prosody, pauses, trailing sentences, vocalisations (sighs, coughs, scoffs), and body language frame the controversy---as solemn devotional sorrow, administrative scrutiny, or partisan combat?
  \item \textbf{Outlet:} Aaj Tak (\textit{Black \& White} / Special Bulletin with Anjana Om Kashyap) | \textbf{Date:} 23 July 2024
  \item \textbf{Video title:} Ram Mandir Donation Controversy / Counting Theft \& Audit Scrutiny Bulletin | \href{https://youtu.be/B4L1pz0o4dc}{https://youtu.be/B4L1pz0o4dc} | \textbf{Timestamp(s):} \texttt{[00:00]}--\texttt{[08:12]}
  \item \textbf{Analyst:} Baani Solanki (\texttt{2025114016}) | \textbf{Language:} Hindi (transcribed in Roman script)
  \item \textbf{Format / Visual setup:} \texttt{[ ]} TV Panel Debate \texttt{[x]} TV Bulletin \texttt{[ ]} Solo YouTube Explainer \texttt{[ ]} Interview / Discussion
  \item \textbf{Primary speaker(s) observed:} Anjana Om Kashyap (standing/seated studio news anchor delivering a dramatic prime-time monologue with CCTV footage inserts, temple visuals, and opposition soundbites)
\end{itemize}

\subsubsection*{Event-specific non-verbal questions (Event 1)}

\noindent\textbf{1.} \textbf{Vocal \& gestural marking of core labels ("scam", "loota", "paap", "gadabdi", "avyavastha"): When the controversy is named, what non-verbal cues accompany the label? (Look for physical air-quotes, vocal scare-quotes, dramatic pauses before the word, hushed/reverential drops in volume, or shouted stress.)}\par\smallskip
\begin{itemize}
  \item \textbf{Notes:}
Anjana Om Kashyap's vocal and gestural marking of core labels is the clearest example in our dataset of \textbf{sacralising individual theft while sanitising institutional failure}:\par\smallskip
  \item \textbf{Moral/Religious Labels (\texttt{"paap"}, \texttt{"aparadh"}):} When describing the counting clerk stealing cash from the offering boxes (\texttt{[00:42]}), she pauses \texttt{(1.2)}, drops her voice to a hushed, trembling register \texttt{$^\circ$...$^\circ$}, and then stresses \texttt{"PAAP"} with a deep, mournful falling pitch \texttt{$\downarrow$} and a slow head shake: \textit{"Prabhu Shri Ram ke charnon mein chadhaye gaye daan ki chori (1.2) ye sirf aparadh nahin (.) \_PAAP\_ hai$\downarrow$"} (\texttt{[00:48]}).
  \item \textbf{Political Accusation Labels (\texttt{"'loota'"}, \texttt{"'ghotala'"}):} When quoting Samajwadi Party leaders (Akhilesh Yadav / Awadhesh Prasad) who alleged that the Temple Trust \textit{"looted"} donations (\texttt{[04:15]}), Anjana pauses \texttt{(1.5)} before the word \textit{"'loota'"} and forms \textbf{physical air-quotes with both hands} \texttt{\{\{gesture: bilateral two-finger air-quotes on 'loota'\}\}} accompanied by raised eyebrows and a sharp, incredulous rising tone \texttt{$\uparrow$}, visually quarantine-marking the word \textit{"loota"} as a partisan slander.
  \item \textbf{Institutional Lapses (\texttt{"avyavastha"}, \texttt{"supervision lacking"}):} When describing the Trust's own lack of frisking and audit protocols (\texttt{[02:35]}), she switches to a brisk, low-arousal administrative cadence with a brief micro-pause \texttt{(.)}: \textit{"Wahan par (.) thoda supervision lacking tha, thodi \_avyavastha\_ thi$\downarrow$"} (\texttt{[02:42]}), accompanied by a small, minimising hand wave \texttt{\{\{gesture: small horizontal hand wave\}\}}.
\end{itemize}

\noindent\textbf{2.} \textbf{Differential demeanour toward actors (Temple, Trust leadership, counting clerks, donors, opposition accusers): How do the speaker's voice pitch, facial expression, and posture shift when referring to Prabhu Shri Ram / devotees vs. top figures (Champat Rai, Nripendra Misra, CM Yogi) vs. arrested counting staff vs. opposition leaders (Akhilesh Yadav, Awadhesh Prasad)?}\par\smallskip
\begin{itemize}
  \item \textbf{Notes:}
Anjana exhibits four sharply differentiated non-verbal personas depending on the actor named:\par\smallskip
  \item \textbf{Prabhu Shri Ram \& Devotees (\texttt{[00:15]}--\texttt{[01:10]}):} Hushed, reverential cadence (\texttt{$^\circ$Prabhu Shri Ram ke Mandir$^\circ$}), slightly bowed head, softened eyes, and hands brought together near her chest/waist in a semi-devotional posture \texttt{\{\{gesture: hands clasped gently at waist + bowed head\}\}}.
  \item \textbf{Arrested Counting Clerk / Bank Staff (\texttt{[01:20]}--\texttt{[02:20]}):} Stern, moralising scowl, pointing a rigid index finger toward the CCTV screen showing the clerk stuffing notes into his pocket (\texttt{[01:45]}), with clipped, disapproving downward stress.
  \item \textbf{Top Figures (Nripendra Misra, Champat Rai, CM Yogi, \texttt{[03:10]}--\texttt{[03:55]}):} Deferential, matter-of-fact institutional tone; upright posture and respectful nods when noting that Nripendra Misra ordered a digital counting overhaul.
  \item \textbf{Opposition Leaders (Akhilesh Yadav, Awadhesh Prasad, \texttt{[04:10]}--\texttt{[06:45]}):} Combative, high-pitched prosecutorial interrogation \texttt{$\uparrow$}, chin raised, right index finger wagging at the camera lens \texttt{\{\{gesture: assertive index-finger wag at lens\}\}} (\textit{"Par Akhilesh ji Ram Mandir kyon$\uparrow$ nahin jaate?!"}, \texttt{[05:18]}).
\end{itemize}

\noindent\textbf{3.} \textbf{Solemn devotional grief vs. high-decibel partisan spectacle: How do pauses, sighs, speech rate, and hand gestures construct the segment's emotional register? Where is delivery slow, hushed, and mournful (sighs, bowed head, folded hands), and where does it escalate into rapid-fire shouting, desk-thumping, or mocking irony?}\par\smallskip
\begin{itemize}
  \item \textbf{Notes:}
The bulletin is structured as a dramatic two-act non-verbal performance:\par\smallskip
  \item \textbf{Act I --- Solemn Devotional Grief (\texttt{[00:00]}--\texttt{[03:50]}):} Speech rate is slow (\textasciitilde{}115 WPM), punctuated by heavy audible \texttt{[sigh]} vocalisations (\texttt{[00:35]}, \texttt{[02:30]}), long \texttt{(1.2)}--\texttt{(1.5)} pauses, mournful head shakes, and hushed volume \texttt{$^\circ$...$^\circ$} lamenting that after \textit{"500 salon ke sangharsh ke baad..."} (\texttt{[00:38]}) a petty clerk betrayed devotees' faith in the sanctum.
  \item \textbf{Act II --- High-Decibel Partisan Counter-Attack (\texttt{[04:00]}--\texttt{[07:50]}):} The moment the Samajwadi Party's allegations are introduced, Anjana's speech rate jumps to \textasciitilde{}155 WPM, her pitch rises sharply \texttt{$\uparrow$}, her volume surges, and her hand gestures shift from mournful clasped hands to aggressive finger-pointing and palm-chopping against opposition "hypocrisy" (\texttt{"Jo log karsevakon par goli chalane walon ke saath the..."}, \texttt{[05:40]}).
\end{itemize}

\noindent\textbf{4.} \textbf{Non-verbal evaluation of financial scrutiny: When questions about unreceipted offerings, lack of frisking, or land deals are raised, how does the host/speaker react non-verbally (earnest nodding, leaning into the camera, open-palm appeals vs. eye-rolls, scoffs, dismissive hand waves, or incredulous rising pitch $\uparrow$)?}\par\smallskip
\begin{itemize}
  \item \textbf{Notes:}
Anjana splits financial scrutiny into two categories:\par\smallskip
  \item \textbf{Clerical Oversight (Lack of Frisking / Manual Counting):} Treated with \textbf{reluctant, maternal sorrow}---she lets out a heavy \texttt{[sigh]} at \texttt{[02:30]}, furrows her brow, and opens her palms (\texttt{"sawal toh banta hai ki wahan frisking kyon nahin hui$\downarrow$"}, \texttt{[02:48]}), framing the lapse as innocent naivety (\textit{"bholapan"} / \textit{"avyavastha"}) by a trusting religious body.
  \item \textbf{Systemic Corruption Allegations by Opposition:} Met with \textbf{incredulous rising pitch \texttt{$\uparrow$}, mocking \texttt{[scoff]} vocalisations, and dismissive hand waves} (\texttt{[04:25]}, \texttt{[05:10]}), treating broader audit demands by SP/Congress leaders as malicious attempts to defame Hindu faith.
\end{itemize}

\noindent\textbf{5.} \textbf{Performing certainty vs. hedging and innuendo (pauses, stops, trailing off): Are allegations delivered with definitive falling intonation $\downarrow$, finger-counting, and pointing at on-screen evidence (CCTV, documents), or hedged with hesitant pauses (.), abrupt self-interruptions (---), throat-clearing [cough], and trailing sentences (...) that imply higher-level complicity without naming it?}\par\smallskip
\begin{itemize}
  \item \textbf{Notes:}
  \item \textbf{Certainty on Individual Clerk's Guilt:} Delivered with definitive falling intonation \texttt{$\downarrow$} and direct pointing at the red-circled CCTV footage (\texttt{"CCTV mein range haathon pakda gaya$\downarrow$"}, \texttt{[01:52]}).
  \item \textbf{Hedging on Trust Leadership's Responsibility:} Delivered with micro-pauses \texttt{(.)}, throat-clearing \texttt{[cough]}, and euphemistic self-repair \texttt{---}: \textit{"Trust ki taraf se--- matlab jo counting agency aur bank ke log the (.) unki monitoring mein thodi kami rahi..."} (\texttt{[02:38]}), deflecting blame away from Champat Rai or Nripendra Misra onto external bank staff.
  \item \textbf{Innuendo against Opposition via Trailing Off (\texttt{...}):} At \texttt{[05:55]}, she trails off with a narrowed gaze: \textit{"Aur aaj wahi log Ram Mandir ke chande par sawaal utha rahe hain jinhone hamesha Mandir ka... aap jaante hain ki inka asli maksad kya hai..."} (\texttt{[06:02]}), leaving the accusation of "anti-Hindu conspiracy" unspoken for the viewer to fill in.
\end{itemize}

\noindent\textbf{6.} \textbf{Turn-taking, interruptions, and acoustic treatment of denials: In debates/interviews: who receives calm, uninterrupted speaking time and sympathetic nods, and who is cut off mid-sentence (---), talked over, or muted? In solo videos: how are clips of official denials framed via reaction pauses, raised eyebrows, or sarcastic chuckles?}\par\smallskip
\begin{itemize}
  \item \textbf{Notes:}
In this studio bulletin, recorded soundbites are framed asymmetrically:\par\smallskip
  \item \textbf{Official Trust / Government explanations} (Nripendra Misra's assurance of automated counting machines, \texttt{[03:25]}) are introduced with a calm, reassuring nod and followed by a supportive summary (\texttt{"ab poori vyavastha high-tech ki ja rahi hai$\downarrow$"}, \texttt{[03:45]}).
  \item \textbf{Opposition soundbites} (Awadhesh Prasad / SP leaders demanding a high-level inquiry, \texttt{[04:30]}) are cut short after 8 seconds, immediately followed by Anjana's mocking \texttt{[scoff]}, raised eyebrows, and a rapid-fire counter-interrogation directed at the screen (\texttt{"Sawaal poochhne se pehle Akhilesh Yadav ye batayein---"}, \texttt{[04:42]}).
\end{itemize}

\noindent\textbf{7.} \textbf{Performing the boundary between religious devotion and institutional critique: How does the presenter use vocal tone, sighs, and body language to signal personal reverence for the Temple (or secular objectivity) while criticising---or defending---the Trust's financial management?}\par\smallskip
\begin{itemize}
  \item \textbf{Notes:}
Anjana explicitly performs \textbf{devotional insiderhood}: she never critiques the Temple management from a secular distance, but rather speaks \textit{as an anguished devotee}. Before mentioning the counting lapse at \texttt{[02:28]}, she lets out an audible \texttt{[sigh]}, lowers her voice \texttt{$^\circ$...$^\circ$}, and places her hand over her heart \texttt{\{\{gesture: right hand pressed to heart + bowed head\}\}}: \textit{"Yeh khabar batate hue bhi (.) mann ko dukh hota hai..."} (\texttt{[02:32]}). This non-verbal performance signals that her mild critique of the counting arrangement comes from love for Lord Ram, insulating the channel from any charge of being anti-Mandir while simultaneously delegitimising outsider/Opposition criticism as hostile.\par\smallskip
\end{itemize}

\noindent\textbf{8.} \textbf{Self-presentation as the "guardian of public trust": How does the speaker use direct eye-gaze, posture, moral sighs, or emphatic gestures (pointing at the viewer, hands on heart/desk) to position themselves as the authentic protector of devotees or public accountability?}\par\smallskip
\begin{itemize}
  \item \textbf{Notes:}
Anjana positions herself as the \textbf{moral guardian of the pious Hindu devotee (\texttt{shraddhalu})} against both petty thieves inside the counting room and political opportunists outside it. She uses intense, unblinking eye-gaze into the teleprompter, places her hand on her chest when invoking \textit{"karodon Ram bhakton ki aastha"} (\texttt{[01:05]}), and points her pen/index finger sharply at the lens (\texttt{[05:18]}, \texttt{[07:10]}) when challenging opposition politicians, visually casting herself as the shield of the Temple's sanctity.\par\smallskip
\end{itemize}

\subsubsection*{Apply the five non-verbal features (Event 1)}

\begin{longtable}{|p{0.31\textwidth}|p{0.31\textwidth}|p{0.31\textwidth}|}
\hline
Non-verbal feature & Timestamp + spoken anchor words + non-verbal notation & What it does in context (stance, framing, implication) \\ \hline \endhead
\textbf{1. Pauses, silences \& stops} & \texttt{[02:38]} \textit{"Trust ki taraf se--- matlab jo bank ke log the (.) unki monitoring mein kami rahi"} / \texttt{[04:15]} \textit{"SP ne kaha ki chanda (1.5) 'loota' gaya$\uparrow$"} & The abrupt self-repair \texttt{---} at \texttt{[02:38]} shifts blame from the Trust to bank staff; the \texttt{(1.5)} pause before \textit{"'loota'"} isolates the opposition's word for ridicule. \\ \hline
\textbf{2. Trailing sentences (\texttt{...})} & \texttt{[00:38]} \textit{"[sigh] \texttt{500} salon ke sangharsh ke baad agar Mandir mein aisa..."} / \texttt{[06:02]} \textit{"Jinhone hamesha Mandir ka... aap samajh hi sakte hain..."} & At \texttt{[00:38]}, trailing off with a head shake evokes unspeakable devotional sorrow; at \texttt{[06:02]}, it insinuates that opposition critics harbour an anti-Ram agenda without stating it outright. \\ \hline
\textbf{3. Intonation, pitch \& volume} & \texttt{[00:48]} \textit{"$^\circ$Prabhu Shri Ram ke charnon mein$^\circ$... ye \_PAAP\_ hai$\downarrow$"} vs. \texttt{[05:18]} \textit{"Par Akhilesh ji Ram Mandir kyon$\uparrow$ nahin jaate?!"} & Hushed reverential cadence \texttt{$^\circ$...$^\circ$} on the Temple contrasts sharply with high-pitched, combative rising interrogation \texttt{$\uparrow$} aimed at Akhilesh Yadav. \\ \hline
\textbf{4. Coughs, sighs \& vocalisations} & \texttt{[02:30]} \texttt{[[sigh]]} \textit{"Mann ko dukh hota hai ki wahan suparavijan lacking tha..."} / \texttt{[04:40]} \texttt{[[scoff]]} (right after SP leader's byte) & The heavy \texttt{[sigh]} frames Trust oversight failures as reluctant family sorrow; the \texttt{[scoff]} instantly invalidates the Opposition's call for an audit. \\ \hline
\textbf{5. Hand gestures, face \& posture} & \texttt{[00:25]} \texttt{\{\{gesture: hands clasped at waist + bowed head\}\}} / \texttt{[04:15]} \texttt{\{\{gesture: bilateral physical air-quotes on 'loota'\}\}} / \texttt{[05:18]} \texttt{\{\{gesture: index-finger wag at lens\}\}} & Clasped hands perform piety; physical air-quotes mark \textit{"loota"} as fake news; finger-wagging turns the bulletin into an indictment of the Opposition. \\ \hline
\end{longtable}

\subsubsection*{Short summary (Event 1 Non-Verbal Synthesis)}

\begin{itemize}
  \item \textbf{What overall interpretation or emotional stance does the non-verbal delivery encourage?}
The non-verbal delivery encourages \textbf{reverential sorrow over an isolated clerical sin (\texttt{"paap"}) paired with fierce partisan indignation against opposition critics}. Viewers are guided to view the theft as the work of one greedy bank clerk, the Trust's lapse as innocent administrative oversight (\textit{"avyavastha"}), and the Opposition's questions as malicious political opportunism.\par\smallskip
  \item \textbf{Whose voice or physical presence is acoustically/visually foregrounded, undermined, or silenced?}
\textbf{Anjana's devotional-prosecutorial persona and the Trust's corrective measures} are foregrounded with calm authority and reverential pauses. \textbf{The arrested clerk} is visually shamed via CCTV loops, while \textbf{Opposition leaders (SP/Congress)} are acoustically and visually undermined through physical air-quotes \texttt{\{\{air-quotes on 'loota'\}\}}, mocking \texttt{[scoff]} vocalisations, and rapid counter-questions.\par\smallskip
  \item \textbf{What evidence supports your non-verbal reading? What complicates it (e.g., speaker baseline, format)?}
The unmistakable acoustic shift between Act I (\textasciitilde{}115 WPM, hushed \texttt{$^\circ$...$^\circ$}, \texttt{[sigh]}, clasped hands on Temple/Trust) and Act II (\textasciitilde{}155 WPM, loud \texttt{WORD}, rising pitch \texttt{$\uparrow$}, finger-wagging on Akhilesh Yadav) proves a strong within-speaker shift across actors. While high theatricality is part of Hindi prime-time TV bulletin conventions, the directional asymmetry of air-quotes and self-repairs (\texttt{Trust ki taraf se--- matlab bank ke log}, \texttt{[02:38]}) clearly protects institutional leadership.\par\smallskip
\end{itemize}

\bigskip\hrule\bigskip

\subsection*{Event 2. Students' protest regarding NEET paper leak (Non-Verbal \& Paralinguistic Layer)}

\begin{itemize}
  \item \textbf{Focus:} How do prosody, pauses, trailing sentences, vocalisations, and visual body language construct the legitimacy of student protesters, police use of force, and state accountability?
  \item \textbf{Outlet:} Aaj Tak (Live Prime-Time Debate / Ground \& Studio Broadcast with Anjana Om Kashyap) | \textbf{Date:} 2024 (NEET Protest \& Parliament March Live Coverage)
  \item \textbf{Video title:} NEET Paper Leak Protest Live: Students' Parliament March, Police Clash \& Political Debate | \href{https://www.youtube.com/live/JXLIs0WtJAM}{https://www.youtube.com/live/JXLIs0WtJAM} | \textbf{Timestamp(s):} \texttt{[00:00]}--\texttt{[32:40]}
  \item \textbf{Analyst:} Baani Solanki (\texttt{2025114016}) | \textbf{Language:} Hindi (transcribed in Roman script)
  \item \textbf{Format / Visual setup:} \texttt{[x]} TV Panel Debate \texttt{[ ]} TV Bulletin \texttt{[ ]} Solo YouTube Explainer \texttt{[x]} Interview / Discussion
  \item \textbf{Primary speaker(s) observed:} Anjana Om Kashyap (studio anchor), Additional DCP / Police Officer (live guest showing bandaged hand), Student / Youth Wing Protesters, and BJP vs. Opposition Spokespersons in multi-window split-screen
\end{itemize}

\subsubsection*{Event-specific non-verbal questions (Event 2)}

\noindent\textbf{1.} \textbf{Prosodic \& gestural marking of participant labels ("students/bachche" vs. "protesters" vs. "upadravi/goons"): How do vocal tone, volume, facial expression, and hand gestures shift across labels? (e.g., softened, empathetic pitch for "students" vs. harsh, staccato stress, scowls, or accusatory finger-pointing on "upadraviyon", "goons", or "anti-social elements".)}\par\smallskip
\begin{itemize}
  \item \textbf{Notes:}
Anjana performs a sharp \textbf{acoustic and gestural bifurcation} between "genuine students" at home and "protesters/rowdies" on the street:\par\smallskip
  \item When referring to \textit{"asli chhatra"} or \textit{"masoom bachche"} studying at home (\texttt{[03:15]}), her voice softens slightly with a concerned nod.
  \item However, the moment she refers to those marching toward Parliament and clashing at barricades (\texttt{[05:40]}, \texttt{[11:20]}, \texttt{[18:50]}), she shifts to the labels \textit{"pradarshankari"}, \textit{"chhatra sangathan ke karyakarta"}, and \textit{"upadravi"} (rowdies), delivering \textit{"UPADRAVIYON"} with \textbf{harsh, staccato stress, a deep scowl, and an accusatory pen-point toward the riot B-roll} \texttt{\{\{gesture: sharp pen-point at screen + furrowed scowl\}\}} (\texttt{"Kya chhatron ke bhes mein \_upadraviyon\_ ne kanoon haath mein liya$\uparrow$?!"}, \texttt{[06:12]}).
\end{itemize}

\noindent\textbf{2.} \textbf{Delivery of students' stated demands: When the protesters' core demands (Education Minister's resignation, compensation, release of Sonam Wangchuk, exam reform) are mentioned, are they delivered slowly with deliberate pauses and finger-counting, rushed through in a flat monotone, mocked with air-quotes/scoffs, or trailed off (...)?}\par\smallskip
\begin{itemize}
  \item \textbf{Notes:}
When protesters' broader political demands (such as the Education Minister's resignation or release of detained activists/Sonam Wangchuk alongside NEET re-examination, \texttt{[14:10]}--\texttt{[15:30]}) are mentioned, Anjana \textbf{rushes through them in a rapid, dismissive list} before stopping on the political demands with \textbf{incredulous rising pitch \texttt{$\uparrow$}, raised eyebrows, and a mocking \texttt{[scoff]}}: \textit{"Ab NEET ke naam par istifa maang rahe hain (.) Sonam Wangchuk ki baat kar rahe hain... matlab mudda exam hai ya \_siyasat\_$\uparrow$?!"} (\texttt{[14:45]}), accompanied by a two-handed dismissive wave \texttt{\{\{gesture: bilateral dismissive hand wave\}\}}.\par\smallskip
\end{itemize}

\noindent\textbf{3.} \textbf{Where does vocal \& kinetic energy peak---grievance, street disruption, or electoral horse-race?: Track where the speaker's volume, pitch range, pace, and gesture intensity reach their climax: during accounts of the paper leak/student suicides, during street clashes and barricade breaches, or during partisan BJP-vs-Opposition sparring.}\par\smallskip
\begin{itemize}
  \item \textbf{Notes:}
Vocal and kinetic energy peaks overwhelmingly during \textbf{street clashes/barricade breaches (\texttt{[05:30]}--\texttt{[12:40]}) and high-decibel BJP-vs-Opposition studio sparring (\texttt{[19:00]}--\texttt{[30:00]})}:\par\smallskip
  \item While initial references to the NEET paper leak are delivered at a moderate pace, Anjana's volume, pitch, and gesture intensity skyrocket when showing barricades being pushed and police officers injured (\texttt{"Dekhiye kaise \_pathar\_ chalaaye gaye! Barricade \_toda\_ gaya!"}, \texttt{[07:15]}), leaning forward over her desk and striking her palm with her pen \texttt{\{\{gesture: striking palm with pen\}\}}.
\end{itemize}

\noindent\textbf{4.} \textbf{Non-verbal attribution of responsibility (sighs, stops, head shakes, pointing): When examination authorities (NTA), the Education Minister, police, opposition politicians, or students are named as responsible, what non-verbal cues accompany the attribution (exasperated sighs, mocking laughs, abrupt self-stops ---, or direct finger-pointing at the lens)?}\par\smallskip
\begin{itemize}
  \item \textbf{Notes:}
  \item \textbf{Students / Opposition Youth Wings:} Named as responsible for street chaos with \textbf{direct finger-pointing at their split-screen windows}, stern head shakes, and loud interruptions (\texttt{"Aapke karyakartaon ne police par haath uthaya$\downarrow$!"}, \texttt{[12:10]}).
  \item \textbf{Police / Administration:} Completely absolved non-verbally via \textbf{sympathetic head nods and pained facial expressions} (\texttt{[08:20]}--\texttt{[10:45]}).
  \item \textbf{NTA / Government:} When an opposition panelist tries to attribute primary blame to the Education Minister (\texttt{[21:15]}), Anjana lets out an exasperated \texttt{[sigh]} and cuts them off mid-sentence \texttt{---} (\texttt{"CBI jaanch toh ho rahi hai na--- phir sadak par hungama kyon$\uparrow$?!"}, \texttt{[21:22]}).
\end{itemize}

\noindent\textbf{5.} \textbf{Prosody of questions---open inquiry vs. loaded rhetorical challenge: Are questions about the protest delivered with level, open intonation, or with sharp rising pitch $\uparrow$, sarcastic stress, raised eyebrows, and trailing clauses (...) that presuppose protesters are violent or politically manipulated?}\par\smallskip
\begin{itemize}
  \item \textbf{Notes:}
Questions directed at student leaders and opposition spokespersons are \textbf{100\% loaded rhetorical challenges} delivered with sharp rising pitch \texttt{$\uparrow$}, sarcastic stress, arched eyebrows, and accusatory trailing clauses (\texttt{...}), presupposing guilt:\par\smallskip
  \item To student/opposition representative (\texttt{[11:45]}): \textit{"Shantipurna pradarshan mein police walon ke haath kaise$\uparrow$ toot gaye?! Kya ye chhatra hain ya phir inke peechhe..."} (\texttt{[11:52]}), trailing off with a pointed stare that presupposes political conspiracy.
  \item Conversely, questions to the Additional DCP (\texttt{[08:30]}) are delivered with a \textbf{softened, solicitous, open intonation}: \textit{"Sir, aapke haath mein chot kaise lagi, zara batayein$\downarrow$"} (\texttt{[08:35]}).
\end{itemize}

\noindent\textbf{6.} \textbf{Visual \& paralinguistic performance of injury and violence: How are claims of violence substantiated non-verbally on screen? (Compare camera close-ups and pained vocal tone on police bandages/fractured hands vs. visual/acoustic treatment of student head injuries, lathi-charge footage, or defensive stammers/pauses when an officer or guest is asked who used force first.)}\par\smallskip
\begin{itemize}
  \item \textbf{Notes:}
This broadcast contains the classic textbook contrast in \textbf{multimodal substantiation of injury}:\par\smallskip
  \item \textbf{Police Injury Foregrounded (\texttt{[08:30]}--\texttt{[10:15]}):} The camera zooms into a tight close-up of the Additional DCP holding his \textbf{white-bandaged hand} up to the lens \texttt{\{\{gesture: injured DCP holds bandaged hand up to camera\}\}} while Anjana lowers her voice into a pained, sympathetic tone (\texttt{$^\circ$dekhiye DCP sahab ke haath mein fracture hai$^\circ$}, \texttt{[08:48]}) and nods gravely.
  \item \textbf{Student Injury \& Police Force Hedged (\texttt{[10:20]}--\texttt{[10:55]}):} When a reporter/question briefly asks whether police lathi-charged students first, the officer stammers and cuts himself off mid-sentence \texttt{---}: \textit{"Dekhiye wo toh--- (.) abhi investigation chal raha hai, humne sirf (.) \_preventive arrangement\_ kiya tha..."} (\texttt{[10:35]}). Instead of pressing him on the stammer \texttt{---}, Anjana nods acceptingly and pivots away \texttt{[cough / throat-clear]} (\texttt{[10:42]}), while visuals of students being dragged by police play only in a small, muted corner window.
\end{itemize}

\noindent\textbf{7.} \textbf{Asymmetric delivery of police actions vs. protester actions: Compare the vocal cadence, pauses, and hand gestures used when describing police actions (e.g., calm, matter-of-fact administrative tone on "tear gas chhode gae" / "preventive arrangement") versus protester actions (alarmist pitch, striking gestures, or ironic softening as "retaliation").}\par\smallskip
\begin{itemize}
  \item \textbf{Notes:}
  \item \textbf{Police Actions:} Described in a \textbf{calm, passive, matter-of-fact administrative cadence} with soothing open-palm gestures: \textit{"Police ki taraf se bas halka bal prayog kiya gaya (.) tear gas chhode gae taaki situation control mein rahe$\downarrow$"} (\texttt{[07:45]}), echoing the officer's euphemism \textit{"preventive arrangement"}.
  \item \textbf{Protester Actions:} Described in an \textbf{alarmist, high-velocity, active cadence} with striking hand gestures: \textit{"Unhone barricade \_tode\_! Police par \_hamla\_ kiya$\uparrow$!"} (\texttt{[07:15]}). When a student panelist claims they only pushed back in self-defence (\textit{"retaliation"}), Anjana repeats \textit{"'retaliation'$\uparrow$?!"} (\texttt{[12:25]}) with a mocking sing-song rising contour and an eye-roll \texttt{\{\{gesture: eye-roll + head toss\}\}}.
\end{itemize}

\noindent\textbf{8.} \textbf{Acoustic \& visual space granted to students' own voices: Are protesting students heard in their own voice with clear audio, steady close-ups, and empathetic listening nods, or are they reduced to noisy background chaos B-roll while studio anchors/politicians speak over them?}\par\smallskip
\begin{itemize}
  \item \textbf{Notes:}
Protesting students are \textbf{acoustically subordinated and frequently talked over \texttt{[overlap]}}. For most of the broadcast, street protesters appear in muted or noisy ambient B-roll loops showing shouting and barricade shoving while studio graphics blare \textit{"NEET KE NAAM PAR HUNGAMA"} (\texttt{[06:00]}). When a student/youth-wing representative speaks in the multi-box panel (\texttt{[11:30]}--\texttt{[13:50]}), Anjana interrupts them within 10--12 seconds \texttt{[overlap]}, raising her voice over theirs (\textit{"[Anjana: Suniye--- suniye! Barricade kisne toda?!] [Student: Ma'am humari baat toh---]"}, \texttt{[12:18]}) and eventually having the producer fade down the student's mic volume so the anchor's reprimand dominates the audio mix.\par\smallskip
\end{itemize}

\noindent\textbf{9.} \textbf{Non-verbal shift between the NEET paper-leak issue and the street protest: Does the speaker's non-verbal demeanour shift from sympathetic concern (sombre tone, sighs, solidarity gestures) when discussing the exam leak itself to combative, authoritarian, or dismissive delivery once students march toward Parliament?}\par\smallskip
\begin{itemize}
  \item \textbf{Notes:}
\textbf{Yes---a dramatic 180-degree non-verbal shift occurs.} In the opening two minutes (\texttt{[01:10]}--\texttt{[03:10]}), when mentioning the NEET paper leak in the abstract, Anjana uses a sombre, concerned tone and a sympathetic \texttt{[sigh]} (*"chhatron ke bhavishya ka sawaal hai$\downarrow$"\texttt{, }[01:45]\texttt{). However, the instant the broadcast transitions to the students marching toward Parliament (}[05:30]\texttt{ onward), her sympathy vanishes and is replaced by an \textbf{authoritarian, law-and-order prosecutorial demeanour}: her posture stiffens, her brow furrows into a scowl, her pitch spikes }$\uparrow$`, and she wags her finger at the protesters, treating collective street action as illegitimate disorder.\par\smallskip
\end{itemize}

\subsubsection*{Apply the five non-verbal features (Event 2)}

\begin{longtable}{|p{0.31\textwidth}|p{0.31\textwidth}|p{0.31\textwidth}|}
\hline
Non-verbal feature & Timestamp + spoken anchor words + non-verbal notation & What it does in context (stance, framing, implication) \\ \hline \endhead
\textbf{1. Pauses, silences \& stops} & \texttt{[10:35]} Addl. DCP: \textit{"Wo toh--- (.) abhi investigation chal raha hai"} / \texttt{[12:18]} \texttt{[overlap]} \textit{"[Anjana: Suniye--- barricade kisne toda?!] [Student: Ma'am lathi---]"} (student mic faded) & The officer's abrupt self-cut-off \texttt{---} betrays evasion when asked if police struck first; Anjana's \texttt{[overlap]} interruption and mic-fading silence the student's counter-testimony. \\ \hline
\textbf{2. Trailing sentences (\texttt{...})} & \texttt{[11:52]} \textit{"Kya ye waqai NEET ke chhatra hain (.) ya phir inke peechhe kisi siyasi party ka..."} \texttt{\{\{gesture: arched eyebrow + pointed pen\}\}} & Leaves the allegation of partisan hijacking unfinished (\texttt{...}) so the audience supplies the "Opposition conspiracy" conclusion without the anchor stating it as fact. \\ \hline
\textbf{3. Intonation, pitch \& volume} & \texttt{[08:48]} \textit{"$^\circ$Dekhiye DCP sahab ke haath mein fracture hai$^\circ$"} vs. \texttt{[06:12]} \textit{"Chhatron ke bhes mein \_UPADRAVIYON\_ ne kanoon haath mein liya$\uparrow$!"} & Softened, pained delivery \texttt{$^\circ$...$^\circ$} on police injury contrasts with loud, staccato stress and rising pitch \texttt{$\uparrow$} on \textit{"UPADRAVIYON"}, inverting victim and aggressor. \\ \hline
\textbf{4. Coughs, sighs \& vocalisations} & \texttt{[10:42]} \texttt{[[cough / throat-clear]]} (pivoting away from police lathi-charge question) / \texttt{[14:45]} \texttt{[[scoff]]} \textit{"Mudda exam hai ya siyasat$\uparrow$?!"} & The \texttt{[cough]} smooths over an awkward moment for the police guest; the \texttt{[scoff]} dismisses protesters' broader accountability demands as pure politics. \\ \hline
\textbf{5. Hand gestures, face \& posture} & \texttt{[08:40]} \texttt{\{\{gesture: injured DCP holds bandaged hand up to camera\}\}} / \texttt{[12:25]} \texttt{\{\{gesture: anchor eye-roll + finger-wag on 'retaliation'\}\}} & The bandaged hand serves as a powerful visual prop legitimising police force, while the anchor's eye-roll and finger-wag delegitimise student self-defence claims. \\ \hline
\end{longtable}

\subsubsection*{Short summary (Event 2 Non-Verbal Synthesis)}

\begin{itemize}
  \item \textbf{What overall interpretation or emotional stance does the non-verbal delivery encourage?}
The non-verbal delivery encourages \textbf{law-and-order alarm and suspicion toward street protesters}. While paying brief lip service to exam grievances, the broadcast's acoustic and visual weight guides the viewer to sympathise with the injured police officer and view the marching students as unruly puppets (\textit{"upadravi"}) of opposition parties.\par\smallskip
  \item \textbf{Whose voice or physical presence is acoustically/visually foregrounded, undermined, or silenced?}
\textbf{The injured Additional DCP (with his bandaged hand in close-up) and Anchor Anjana Om Kashyap} are visually and acoustically foregrounded, given uninterrupted time and sympathetic nods. \textbf{Protesting students} are visually reduced to chaotic B-roll and acoustically \textbf{undermined/silenced} via constant \texttt{[overlap]} interruptions, finger-wagging, and faded microphones (\texttt{[12:18]}).\par\smallskip
  \item \textbf{What evidence supports your non-verbal reading? What complicates it (e.g., speaker baseline, format)?}
The stark contrast between how Anjana treats the DCP's defensive stammer (\texttt{"wo toh--- abhi investigation..."}, \texttt{[10:35]}, met with a gentle nod and pivot) versus the student's explanation (\texttt{[12:18]}, met with an \texttt{[overlap]} shout and eye-roll) provides concrete turn-taking evidence of institutional bias. Format factors (live multi-panel TV debate time pressure) partially explain rapid interruptions, but the asymmetry in who gets interrupted proves ideological alignment.\par\smallskip
  \item \textbf{How do the non-verbal cues interact with the spoken words? (Do they reinforce, dramatise, soften/hedge, invert meaning via sarcasm/irony, or leave an accusation unspoken via trailing off?)}
The non-verbal cues \textbf{soften state violence} (calm, low-volume delivery on \textit{"tear gas chhode gae"} and \textit{"preventive arrangement"} at \texttt{[07:45]}, \texttt{[10:35]}), \textbf{dramatise protester disorder} via loud stress on \textit{"UPADRAVIYON"} (\texttt{[06:12]}) and close-ups of the DCP's bandage (\texttt{[08:40]}), \textbf{invert student claims via sarcasm} (mocking pitch on \textit{"'retaliation'$\uparrow$?!"}, \texttt{[12:25]}), and \textbf{insinuate partisan manipulation via trailing off} (\texttt{...} at \texttt{[11:52]}).\par\smallskip
  \item \textbf{What context, speaker baseline, or format factors support or complicate this interpretation? (e.g., Is this the speaker's normal style across all videos, or a marked shift on this specific actor/claim?)}
Across both Event 1 and Event 2, Anjana Om Kashyap's baseline is high-energy, authoritative Hindi TV anchoring. Yet within both videos, her delivery shifts systematically across actors: \textbf{softened, sighing, and protective toward state/temple authorities} (\texttt{$^\circ$Prabhu Shri Ram$^\circ$}, \texttt{$^\circ$DCP sahab ke haath mein fracture$^\circ$}), and \textbf{combative, high-pitched \texttt{$\uparrow$}, and accusatory toward opposition leaders and street protesters}.\par\smallskip
  \item \textbf{How does another outlet perform the same moment, actor, or claim non-verbally?}
Compare \textbf{Aaj Tak} (\texttt{Event 2}) directly with \textbf{The Wire} (\texttt{Event 2}) and \textbf{Dhruv Rathee} (\texttt{Event 2}): whereas Aaj Tak zooms in on a police officer's bandaged hand with pained vocal sympathy and cuts off protesting students \texttt{[overlap]}, The Wire and Dhruv Rathee give uninterrupted close-up microphone space to bruised and exhausted students while using skeptical rising intonation \texttt{$\uparrow$} and documentary FIRs to challenge police and NTA claims.\par\smallskip
  \item \textbf{My interpretation, with reasons:}
Aaj Tak's multimodal discourse across Events 1 and 2 operates as a \textbf{shield for institutional authority}. In Event 1, physical air-quotes (\texttt{\{\{air-quotes on 'loota'\}\}}), devotional sighs \texttt{[sigh]}, and self-repairs (\texttt{---}) protect the Ram Mandir Trust by isolating blame onto a single clerk and attacking the Opposition. In Event 2, visual props (the DCP's bandaged hand), asymmetric turn-taking \texttt{[overlap]}, and prosodic shifts from \textit{"chhatra"} to \textit{"\_UPADRAVI\_"} transform a crisis of state exam failure into a spectacle of street indiscipline.\par\smallskip
\end{itemize}
\newpage
\section*{Non-Verbal \& Multimodal Discourse Analysis Worksheet}

\textit{Qualitative analysis of pauses, trailing sentences, prosody/intonation, vocalisations (coughs, sighs, scoffs), and kinesics (hand gestures, facial expressions, posture) in political news discourse}\par\smallskip

\bigskip\hrule\bigskip

\subsection*{Guidelines: Five non-verbal \& paralinguistic features to look for in every event}

\begin{longtable}{|p{0.23\textwidth}|p{0.23\textwidth}|p{0.23\textwidth}|p{0.23\textwidth}|}
\hline
Feature & What to look for & Illustrative example (from sample) & Questions to answer \\ \hline \endhead
\textbf{1. Pauses, Silences \& Abrupt Stops} & Micro-pauses \texttt{(.)}, timed silences \texttt{(1.5)}, pregnant/dramatic pauses before a loaded word, hesitant stumbles, abrupt self-interruptions or cut-offs (\texttt{---}), and forced silences (muting a panelist's mic or talking over them). & Anchor pauses \texttt{(1.5)} before saying \textit{"'loota'"} for dramatic weight; Addl. DCP cuts off mid-sentence (\textit{"wo toh--- abhi investigation chal raha hai"}) when asked if police struck first; panelist's mic faded during counter-question. & Where does the pause or stop occur? Is it deliberate/rhetorical (for suspense, emphasis, or gravitas) or defensive/evasive (cognitive load, self-repair, avoiding blame)? Who controls turn-taking and silence? \\ \hline
\textbf{2. Trailing Sentences \& Unfinished Utterances} & Clauses that trail off (\texttt{...}), drop in volume at the end, or leave the conclusion unspoken (aposiopesis) so the audience supplies the implication without the speaker taking explicit responsibility. & \textit{"Aur phir wahan jo log ghuse the... aap samajh hi sakte hain..."}; \textit{"500 salon ke sangharsh ke baad agar aisa..."} trailing off with a slow head shake. & What exact thought is left unfinished? What presupposition or shared bias is the viewer invited to fill in? Does trailing off act as a safe hedge to imply higher-level complicity or "anti-national" motive without stating it? \\ \hline
\textbf{3. Intonation, Pitch, Volume \& Prosody} & Pitch contour (rising \texttt{$\uparrow$} for incredulity/challenge; falling \texttt{$\downarrow$} for finality), volume shifts (shouting vs. hushed whisper), emphatic word stress, speech rate (rapid-fire urgency vs. slow solemnity), and tonal register (sarcastic, reverential, combative). & Hushed, reverential cadence on \textit{"Prabhu Shri Ram ke Mandir"} vs. sharp, high-pitched interrogation on \textit{"par Akhilesh ji Ram Mandir kyon$\uparrow$ nahin jate?!"}; satirical sing-song inflection on \textit{"priventiv arenjment"}. & Which exact words receive pitch peaks or volume stress? How does vocal tone turn neutral words into irony, mockery, or moral outrage? Does the speaker shift prosodic register depending on which actor is discussed? \\ \hline
\textbf{4. Coughs, Sighs \& Non-Lexical Vocalisations} & Audible sighs (\texttt{[sigh]}), scoffs/snorts (\texttt{[scoff]}), ironic chuckles (\texttt{[chuckle]}), throat-clearing/coughs (\texttt{[cough]}), sharp intakes of breath/gasps (\texttt{[gasp]}), and tongue-clicks (\texttt{[tsk]}). & Heavy \texttt{[sigh]} before reporting Trust oversight lapses (signalling reluctant sorrow rather than hostility); mocking \texttt{[scoff]} or \texttt{[chuckle]} when reading an opponent's denial; \texttt{[cough / throat-clear]} when pivoting away from an awkward fact. & What emotion or epistemic stance does the vocalisation perform (exasperation, moral disgust, disbelief, ridicule, discomfort)? Does a sigh or scoff pre-judge a speaker's claim before any verbal rebuttal is made? \\ \hline
\textbf{5. Hand Gestures, Facial Expressions \& Posture (Kinesics)} & Hand movements (air-quotes, finger-pointing/wagging, desk-thumping, open palms, counting on fingers, dismissive wave), facial cues (raised eyebrows, eye-rolls, smirk, furrowed brow), head nods/shakes, leaning into camera, and visual props (bandaged limbs, documents). & Anchor forms physical air-quotes on \textit{"'loota'"} or \textit{"'peaceful' protest"}; host points finger/pen at lens during an ultimatum; injured DCP holds bandaged hand up to camera; creator raises one eyebrow and leans in for a conspiratorial reveal. & How do gestures reinforce, contradict, or sharpen the spoken words? During split-screen listening, does the host's face (nodding vs. smirking/eye-rolling) cue the audience on whom to trust? How are bodily visuals used as evidence? \\ \hline
\end{longtable}

\subsection*{Quick Multimodal Transcription Notation Key (use inside transcript excerpts)}

\begin{longtable}{|p{0.15\textwidth}|p{0.15\textwidth}|p{0.15\textwidth}|p{0.15\textwidth}|p{0.15\textwidth}|p{0.15\textwidth}|}
\hline
Symbol & Meaning & Example in transcript & Symbol & Meaning & Example in transcript \\ \hline \endhead
\texttt{(.)} & Brief micro-pause (< 0.5 sec) & \texttt{woh (.) suparavijan lacking tha} & \texttt{$\uparrow$ / $\downarrow$} & Marked rising / falling pitch & \texttt{kyon$\uparrow$ FIR nahin karaee$\downarrow$} \\ \hline
\texttt{(1.5)} & Timed pause in seconds & \texttt{kisne (1.5) 'loota'?} & \texttt{WORD / \_word\_} & Loud volume / emphatic stress & \texttt{ye PAAP hai / \_bilkul\_ galat} \\ \hline
\texttt{---} & Abrupt stop, cut-off, or self-repair & \texttt{police ki taraf se--- abhi jaanch...} & \texttt{$^\circ$word$^\circ$} & Softened / whispered delivery & \texttt{$^\circ$ander ki baat ye hai$^\circ$} \\ \hline
\texttt{...} & Trailing off / unfinished sentence & \texttt{aap jaante hain ki kaun log...} & \texttt{[sigh] [cough] [scoff]} & Non-lexical vocalisations & \texttt{[sigh] 500 salon ke baad...} \\ \hline
\texttt{[overlap]} & Simultaneous / overlapping speech & \texttt{[Anjana: suniye---] [Guest: nahin!]} & \texttt{\{\{gesture: ...\}\}} & Hand gesture / facial cue / posture & \texttt{\{\{air-quotes on 'loota'\}\}} \\ \hline
\end{longtable}

\textbf{Methodological Reminder:} A non-verbal cue is not automatically manipulative or ideological---a cough, pause, or rapid cadence may reflect studio audio delay, live time pressure, or a speaker's habitual style. Always check (1) whether the cue clusters with verbal features, and (2) whether the speaker's delivery shifts noticeably across different actors or events (within-outlet baseline).\par\smallskip

\bigskip\hrule\bigskip

\subsection*{Event 1. Ram Mandir donation controversy (Non-Verbal \& Paralinguistic Layer)}

\begin{itemize}
  \item \textbf{Focus:} How do prosody, pauses, trailing sentences, vocalisations (sighs, coughs, scoffs), and body language frame the controversy---as solemn devotional sorrow, administrative scrutiny, or partisan combat?
  \item \textbf{Outlet:} Republic TV / Republic Bharat (Prime-Time Bulletin \& Debate Segment) | \textbf{Date:} 14 June 2021
  \item \textbf{Video title:} Ram Mandir Land Deal Allegations: Truth vs Conspiracy / Champat Rai \& Trust Response | \href{https://youtu.be/0H-bruoEkaY}{https://youtu.be/0H-bruoEkaY} | \textbf{Timestamp(s):} \texttt{[00:00]}--\texttt{[09:45]}
  \item \textbf{Analyst:} Baani Solanki (\texttt{2025114016}) | \textbf{Language:} Hindi / English mix (transcribed in Roman script)
  \item \textbf{Format / Visual setup:} \texttt{[x]} TV Panel Debate \texttt{[x]} TV Bulletin \texttt{[ ]} Solo YouTube Explainer \texttt{[ ]} Interview / Discussion
  \item \textbf{Primary speaker(s) observed:} Republic Anchor (standing/seated in high-alert red-and-yellow breaking-news studio), Champat Rai \& Anil Mishra (Trust spokespersons), and BJP vs. SP/AAP panelists
\end{itemize}

\subsubsection*{Event-specific non-verbal questions (Event 1)}

\noindent\textbf{1.} \textbf{Vocal \& gestural marking of core labels ("scam", "loota", "paap", "gadabdi", "avyavastha"): When the controversy is named, what non-verbal cues accompany the label? (Look for physical air-quotes, vocal scare-quotes, dramatic pauses before the word, hushed/reverential drops in volume, or shouted stress.)}\par\smallskip
\begin{itemize}
  \item \textbf{Notes:}
Republic TV never uses \texttt{"scam"} or \texttt{"ghotala"} as a factual description of the land transaction; instead, the word \texttt{"ghotala"} is always accompanied by \textbf{mocking vocal scare-quotes, physical air-quotes, or shouted counter-labels (\texttt{"SHADYANTRA"}, \texttt{"PROPAGANDA"}, \texttt{"JHOOTH"})}:\par\smallskip
  \item At \texttt{[00:25]}, the anchor introduces the opposition's allegation with a \texttt{(1.0)} pause and exaggerated two-handed air-quotes: \textit{"Ram Mandir ke naam par (1.0) 'ghotala'$\uparrow$?! Ya phir Mandir nirmaan rokne ki \_SAZISH\_$\downarrow$!"} (\texttt{[00:29]}).
  \item The word \textit{"\_SAZISH\_"} (conspiracy) receives maximum volume stress (\texttt{WORD}), a sharp downward fist-thump on the desk \texttt{\{\{gesture: right fist thumping desk\}\}}, and a fierce scowl, immediately inverting the story from a financial audit into an ideological war against the Temple.
\end{itemize}

\noindent\textbf{2.} \textbf{Differential demeanour toward actors (Temple, Trust leadership, counting clerks, donors, opposition accusers): How do the speaker's voice pitch, facial expression, and posture shift when referring to Prabhu Shri Ram / devotees vs. top figures (Champat Rai, Nripendra Misra, CM Yogi) vs. arrested counting staff vs. opposition leaders (Akhilesh Yadav, Awadhesh Prasad)?}\par\smallskip
\begin{itemize}
  \item \textbf{Notes:}
Republic exhibits extreme non-verbal polarisation between two moral camps:\par\smallskip
  \item \textbf{Prabhu Shri Ram, Devotees \& Trust Leadership (Champat Rai, Anil Mishra):} Treated as a single, indivisible sacred entity. When introducing Champat Rai's clarification (\texttt{[02:15]}), the anchor's pitch lowers into an \textbf{earnest, respectful, affirmative register}, posture upright, nodding continuously as Trust documents are shown (\texttt{"Trust ne ek-ek paise ka hisaab diya hai$\downarrow$"}, \texttt{[02:22]}).
  \item \textbf{Opposition Accusers (Sanjay Singh, Pawan Pandey, Akhilesh Yadav, Congress):} Addressed with \textbf{high-decibel combative shouting, mocking \texttt{[scoff]} vocalisations, rapid-fire finger-wagging directly at the lens, and aggressive forward leaning} (\texttt{"Jo log kal tak Ram ke astitva par sawaal uthate the, aaj unhe chande ki chinta ho rahi hai$\uparrow$?!"}, \texttt{[03:40]}).
\end{itemize}

\noindent\textbf{3.} \textbf{Solemn devotional grief vs. high-decibel partisan spectacle: How do pauses, sighs, speech rate, and hand gestures construct the segment's emotional register? Where is delivery slow, hushed, and mournful (sighs, bowed head, folded hands), and where does it escalate into rapid-fire shouting, desk-thumping, or mocking irony?}\par\smallskip
\begin{itemize}
  \item \textbf{Notes:}
Unlike Aaj Tak (which opens with slow, mournful devotional grief), Republic TV operates almost entirely in \textbf{high-decibel partisan spectacle} (\textasciitilde{}170--185 words per minute) right from \texttt{[00:00]}. There are virtually no quiet, mournful silences; pauses are short, sharp rhetorical gasps \texttt{[gasp]} used to reload before the next shouted clause (\texttt{[01:10]}, \texttt{[04:15]}). Hand gestures are large, percussive, and militaristic: continuous \textbf{desk-thumping, karate-chops in the air, and stabbing index-finger points} at the camera and at opposition split-screen boxes (\texttt{[05:20]}--\texttt{[08:30]}).\par\smallskip
\end{itemize}

\noindent\textbf{4.} \textbf{Non-verbal evaluation of financial scrutiny: When questions about unreceipted offerings, lack of frisking, or land deals are raised, how does the host/speaker react non-verbally (earnest nodding, leaning into the camera, open-palm appeals vs. eye-rolls, scoffs, dismissive hand waves, or incredulous rising pitch $\uparrow$)?}\par\smallskip
\begin{itemize}
  \item \textbf{Notes:}
Any question about the 5-minute price jump from Rs 2 crore to Rs 18.5 crore is met non-verbally with \textbf{dismissive hand waves, eye-rolls, loud \texttt{[scoff]} vocalisations, and incredulous rising pitch \texttt{$\uparrow$}}:\par\smallskip
  \item When an opposition panelist attempts to read out the \texttt{7:10 PM} vs. \texttt{7:15 PM} registry timestamps (\texttt{[05:45]}), the anchor rolls their eyes, waves the back of their hand dismissively across the desk \texttt{\{\{gesture: backhanded dismissive wave + eye-roll\}\}}, and shouts over them \texttt{[overlap]}: \textit{"Purana registered agreement aapko dikhai nahin deta$\uparrow$?! Ya padhna nahin aata$\uparrow$?!"} (\texttt{[05:52]}).
  \item Conversely, when the Trust spokesperson explains that the price was below market rate and paid digitally via RTGS (\texttt{[02:40]}--\texttt{[03:15]}), the anchor leans in with \textbf{earnest, vigorous head nods} \texttt{\{\{gesture: vigorous affirmative nodding\}\}} and taps the desk to endorse the point.
\end{itemize}

\noindent\textbf{5.} \textbf{Performing certainty vs. hedging and innuendo (pauses, stops, trailing off): Are allegations delivered with definitive falling intonation $\downarrow$, finger-counting, and pointing at on-screen evidence (CCTV, documents), or hedged with hesitant pauses (.), abrupt self-interruptions (---), throat-clearing [cough], and trailing sentences (...) that imply higher-level complicity without naming it?}\par\smallskip
\begin{itemize}
  \item \textbf{Notes:}
Republic performs \textbf{100\% epistemic certainty in exonerating the Trust} and \textbf{conspiratorial certainty in indicting the Opposition}:\par\smallskip
  \item \textbf{Exoneration of the Trust:} Delivered with booming falling intonation \texttt{$\downarrow$} and finger-counting (\texttt{"Number ek--- saara paisa bank se gaya$\downarrow$! Number do--- purana agreement 2011 se chal raha tha$\downarrow$!"}, \texttt{[03:05]}).
  \item \textbf{Innuendo / Conspiracy against Opposition:} Delivered via dramatic trailing clauses (\texttt{...}) and lowered, menacing pitch: \textit{"Chunav se theek pehle (.) Ayodhya ko badnaam karne ke liye ye khel khela ja raha hai... aur iske peechhe kaun si taqatein hain..."} (\texttt{[04:50]}), trailing off with a slow, menacing point at the lens \texttt{\{\{gesture: sustained finger-point at lens + narrowed eyes\}\}} to imply an "anti-Hindu / anti-national" cabal.
\end{itemize}

\noindent\textbf{6.} \textbf{Turn-taking, interruptions, and acoustic treatment of denials: In debates/interviews: who receives calm, uninterrupted speaking time and sympathetic nods, and who is cut off mid-sentence (---), talked over, or muted? In solo videos: how are clips of official denials framed via reaction pauses, raised eyebrows, or sarcastic chuckles?}\par\smallskip
\begin{itemize}
  \item \textbf{Notes:}
Turn-taking is \textbf{radically asymmetric and weaponised}:\par\smallskip
  \item \textbf{Trust Spokespersons \& BJP Defenders (\texttt{[02:15]}--\texttt{[03:30]}, \texttt{[06:30]}):} Receive 60--90 seconds of completely uninterrupted speaking time, clean solo audio, full-screen or dominant window placement, and affirmative verbal backchannels (\textit{"Bilkul sahi"}, \textit{"Yes"}).
  \item \textbf{Opposition Panelists (SP / AAP, \texttt{[05:35]}--\texttt{[08:15]}):} Cut off mid-sentence \texttt{---} within 5--8 seconds of speaking. The anchor repeatedly talks over them \texttt{[overlap]} at superior volume (\textit{"[Anchor: Jhooth mat boliye! Ek minute--- jhooth mat boliye!] [SP Guest: Kagaz toh dekhiye---]"}, \texttt{[06:05]}), while the studio audio engineer ducks the opposition guest's microphone volume by \textasciitilde{}10 dB so their rebuttal is reduced to unintelligible background mumbles.
\end{itemize}

\noindent\textbf{7.} \textbf{Performing the boundary between religious devotion and institutional critique: How does the presenter use vocal tone, sighs, and body language to signal personal reverence for the Temple (or secular objectivity) while criticising---or defending---the Trust's financial management?}\par\smallskip
\begin{itemize}
  \item \textbf{Notes:}
Republic \textbf{erases the boundary between religious devotion and institutional management} entirely. Through non-verbal fusion---placing a hand over the heart when saying \textit{"Shri Ram Janmabhoomi Teerth Kshetra Trust"} (\texttt{[01:35]}) and immediately pointing an angry finger when saying \textit{"Mandir par hamla"} (\texttt{[01:42]})---the anchor equates questioning Champat Rai's real-estate signatory decisions with attacking Lord Ram Himself. Any institutional audit is performed non-verbally as sacrilege.\par\smallskip
\end{itemize}

\noindent\textbf{8.} \textbf{Self-presentation as the "guardian of public trust": How does the speaker use direct eye-gaze, posture, moral sighs, or emphatic gestures (pointing at the viewer, hands on heart/desk) to position themselves as the authentic protector of devotees or public accountability?}\par\smallskip
\begin{itemize}
  \item \textbf{Notes:}
The anchor presents themselves as a \textbf{belligerent civilisational warrior defending Hindu devotees (\texttt{"karodon Ram bhakton ki aastha"}) against a hostile political cabal}. Standing/leaning aggressively into the camera, gripping a sheaf of papers in one hand and slapping them on the desk \texttt{\{\{gesture: slapping sheaf of papers on desk\}\}} (\texttt{[07:10]}), the host stares directly into the lens and shouts: \textit{"Desh ki janta is saazish ko \_kabhi\_ maaf nahin karegi$\downarrow$!"} (\texttt{[08:55]}).\par\smallskip
\end{itemize}

\subsubsection*{Apply the five non-verbal features (Event 1)}

\begin{longtable}{|p{0.31\textwidth}|p{0.31\textwidth}|p{0.31\textwidth}|}
\hline
Non-verbal feature & Timestamp + spoken anchor words + non-verbal notation & What it does in context (stance, framing, implication) \\ \hline \endhead
\textbf{1. Pauses, silences \& stops} & \texttt{[00:25]} \textit{"Ram Mandir ke naam par (1.0) 'ghotala'$\uparrow$?!"} / \texttt{[06:05]} \texttt{[overlap]} \textit{"[Anchor: Jhooth mat boliye--- chup rahiye!] [Guest: Do crore ka---]"} (guest mic faded) & The \texttt{(1.0)} pause sets up sarcastic air-quotes on \texttt{'ghotala'}, while forced silence via \texttt{[overlap]} shouting and audio ducking blocks the guest from stating the Rs 2 crore figure. \\ \hline
\textbf{2. Trailing sentences (\texttt{...})} & \texttt{[04:50]} \textit{"Chunav se theek pehle Ayodhya ko badnaam karne ke liye... aap jaante hain kaun log baukhlaye hue hain..."} & Trails off with a menacing stare to presuppose that anyone questioning the land deal is driven by electoral panic and anti-temple malice. \\ \hline
\textbf{3. Intonation, pitch \& volume} & \texttt{[00:29]} \textit{"Mandir nirmaan rokne ki \_SAZISH\_$\downarrow$!"} (maximum volume \texttt{WORD} + falling finality \texttt{$\downarrow$}) vs. \texttt{[05:52]} \textit{"Padhna nahin aata$\uparrow$?!"} (sharp mocking rise \texttt{$\uparrow$}) & Uses booming falling stress \texttt{$\downarrow$} to declare an opposition conspiracy as established fact, and mocking rising pitch \texttt{$\uparrow$} to humiliate dissenting panelists. \\ \hline
\textbf{4. Coughs, sighs \& vocalisations} & \texttt{[01:10]} \texttt{[[gasp]]} (sharp intake of breath between rapid shouted clauses) / \texttt{[03:45]} \texttt{[[scoff]]} \textit{"Aaj inhe chande ki chinta ho rahi hai!"} & Rapid \texttt{[gasp]} intakes sustain high-arousal urgency; the loud \texttt{[scoff]} ridicules the sincerity of opposition leaders' concern for donors. \\ \hline
\textbf{5. Hand gestures, face \& posture} & \texttt{[00:25]} \texttt{\{\{gesture: two-handed air-quotes on 'ghotala'\}\}} / \texttt{[00:29]} \texttt{\{\{gesture: right fist thumping desk\}\}} / \texttt{[07:10]} \texttt{\{\{gesture: slapping papers on desk + pointing at lens\}\}} & Air-quotes delegitimise the allegation; desk-thumping and paper-slapping perform prosecutorial dominance and righteous nationalist anger. \\ \hline
\end{longtable}

\subsubsection*{Short summary (Event 1 Non-Verbal Synthesis)}

\begin{itemize}
  \item \textbf{What overall interpretation or emotional stance does the non-verbal delivery encourage?}
The non-verbal delivery encourages \textbf{combative, defensive outrage against the Opposition}. By pairing air-quotes on \textit{"'ghotala'"} with fist-thumping stress on \textit{"\_SAZISH\_"} and relentless \texttt{[overlap]} interruptions of critics, the broadcast leaves zero room for administrative doubt, framing the entire controversy as a fabricated attack on the Ram Mandir.\par\smallskip
  \item \textbf{Whose voice or physical presence is acoustically/visually foregrounded, undermined, or silenced?}
\textbf{The studio anchor and Trust leadership (Champat Rai / Anil Mishra)} are acoustically and visually foregrounded with uninterrupted solo segments and vigorous affirmative nods. \textbf{Opposition spokespersons (SP/AAP)} are systematically \textbf{silenced}---cut off mid-sentence \texttt{---} within seconds, shouted over \texttt{[overlap]}, and acoustically ducked in the studio mix (\texttt{[06:05]}).\par\smallskip
  \item \textbf{What evidence supports your non-verbal reading? What complicates it (e.g., speaker baseline, format)?}
The turn-taking data---60+ seconds of uninterrupted calm audio for Trust defenders (\texttt{[02:15]}--\texttt{[03:30]}) versus <8 seconds before \texttt{[overlap]} interruption for critics (\texttt{[05:35]}--\texttt{[08:15]})---is objective acoustic proof of gatekeeping. While high-decibel shouting (\textasciitilde{}175 WPM) is Republic TV's general prime-time baseline, the complete absence of interruptions when Trust representatives speak shows that the anchor's aggression is selectively deployed.\par\smallskip
\end{itemize}

\bigskip\hrule\bigskip

\subsection*{Event 2. Students' protest regarding NEET paper leak (Non-Verbal \& Paralinguistic Layer)}

\begin{itemize}
  \item \textbf{Focus:} How do prosody, pauses, trailing sentences, vocalisations, and visual body language construct the legitimacy of student protesters, police use of force, and state accountability?
  \item \textbf{Outlet:} Republic TV (\textit{The Debate with Arnab Goswami}) | \textbf{Date:} June 2024
  \item \textbf{Video title:} NEET Paper Leak \& Student Protest Debate | Arnab Goswami on NTA \& Political Row | \href{https://youtu.be/HTvb6E\_1Qdc}{https://youtu.be/HTvb6E\_1Qdc} | \textbf{Timestamp(s):} \texttt{[00:00]}--\texttt{[36:20]}
  \item \textbf{Analyst:} Baani Solanki (\texttt{2025114016}) | \textbf{Language:} English (with occasional Hindi terms in Roman script)
  \item \textbf{Format / Visual setup:} \texttt{[x]} TV Panel Debate \texttt{[ ]} TV Bulletin \texttt{[ ]} Solo YouTube Explainer \texttt{[ ]} Interview / Discussion
  \item \textbf{Primary speaker(s) observed:} Arnab Goswami (studio anchor seated/leaning forward at desk), surrounded by an 8-window split-screen grid of Student Representatives, Educators, BJP/NDA Spokespersons, and Opposition (Congress/RJD/SP) Panelists
\end{itemize}

\subsubsection*{Event-specific non-verbal questions (Event 2)}

\noindent\textbf{1.} \textbf{Prosodic \& gestural marking of participant labels ("students/bachche" vs. "protesters" vs. "upadravi/goons"): How do vocal tone, volume, facial expression, and hand gestures shift across labels? (e.g., softened, empathetic pitch for "students" vs. harsh, staccato stress, scowls, or accusatory finger-pointing on "upadraviyon", "goons", or "anti-social elements".)}\par\smallskip
\begin{itemize}
  \item \textbf{Notes:}
Arnab Goswami creates a dramatic \textbf{three-way prosodic and gestural split} across labels:\par\smallskip
  \item \textbf{Genuine NEET Aspirants (\texttt{"young students"}, \texttt{"24 lakh children"}):} Delivered with a \textbf{deliberately slowed, earnest, protective cadence} (\texttt{[00:35]}, \texttt{[04:10]}), both hands held open toward the camera \texttt{\{\{gesture: both palms open and extended toward lens\}\}} and a deeply furrowed brow of paternal concern.
  \item \textbf{Bureaucrats (\texttt{"NTA babus"}, \texttt{"faceless bureaucrats"}):} Delivered with biting, staccato contempt (\texttt{[02:15]}), jabbing his pen downward into the desk \texttt{\{\{gesture: repeated pen-jabs onto desk surface\}\}}.
  \item \textbf{Opposition Youth-Wing Protesters \& Politicians (\texttt{"political vultures"}, \texttt{"hijackers"}, \texttt{"RJD-Congress lobby"}):} Delivered at \textbf{maximum shouted volume (\texttt{WORD}), rapid-fire pace (\textasciitilde{}190 WPM), red-faced scowls, and aggressive two-handed finger-pointing} (\texttt{"Stop hijacking the students' pain! You are \_POLITICAL VULTURES\_$\downarrow$!"}, \texttt{[14:25]}).
\end{itemize}

\noindent\textbf{2.} \textbf{Delivery of students' stated demands: When the protesters' core demands (Education Minister's resignation, compensation, release of Sonam Wangchuk, exam reform) are mentioned, are they delivered slowly with deliberate pauses and finger-counting, rushed through in a flat monotone, mocked with air-quotes/scoffs, or trailed off (...)?}\par\smallskip
\begin{itemize}
  \item \textbf{Notes:}
Arnab selectively \textbf{splits the students' demands into "acceptable administrative demands" vs. "unacceptable political demands"}:\par\smallskip
  \item \textbf{Administrative Demands (Scrapping grace marks, firing NTA officials, CBI probe, \texttt{[03:10]}--\texttt{[04:45]}):} Delivered \textbf{slowly and forcefully with finger-counting} (\texttt{"Point one (.) scrap the grace marks$\downarrow$! Point two (.) overhaul the NTA from top to bottom$\downarrow$!"}, \texttt{[03:22]}), thumping his fist on each number.
  \item \textbf{Political Demands (Resignation of Education Minister Dharmendra Pradhan or street blockades by youth wings, \texttt{[15:10]}):} Immediately \textbf{mocked with a loud \texttt{[scoff]}, dismissive hand-waving, and shouted interruptions} (\texttt{"You don't care about NEET! You only want a political resignation for your election narrative!"}, \texttt{[15:28]}).
\end{itemize}

\noindent\textbf{3.} \textbf{Where does vocal \& kinetic energy peak---grievance, street disruption, or electoral horse-race?: Track where the speaker's volume, pitch range, pace, and gesture intensity reach their climax: during accounts of the paper leak/student suicides, during street clashes and barricade breaches, or during partisan BJP-vs-Opposition sparring.}\par\smallskip
\begin{itemize}
  \item \textbf{Notes:}
While Arnab opens with high vocal energy against the NTA (\texttt{[00:00]}--\texttt{[06:00]}), the \textbf{absolute acoustic and kinetic climax of the 36-minute broadcast occurs during partisan BJP-vs-Opposition sparring over the Bihar paper-leak nexus (\texttt{[12:00]}--\texttt{[28:00]})}:\par\smallskip
  \item When linking the arrested Bihar mastermind Sikandar Yadavendu to RJD leader Tejashwi Yadav's personal secretary Pritam Kumar (\texttt{[13:40]}--\texttt{[18:50]}), Arnab's volume hits peak saturation (\texttt{WORD}), he rises halfway out of his chair \texttt{\{\{gesture: rising half-standing from chair, leaning over desk\}\}}, waves printed call-detail/guest-house papers at the camera, and shouts non-stop for 4 minutes at the RJD/Congress panelists.
\end{itemize}

\noindent\textbf{4.} \textbf{Non-verbal attribution of responsibility (sighs, stops, head shakes, pointing): When examination authorities (NTA), the Education Minister, police, opposition politicians, or students are named as responsible, what non-verbal cues accompany the attribution (exasperated sighs, mocking laughs, abrupt self-stops ---, or direct finger-pointing at the lens)?}\par\smallskip
\begin{itemize}
  \item \textbf{Notes:}
Arnab performs a classic \textbf{non-verbal blame-deflection manoeuvre} that shields the Union Government while attacking both the NTA bureaucracy and the Opposition:\par\smallskip
  \item \textbf{NTA Bureaucrats:} Blamed with exasperated head shakes, \texttt{[sigh]} vocalisations, and pen-pointing (\texttt{"NTA has become a law unto itself"}, \texttt{[02:45]}), isolating failure to an "autonomous agency".
  \item \textbf{Union Education Minister / Modi Government:} Whenever an opposition panelist names Dharmendra Pradhan or PM Modi (\texttt{[16:10]}), Arnab abruptly stops them \texttt{---}, shakes his head vigorously, and pivots (\texttt{"The government has already ordered a CBI probe and enacted a strict anti-paper-leak law$\downarrow$"}, \texttt{[16:22]}).
  \item \textbf{Opposition Politicians (Tejashwi Yadav / RJD):} Blamed for the actual Patna paper leak with \textbf{direct finger-pointing at the lens, mocking \texttt{[chuckle]}/\texttt{[scoff]} vocalisations, and desk-slapping} (\texttt{"Who booked the NHAI guest house in Patna$\uparrow$?! \_PRITAM KUMAR\_$\downarrow$!"}, \texttt{[14:05]}).
\end{itemize}

\noindent\textbf{5.} \textbf{Prosody of questions---open inquiry vs. loaded rhetorical challenge: Are questions about the protest delivered with level, open intonation, or with sharp rising pitch $\uparrow$, sarcastic stress, raised eyebrows, and trailing clauses (...) that presuppose protesters are violent or politically manipulated?}\par\smallskip
\begin{itemize}
  \item \textbf{Notes:}
  \item \textbf{To Student / Educator Guests (\texttt{[07:15]}, \texttt{[09:30]}):} Arnab switches to a \textbf{remarkably soft, solicitous, open-inquiry prosody} (\texttt{"Tell me, young man, what did you feel when you saw 67 toppers$\uparrow$?"}, \texttt{[07:20]}), leaning his chin on his hand and nodding empathetically \texttt{\{\{gesture: chin resting on hand + slow empathetic nod\}\}}.
  \item \textbf{To Opposition / Youth-Wing Protesters (\texttt{[13:50]}, \texttt{[19:20]}):} Delivered as \textbf{explosive, loaded rhetorical challenges} with sharp rising pitch \texttt{$\uparrow$}, sarcastic stress, and accusatory trailing off (\texttt{...}): \textit{"Why is Rahul Gandhi silent on the Bihar guest house link$\uparrow$?! Because the trail leads straight to..."} (\texttt{[14:48]}).
\end{itemize}

\noindent\textbf{6.} \textbf{Visual \& paralinguistic performance of injury and violence: How are claims of violence substantiated non-verbally on screen? (Compare camera close-ups and pained vocal tone on police bandages/fractured hands vs. visual/acoustic treatment of student head injuries, lathi-charge footage, or defensive stammers/pauses when an officer or guest is asked who used force first.)}\par\smallskip
\begin{itemize}
  \item \textbf{Notes:}
Republic splits visual and paralinguistic treatment between \textbf{"innocent aspirants in studios"} and \textbf{"rowdy political cadres on streets"}:\par\smallskip
  \item Ordinary NEET candidates in the studio grid receive steady close-up frames and pained vocal sympathy from Arnab (\texttt{[07:15]}--\texttt{[10:30]}).
  \item Street protest clashes (\texttt{[20:10]}) are framed via split-screen loops showing Congress/NSUI/Left flags at barricades while Arnab points at the party flags on screen \texttt{\{\{gesture: pointing at party flags in B-roll window\}\}} and shouts: \textit{"Look at the flags! These are \_NOT\_ students, these are Congress workers creating anarchy on the streets$\downarrow$!"} (\texttt{[20:25]}), completely dismissing any claims of police excess against youth-wing marchers.
\end{itemize}

\noindent\textbf{7.} \textbf{Asymmetric delivery of police actions vs. protester actions: Compare the vocal cadence, pauses, and hand gestures used when describing police actions (e.g., calm, matter-of-fact administrative tone on "tear gas chhode gae" / "preventive arrangement") versus protester actions (alarmist pitch, striking gestures, or ironic softening as "retaliation").}\par\smallskip
\begin{itemize}
  \item \textbf{Notes:}
  \item \textbf{Police \& CBI Actions:} Narrated in a \textbf{resolute, authoritative, approving cadence} with firm downward hand-chops (\texttt{"Bihar Police and EOU have cracked the nexus$\downarrow$! CBI is taking over$\downarrow$!"}, \texttt{[13:15]}), portraying law enforcement as heroic investigators uncovering opposition-linked mafia rings.
  \item \textbf{Street Protester Actions (Opposition Youth Wings):} Narrated with \textbf{alarmist pitch \texttt{$\uparrow$}, mocking \texttt{[scoff]} vocalisations, and striking gestures} (\texttt{"Pelting stones and climbing barricades will not solve NEET---it is cheap political drama!"}, \texttt{[20:40]}).
\end{itemize}

\noindent\textbf{8.} \textbf{Acoustic \& visual space granted to students' own voices: Are protesting students heard in their own voice with clear audio, steady close-ups, and empathetic listening nods, or are they reduced to noisy background chaos B-roll while studio anchors/politicians speak over them?}\par\smallskip
\begin{itemize}
  \item \textbf{Notes:}
Arnab grants \textbf{clean, protected acoustic space to non-political NEET students/educators} in the first third of the show (\texttt{[07:10]}--\texttt{[11:45]}), explicitly shouting at politicians to stay silent while the student speaks (\textit{"Quiet! Let the student speak! Don't interrupt the student!"}, \texttt{[08:05]}, holding up a flat palm like a traffic stop sign \texttt{\{\{gesture: right palm held up vertically to silence politicians\}\}}). However, once the debate shifts to the street protests and the Bihar RJD link (\texttt{[12:00]}--\texttt{[36:20]}), students in the 8-box grid are \textbf{visually sidelined and acoustically drowned out} for over 20 minutes by non-stop \texttt{[overlap]} shouting between Arnab, BJP, and RJD/Congress panelists.\par\smallskip
\end{itemize}

\noindent\textbf{9.} \textbf{Non-verbal shift between the NEET paper-leak issue and the street protest: Does the speaker's non-verbal demeanour shift from sympathetic concern (sombre tone, sighs, solidarity gestures) when discussing the exam leak itself to combative, authoritarian, or dismissive delivery once students march toward Parliament?}\par\smallskip
\begin{itemize}
  \item \textbf{Notes:}
\textbf{Yes---an extreme bimodal shift occurs.} During the exam-anomaly segment (\texttt{[00:00]}--\texttt{[11:45]}), Arnab performs \textbf{paternal solidarity} with students and anger at NTA bureaucrats (leaning in softly to listen to students, furrowing his brow). The moment the topic shifts to the street protests and political accountability (\texttt{[12:00]} onward), he pivots into \textbf{hyper-combative partisan warfare}: shouting at \textasciitilde{}190 WPM, slapping papers on the desk, and accusing street protesters of being Congress/RJD puppets trying to destabilise the government.\par\smallskip
\end{itemize}

\subsubsection*{Apply the five non-verbal features (Event 2)}

\begin{longtable}{|p{0.31\textwidth}|p{0.31\textwidth}|p{0.31\textwidth}|}
\hline
Non-verbal feature & Timestamp + spoken anchor words + non-verbal notation & What it does in context (stance, framing, implication) \\ \hline \endhead
\textbf{1. Pauses, silences \& stops} & \texttt{[08:05]} \textit{"Quiet! (1.0) Let the student speak$\downarrow$!"} / \texttt{[14:02]} \textit{"Who booked the NHAI guest house in Patna? (1.5) \_PRITAM KUMAR\_, PA to Tejashwi Yadav$\downarrow$!"} & The \texttt{(1.0)} command silence positions Arnab as the protector of students; the \texttt{(1.5)} dramatic pause before naming Tejashwi's aide turns the NEET debate into an indictment of the Opposition. \\ \hline
\textbf{2. Trailing sentences (\texttt{...})} & \texttt{[14:48]} \textit{"Why is Rahul Gandhi suddenly silent on the Patna arrest$\uparrow$?! Because the paper-leak trail leads straight to his alliance partner..."} & Trails off while pointing at the lens, inviting the viewer to conclude that the Opposition itself engineered the paper leak. \\ \hline
\textbf{3. Intonation, pitch \& volume} & \texttt{[07:20]} \textit{"$^\circ$Tell me, young man, how hard did you prepare$^\circ$?"} vs. \texttt{[14:25]} \textit{"You are \_POLITICAL VULTURES\_ hijacking the students!"} & Extreme prosodic contrast between hushed, gentle pitch \texttt{$^\circ$...$^\circ$} for individual students and maximum shouted volume \texttt{WORD} against opposition protesters. \\ \hline
\textbf{4. Coughs, sighs \& vocalisations} & \texttt{[02:45]} \texttt{[[sigh]]} \textit{"My heart goes out to the \texttt{24 lakh} children..."} / \texttt{[15:25]} \texttt{[[scoff] + [chuckle]]} (at demand for Education Minister's resignation) & The opening \texttt{[sigh]} performs populist empathy; the mocking \texttt{[scoff]} dismisses calls for ministerial accountability as partisan opportunism. \\ \hline
\textbf{5. Hand gestures, face \& posture} & \texttt{[08:05]} \texttt{\{\{gesture: vertical stop-palm toward politicians, open palm toward student\}\}} / \texttt{[14:10]} \texttt{\{\{gesture: rising half-standing, waving Bihar guest-house papers\}\}} & The two-handed traffic-cop gesture visually separates "pure students" from "dirty politicians"; waving the Bihar papers redirects anger away from the Centre toward the RJD. \\ \hline
\end{longtable}

\subsubsection*{Short summary (Event 2 Non-Verbal Synthesis)}

\begin{itemize}
  \item \textbf{What overall interpretation or emotional stance does the non-verbal delivery encourage?}
The non-verbal delivery encourages \textbf{populist sympathy for individual exam-takers, anger at nameless "NTA bureaucrats", and fierce partisan hostility toward street protesters and the Opposition}. By splitting "innocent students" from "political hijackers", Arnab channels public anger over the NEET scam away from the Union Government and onto the Bihar Opposition (\texttt{RJD/Congress}).\par\smallskip
  \item \textbf{Whose voice or physical presence is acoustically/visually foregrounded, undermined, or silenced?}
\textbf{Arnab Goswami} dominates the acoustic floor throughout. \textbf{Individual non-political students} are foregrounded briefly (\texttt{[07:10]}--\texttt{[11:45]}) with clean audio and empathetic nods, after which they sit silently in small grid boxes while \textbf{Opposition spokespersons and street protesters} are shouted over \texttt{[overlap]}, mocked with scoffs, and visually tagged with party flags.\par\smallskip
  \item \textbf{What evidence supports your non-verbal reading? What complicates it (e.g., speaker baseline, format)?}
The stark acoustic contrast between \texttt{[07:20]} (\textasciitilde{}115 WPM, softened volume \texttt{$^\circ$...$^\circ$}, chin-on-hand listening posture for a student) and \texttt{[14:05]} (\textasciitilde{}190 WPM, shouted \texttt{WORD}, half-standing desk-slapping on the RJD guest-house link) proves a strategic within-broadcast shift. While Arnab's shouting baseline is famous across all \textit{Debate} episodes, his ability to drop into quiet, uninterrupted listening for a student guest shows that his high-decibel \texttt{[overlap]} interruptions of opposition guests are deliberate turn-taking control.\par\smallskip
  \item \textbf{How do the non-verbal cues interact with the spoken words? (Do they reinforce, dramatise, soften/hedge, invert meaning via sarcasm/irony, or leave an accusation unspoken via trailing off?)}
His non-verbal cues \textbf{dramatise bureaucratic failure} (pen-jabbing on \textit{"NTA babus"}), \textbf{soften central government culpability} (calm, reassuring nods when mentioning the CBI probe at \texttt{[16:22]}), \textbf{invert street protests via visual pointing} (pointing at party flags to rebrand protesters as \textit{"Congress cadres"} at \texttt{[20:25]}), and \textbf{implicate the Opposition via dramatic \texttt{(1.5)} pauses and trailing off} (\texttt{...} at \texttt{[14:48]}).\par\smallskip
  \item \textbf{What context, speaker baseline, or format factors support or complicate this interpretation? (e.g., Is this the speaker's normal style across all videos, or a marked shift on this specific actor/claim?)}
Comparing Event 1 and Event 2 shows Republic TV's signature \textbf{"populist-nationalist shield"}: in Event 1 (where a sacred institution is accused), the channel denies any wrongdoing whatsoever (\texttt{"100\% sazish"}). In Event 2 (where \texttt{24 lakh} students and middle-class families are angry and denial is impossible), Arnab adopts a softer, empathetic posture toward the victims for the first 10 minutes before deploying his standard high-decibel prosecutorial arsenal to blame autonomous bureaucrats and an Opposition-linked paper-leak nexus in Bihar.\par\smallskip
  \item \textbf{How does another outlet perform the same moment, actor, or claim non-verbally?}
Whereas \textbf{Dhruv Rathee} (\texttt{Event 2}) and \textbf{The Wire} (\texttt{Event 2}) use document-pointing and rising pitch \texttt{$\uparrow$} to connect the paper leaks in Bihar and Gujarat (Godhra) directly to the Union Education Minister and ruling-party states, \textbf{Republic TV} (\texttt{Event 2}) waves documents exclusively about the Patna NHAI guest house (\texttt{[14:05]}) with booming stress on \textit{"\_PRITAM KUMAR\_"} (Tejashwi Yadav's aide), using non-verbal spectacle to turn a central exam failure into an opposition scandal.\par\smallskip
  \item \textbf{My interpretation, with reasons:}
Republic TV's multimodal discourse across both events demonstrates \textbf{acoustic domination and strategic blame redirection}. Through \texttt{[overlap]} shouting, selective mic-ducking, desk-thumping, and split-screen visual framing, the channel neutralises threats to establishment legitimacy---either by branding financial audit questions as an "anti-Mandir conspiracy" (Event 1) or by splitting sympathetic exam aspirants from "vulture" street protesters and pinning the paper leak on the Opposition (Event 2).\par\smallskip
\end{itemize}
\newpage
\section*{Non-Verbal \& Multimodal Discourse Analysis Worksheet}

\textit{Qualitative analysis of pauses, trailing sentences, prosody/intonation, vocalisations (coughs, sighs, scoffs), and kinesics (hand gestures, facial expressions, posture) in political news discourse}\par\smallskip

\bigskip\hrule\bigskip

\subsection*{Guidelines: Five non-verbal \& paralinguistic features to look for in every event}

\begin{longtable}{|p{0.23\textwidth}|p{0.23\textwidth}|p{0.23\textwidth}|p{0.23\textwidth}|}
\hline
Feature & What to look for & Illustrative example (from sample) & Questions to answer \\ \hline \endhead
\textbf{1. Pauses, Silences \& Abrupt Stops} & Micro-pauses \texttt{(.)}, timed silences \texttt{(1.5)}, pregnant/dramatic pauses before a loaded word, hesitant stumbles, abrupt self-interruptions or cut-offs (\texttt{---}), and forced silences (muting a panelist's mic or talking over them). & Anchor pauses \texttt{(1.5)} before saying \textit{"'loota'"} for dramatic weight; Addl. DCP cuts off mid-sentence (\textit{"wo toh--- abhi investigation chal raha hai"}) when asked if police struck first; panelist's mic faded during counter-question. & Where does the pause or stop occur? Is it deliberate/rhetorical (for suspense, emphasis, or gravitas) or defensive/evasive (cognitive load, self-repair, avoiding blame)? Who controls turn-taking and silence? \\ \hline
\textbf{2. Trailing Sentences \& Unfinished Utterances} & Clauses that trail off (\texttt{...}), drop in volume at the end, or leave the conclusion unspoken (aposiopesis) so the audience supplies the implication without the speaker taking explicit responsibility. & \textit{"Aur phir wahan jo log ghuse the... aap samajh hi sakte hain..."}; \textit{"500 salon ke sangharsh ke baad agar aisa..."} trailing off with a slow head shake. & What exact thought is left unfinished? What presupposition or shared bias is the viewer invited to fill in? Does trailing off act as a safe hedge to imply higher-level complicity or "anti-national" motive without stating it? \\ \hline
\textbf{3. Intonation, Pitch, Volume \& Prosody} & Pitch contour (rising \texttt{$\uparrow$} for incredulity/challenge; falling \texttt{$\downarrow$} for finality), volume shifts (shouting vs. hushed whisper), emphatic word stress, speech rate (rapid-fire urgency vs. slow solemnity), and tonal register (sarcastic, reverential, combative). & Hushed, reverential cadence on \textit{"Prabhu Shri Ram ke Mandir"} vs. sharp, high-pitched interrogation on \textit{"par Akhilesh ji Ram Mandir kyon$\uparrow$ nahin jate?!"}; satirical sing-song inflection on \textit{"priventiv arenjment"}. & Which exact words receive pitch peaks or volume stress? How does vocal tone turn neutral words into irony, mockery, or moral outrage? Does the speaker shift prosodic register depending on which actor is discussed? \\ \hline
\textbf{4. Coughs, Sighs \& Non-Lexical Vocalisations} & Audible sighs (\texttt{[sigh]}), scoffs/snorts (\texttt{[scoff]}), ironic chuckles (\texttt{[chuckle]}), throat-clearing/coughs (\texttt{[cough]}), sharp intakes of breath/gasps (\texttt{[gasp]}), and tongue-clicks (\texttt{[tsk]}). & Heavy \texttt{[sigh]} before reporting Trust oversight lapses (signalling reluctant sorrow rather than hostility); mocking \texttt{[scoff]} or \texttt{[chuckle]} when reading an opponent's denial; \texttt{[cough / throat-clear]} when pivoting away from an awkward fact. & What emotion or epistemic stance does the vocalisation perform (exasperation, moral disgust, disbelief, ridicule, discomfort)? Does a sigh or scoff pre-judge a speaker's claim before any verbal rebuttal is made? \\ \hline
\textbf{5. Hand Gestures, Facial Expressions \& Posture (Kinesics)} & Hand movements (air-quotes, finger-pointing/wagging, desk-thumping, open palms, counting on fingers, dismissive wave), facial cues (raised eyebrows, eye-rolls, smirk, furrowed brow), head nods/shakes, leaning into camera, and visual props (bandaged limbs, documents). & Anchor forms physical air-quotes on \textit{"'loota'"} or \textit{"'peaceful' protest"}; host points finger/pen at lens during an ultimatum; injured DCP holds bandaged hand up to camera; creator raises one eyebrow and leans in for a conspiratorial reveal. & How do gestures reinforce, contradict, or sharpen the spoken words? During split-screen listening, does the host's face (nodding vs. smirking/eye-rolling) cue the audience on whom to trust? How are bodily visuals used as evidence? \\ \hline
\end{longtable}

\subsection*{Quick Multimodal Transcription Notation Key (use inside transcript excerpts)}

\begin{longtable}{|p{0.15\textwidth}|p{0.15\textwidth}|p{0.15\textwidth}|p{0.15\textwidth}|p{0.15\textwidth}|p{0.15\textwidth}|}
\hline
Symbol & Meaning & Example in transcript & Symbol & Meaning & Example in transcript \\ \hline \endhead
\texttt{(.)} & Brief micro-pause (< 0.5 sec) & \texttt{woh (.) suparavijan lacking tha} & \texttt{$\uparrow$ / $\downarrow$} & Marked rising / falling pitch & \texttt{kyon$\uparrow$ FIR nahin karaee$\downarrow$} \\ \hline
\texttt{(1.5)} & Timed pause in seconds & \texttt{kisne (1.5) 'loota'?} & \texttt{WORD / \_word\_} & Loud volume / emphatic stress & \texttt{ye PAAP hai / \_bilkul\_ galat} \\ \hline
\texttt{---} & Abrupt stop, cut-off, or self-repair & \texttt{police ki taraf se--- abhi jaanch...} & \texttt{$^\circ$word$^\circ$} & Softened / whispered delivery & \texttt{$^\circ$ander ki baat ye hai$^\circ$} \\ \hline
\texttt{...} & Trailing off / unfinished sentence & \texttt{aap jaante hain ki kaun log...} & \texttt{[sigh] [cough] [scoff]} & Non-lexical vocalisations & \texttt{[sigh] 500 salon ke baad...} \\ \hline
\texttt{[overlap]} & Simultaneous / overlapping speech & \texttt{[Anjana: suniye---] [Guest: nahin!]} & \texttt{\{\{gesture: ...\}\}} & Hand gesture / facial cue / posture & \texttt{\{\{air-quotes on 'loota'\}\}} \\ \hline
\end{longtable}

\textbf{Methodological Reminder:} A non-verbal cue is not automatically manipulative or ideological---a cough, pause, or rapid cadence may reflect studio audio delay, live time pressure, or a speaker's habitual style. Always check (1) whether the cue clusters with verbal features, and (2) whether the speaker's delivery shifts noticeably across different actors or events (within-outlet baseline).\par\smallskip

\bigskip\hrule\bigskip

\subsection*{Event 1. Ram Mandir donation controversy (Non-Verbal \& Paralinguistic Layer)}

\begin{itemize}
  \item \textbf{Focus:} How do prosody, pauses, trailing sentences, vocalisations (sighs, coughs, scoffs), and body language frame the controversy---as solemn devotional sorrow, administrative scrutiny, or partisan combat?
  \item \textbf{Outlet:} The Wire (Senior Editorial Commentary / Analysis by Arfa Khanum Sherwani / Siddharth Varadarajan team) | \textbf{Date:} June 2021
  \item \textbf{Video title:} Ram Mandir Trust Land Deal Controversy: Serious Questions of Corruption \& Accountability | \href{https://youtu.be/7rUL8fGgRJA}{https://youtu.be/7rUL8fGgRJA} | \textbf{Timestamp(s):} \texttt{[00:00]}--\texttt{[14:28]}
  \item \textbf{Analyst:} Baani Solanki (\texttt{2025114016}) | \textbf{Language:} Hindi / Urdu-inflected Hindustani \& English (transcribed in Roman script)
  \item \textbf{Format / Visual setup:} \texttt{[ ]} TV Panel Debate \texttt{[ ]} TV Bulletin \texttt{[x]} Solo YouTube Explainer \texttt{[x]} Interview / Discussion
  \item \textbf{Primary speaker(s) observed:} Senior Journalist / Presenter at \textit{The Wire} (seated in a quiet, sober newsroom/studio setting, delivering measured editorial analysis with deed overlays and press clips)
\end{itemize}

\subsubsection*{Event-specific non-verbal questions (Event 1)}

\noindent\textbf{1.} \textbf{Vocal \& gestural marking of core labels ("scam", "loota", "paap", "gadabdi", "avyavastha"): When the controversy is named, what non-verbal cues accompany the label? (Look for physical air-quotes, vocal scare-quotes, dramatic pauses before the word, hushed/reverential drops in volume, or shouted stress.)}\par\smallskip
\begin{itemize}
  \item \textbf{Notes:}
\textit{The Wire} avoids both devotional melodrama (\texttt{"paap"}) and tabloid screaming, marking the core labels---\textit{"bhrashtachar"} (corruption), \textit{"ghotala"} (scam), and \textit{"financial impropriety"}---with \textbf{low-pitched, grave, deliberate falling intonation \texttt{$\downarrow$} and measured pauses}:\par\smallskip
  \item At \texttt{[00:45]}, the presenter pauses \texttt{(1.0)} before naming the charge: \textit{"Shri Ram Janmabhoomi Teerth Kshetra Trust ke zameen saude mein (1.0) \_gambhir\_ bhrashtachar ke aarop lage hain$\downarrow$"} (\texttt{[00:52]}), accompanied by a slow, serious head nod \texttt{\{\{gesture: slow grave nod + furrowed brow\}\}}.
  \item Conversely, when quoting the Trust's claim of \textit{"paaradarshita"} (transparency, \texttt{[06:15]}), the presenter uses \textbf{vocal scare-quotes}---raising one eyebrow, pausing \texttt{(0.5)}, and inflecting \textit{"paaradarshita$\uparrow$"} with dry, skeptical rising pitch and a subtle ironic smirk \texttt{\{\{gesture: raised eyebrow + dry half-smirk on 'paaradarshita'\}\}}.
\end{itemize}

\noindent\textbf{2.} \textbf{Differential demeanour toward actors (Temple, Trust leadership, counting clerks, donors, opposition accusers): How do the speaker's voice pitch, facial expression, and posture shift when referring to Prabhu Shri Ram / devotees vs. top figures (Champat Rai, Nripendra Misra, CM Yogi) vs. arrested counting staff vs. opposition leaders (Akhilesh Yadav, Awadhesh Prasad)?}\par\smallskip
\begin{itemize}
  \item \textbf{Notes:}
The presenter maintains a \textbf{sober, secular-institutional demeanour} across all actors, differentiating them through analytical tone rather than theatrical volume shifts:\par\smallskip
  \item \textbf{Donors / Ordinary Devotees:} Referred to with quiet civic sympathy (\texttt{"crore-on aam nagrikon aur bhakton ka paisa"}, \texttt{[01:30]}), opening both hands palm-up on the desk \texttt{\{\{gesture: both palms open resting on desk\}\}} to frame them as public stakeholders cheated by elites.
  \item \textbf{Trust Leadership \& BJP Figures (Champat Rai, Mayor Rishikesh Upadhyay, Anil Mishra):} Addressed with a \textbf{steady, unblinking prosecutorial gaze and clinical precision} (\texttt{[03:20]}--\texttt{[05:45]}). There is zero deference to their religious status; Champat Rai is treated strictly as the chief executive of a public trust created by a Supreme Court mandate (\texttt{[07:10]}).
  \item \textbf{Opposition Accusers \& Local Journalists (Sanjay Singh, Pawan Pandey, Ayodhya press):} Referred to with respectful, collegial nods (\texttt{[02:10]}), crediting them for bringing primary registration deeds into the public domain.
\end{itemize}

\noindent\textbf{3.} \textbf{Solemn devotional grief vs. high-decibel partisan spectacle: How do pauses, sighs, speech rate, and hand gestures construct the segment's emotional register? Where is delivery slow, hushed, and mournful (sighs, bowed head, folded hands), and where does it escalate into rapid-fire shouting, desk-thumping, or mocking irony?}\par\smallskip
\begin{itemize}
  \item \textbf{Notes:}
\textit{The Wire} operates in a \textbf{slow, grave, seminar-room / editorial-board register} (\textasciitilde{}115--125 words per minute) throughout the entire 14-minute video. There is zero shouting, zero desk-thumping, and zero devotional hand-folding. Instead, emotional weight is constructed through \textbf{long, reflective silences \texttt{(1.0)}--\texttt{(1.5)}, weary civic sighs \texttt{[sigh]}, and bitingly understated irony}. For example, at \texttt{[08:20]}, when discussing how ruling-party leaders invoke Lord Ram to shield themselves from audit questions, the speaker lets out a quiet \texttt{[sigh]}, shakes their head slowly, and speaks in a lowered, sorrowful register \texttt{$^\circ$...$^\circ$}: \textit{"$^\circ$Aastha ki aad mein (.) sawaal poochhne walon ko hi deshdrohi ya Ram-virodhi keh diya jaata hai$^\circ$$\downarrow$"} (\texttt{[08:28]}).\par\smallskip
\end{itemize}

\noindent\textbf{4.} \textbf{Non-verbal evaluation of financial scrutiny: When questions about unreceipted offerings, lack of frisking, or land deals are raised, how does the host/speaker react non-verbally (earnest nodding, leaning into the camera, open-palm appeals vs. eye-rolls, scoffs, dismissive hand waves, or incredulous rising pitch $\uparrow$)?}\par\smallskip
\begin{itemize}
  \item \textbf{Notes:}
Financial scrutiny is treated as a \textbf{vital constitutional and public-interest duty}. During the examination of the \texttt{18 March 2021} sale deed vs. agreement to sell (\texttt{[02:40]}--\texttt{[05:30]}), the presenter leans slightly forward toward the camera, holds up the printed registration papers / points to the screen overlay \texttt{\{\{gesture: holding printed document + tapping key clauses with pen\}\}}, and uses deliberate finger-counting (\texttt{[03:15]}) to highlight three unanswered questions: (1) the 825\% markup in 300 seconds, (2) the identical witnesses (Anil Mishra and Mayor Rishikesh Upadhyay), and (3) why the Trust did not buy directly from the original owner Kusum Pathak (\texttt{[04:50]}).\par\smallskip
\end{itemize}

\noindent\textbf{5.} \textbf{Performing certainty vs. hedging and innuendo (pauses, stops, trailing off): Are allegations delivered with definitive falling intonation $\downarrow$, finger-counting, and pointing at on-screen evidence (CCTV, documents), or hedged with hesitant pauses (.), abrupt self-interruptions (---), throat-clearing [cough], and trailing sentences (...) that imply higher-level complicity without naming it?}\par\smallskip
\begin{itemize}
  \item \textbf{Notes:}
\textit{The Wire} combines \textbf{documentary certainty} on the paper trail with \textbf{grave, trailing rhetorical reflections (\texttt{...})} on the political ecosystem:\par\smallskip
  \item \textbf{Documentary Certainty:} The dates, stamp values, and witness signatures are delivered with crisp falling intonation \texttt{$\downarrow$} and rhythmic pen-taps on the desk (\texttt{"Yeh kisi vipakshi party ka aarop nahin, balki Uttar Pradesh sarkar ke registration vibhag ke kagaz hain$\downarrow$"}, \texttt{[03:55]}).
  \item \textbf{Implicating Higher Power via Trailing Off (\texttt{...}):} At \texttt{[10:15]}, when asking how a local BJP Mayor and an RSS-nominated Trustee could execute such a transaction without top-level clearance, the presenter slows down and trails off: \textit{"Bina upar ki manzoori ke (.) Ayodhya jaise high-profile project mein itna bada khel ho jaaye... yeh baat kisi ke gale nahin utarti..."} (\texttt{[10:25]}), accompanied by a slow head shake and raised eyebrows \texttt{\{\{gesture: slow head shake + arched eyebrows\}\}}.
\end{itemize}

\noindent\textbf{6.} \textbf{Turn-taking, interruptions, and acoustic treatment of denials: In debates/interviews: who receives calm, uninterrupted speaking time and sympathetic nods, and who is cut off mid-sentence (---), talked over, or muted? In solo videos: how are clips of official denials framed via reaction pauses, raised eyebrows, or sarcastic chuckles?}\par\smallskip
\begin{itemize}
  \item \textbf{Notes:}
In this editorial commentary format, there is no live shouting match. However, when Champat Rai's brief video clip dismissing the allegations (\texttt{"hum is par koi tippani nahin karenge"}, \texttt{[07:35]}) is played, the presenter follows it with a \textbf{pregnant \texttt{(1.5)} second silence}, looking directly into the camera with a raised eyebrow and a quiet, dry \texttt{[scoff]} (\texttt{[07:48]}): \textit{"Jawab dene ke bajaye (1.5) maun ya phir dhamki$\downarrow$"} (\texttt{[07:52]}). The silence frames the Trust's refusal to answer questions as an admission of guilt.\par\smallskip
\end{itemize}

\noindent\textbf{7.} \textbf{Performing the boundary between religious devotion and institutional critique: How does the presenter use vocal tone, sighs, and body language to signal personal reverence for the Temple (or secular objectivity) while criticising---or defending---the Trust's financial management?}\par\smallskip
\begin{itemize}
  \item \textbf{Notes:}
Unlike Aaj Tak or Republic TV (which perform Hindu devotional insiderhood), \textit{The Wire} explicitly performs \textbf{secular, legal-constitutional objectivity}. The presenter sits with composed, journalistic posture, keeps hands folded calmly on the desk between points \texttt{\{\{gesture: hands lightly clasped on desk in newsroom posture\}\}}, and grounds the critique in legal accountability: because the Shri Ram Janmabhoomi Teerth Kshetra Trust was constituted under a \textbf{Supreme Court judgment} (\texttt{[07:10]}) and receives tax-exempt public donations, its finances are a matter of secular public law, not private religious faith (\texttt{[07:25]}).\par\smallskip
\end{itemize}

\noindent\textbf{8.} \textbf{Self-presentation as the "guardian of public trust": How does the speaker use direct eye-gaze, posture, moral sighs, or emphatic gestures (pointing at the viewer, hands on heart/desk) to position themselves as the authentic protector of devotees or public accountability?}\par\smallskip
\begin{itemize}
  \item \textbf{Notes:}
The presenter positions \textit{The Wire} as the \textbf{uncompromising fourth estate ("speaking truth to power")} at a time when mainstream television has surrendered its watchdog role. Using steady, unhurried eye-gaze into the lens, a weary moral \texttt{[sigh]} over attacks on local Ayodhya journalists who reported the scam (\texttt{[11:40]}), and a calm open-palm appeal (\texttt{[13:50]}), the host projects quiet moral authority rather than theatrical rage.\par\smallskip
\end{itemize}

\subsubsection*{Apply the five non-verbal features (Event 1)}

\begin{longtable}{|p{0.31\textwidth}|p{0.31\textwidth}|p{0.31\textwidth}|}
\hline
Non-verbal feature & Timestamp + spoken anchor words + non-verbal notation & What it does in context (stance, framing, implication) \\ \hline \endhead
\textbf{1. Pauses, silences \& stops} & \texttt{[07:48]} \texttt{[(1.5) silent reaction pause after Champat Rai clip]} \textit{"Jawab dene ke bajaye (1.5) maun ya phir dhamki$\downarrow$"} & The \texttt{(1.5)} silence lets Champat Rai's evasive response hang in the air, framing his refusal to answer as arrogant stonewalling. \\ \hline
\textbf{2. Trailing sentences (\texttt{...})} & \texttt{[10:25]} \textit{"Bina upar ki manzoori ke Ayodhya jaise project mein itna bada khel ho jaaye... yeh baat kisi ke gale nahin utarti..."} & Leaves the identity of \textit{"upar"} (top state/national leadership) unspoken while inviting the viewer to infer high-level political sanction. \\ \hline
\textbf{3. Intonation, pitch \& volume} & \texttt{[00:52]} \textit{"\_Gambhir\_ bhrashtachar ke aarop lage hain$\downarrow$"} (low, grave falling pitch \texttt{$\downarrow$}) vs. \texttt{[06:15]} \textit{"Trust ki 'paaradarshita'$\uparrow$"} (skeptical rising inflection \texttt{$\uparrow$}) & Grave falling intonation lends institutional weight to the corruption charge, while rising inflection on \texttt{'paaradarshita'} turns the Trust's claim into irony. \\ \hline
\textbf{4. Coughs, sighs \& vocalisations} & \texttt{[07:49]} \texttt{[[scoff]]} (soft dry scoff after Trust denial) / \texttt{[08:20]} \texttt{[[sigh]]} \textit{"$^\circ$Aastha ki aad mein sawaal poochhne walon ko hi...$^\circ$"} & The soft \texttt{[scoff]} punctures the Trust's moral posturing; the weary \texttt{[sigh]} performs journalistic sorrow over the weaponisation of faith. \\ \hline
\textbf{5. Hand gestures, face \& posture} & \texttt{[03:15]} \texttt{\{\{gesture: holding printed registry papers + tapping 3 points with pen\}\}} / \texttt{[06:15]} \texttt{\{\{gesture: single raised eyebrow + dry half-smirk\}\}} & Holding physical papers grounds the critique in public records; the raised eyebrow cues skeptical distance from official statements. \\ \hline
\end{longtable}

\subsubsection*{Short summary (Event 1 Non-Verbal Synthesis)}

\begin{itemize}
  \item \textbf{What overall interpretation or emotional stance does the non-verbal delivery encourage?}
The non-verbal delivery encourages \textbf{sober, critical, secular-constitutional scrutiny}. By eschewing both devotional emotionalism and studio shouting in favour of a slow (\textasciitilde{}120 WPM), low-decibel newsroom cadence, long \texttt{(1.5)} reaction silences, and methodical document-tapping, \textit{The Wire} frames the land transaction as a serious abuse of public trust shielded by religious majoritarianism.\par\smallskip
  \item \textbf{Whose voice or physical presence is acoustically/visually foregrounded, undermined, or silenced?}
\textbf{The registration deeds, local Ayodhya investigative journalists, and opposition whistleblowers} are visually and acoustically foregrounded with calm authority. \textbf{Champat Rai and the Trust leadership} are acoustically and visually undermined---their brief denial clips are followed by pregnant \texttt{(1.5)} silences, raised eyebrows, and dry scoffs that expose their refusal to engage with the documentary facts.\par\smallskip
  \item \textbf{What evidence supports your non-verbal reading? What complicates it (e.g., speaker baseline, format)?}
The contrast between the crisp falling intonation \texttt{$\downarrow$} and pen-tapping on registration facts (\texttt{[03:15]}--\texttt{[03:55]}) versus the rising skeptical inflection \texttt{$\uparrow$} on \textit{"paaradarshita"} (\texttt{[06:15]}) and trailing off \texttt{...} on higher political approval (\texttt{[10:25]}) shows how subtle prosody replaces loud volume in serious editorial commentary. A complicating factor is \textbf{speaker baseline}: \textit{The Wire}'s presenters consistently use this grave, understated, slightly mournful cadence across all political broadcasts, meaning the low-key delivery is an institutional house style designed to signal "serious journalism" in contrast to mainstream TV.\par\smallskip
\end{itemize}

\bigskip\hrule\bigskip

\subsection*{Event 2. Students' protest regarding NEET paper leak (Non-Verbal \& Paralinguistic Layer)}

\begin{itemize}
  \item \textbf{Focus:} How do prosody, pauses, trailing sentences, vocalisations, and visual body language construct the legitimacy of student protesters, police use of force, and state accountability?
  \item \textbf{Outlet:} The Wire (Ground Report \& Studio Analysis on NEET Paper Leak \& Student Protests) | \textbf{Date:} June 2024
  \item \textbf{Video title:} NEET Paper Leak Protest: Students March, Police Crackdown \& NTA Accountability | \href{https://youtu.be/pYaTVvHeGoM}{https://youtu.be/pYaTVvHeGoM} | \textbf{Timestamp(s):} \texttt{[00:00]}--\texttt{[16:50]}
  \item \textbf{Analyst:} Baani Solanki (\texttt{2025114016}) | \textbf{Language:} Hindi \& English (transcribed in Roman script)
  \item \textbf{Format / Visual setup:} \texttt{[ ]} TV Panel Debate \texttt{[ ]} TV Bulletin \texttt{[ ]} Solo YouTube Explainer \texttt{[x]} Interview / Discussion
  \item \textbf{Primary speaker(s) observed:} \textit{The Wire} Ground Reporter \& Studio Anchor, alongside prolonged, uninterrupted on-street interviews with protesting NEET Students, Youth Activists, and Parents at Jantar Mantar / Parliament Street
\end{itemize}

\subsubsection*{Event-specific non-verbal questions (Event 2)}

\noindent\textbf{1.} \textbf{Prosodic \& gestural marking of participant labels ("students/bachche" vs. "protesters" vs. "upadravi/goons"): How do vocal tone, volume, facial expression, and hand gestures shift across labels? (e.g., softened, empathetic pitch for "students" vs. harsh, staccato stress, scowls, or accusatory finger-pointing on "upadraviyon", "goons", or "anti-social elements".)}\par\smallskip
\begin{itemize}
  \item \textbf{Notes:}
\textit{The Wire} never uses state-policing labels like \textit{"upadravi"} or \textit{"goons"}. Protesters are consistently referred to as \textit{"chhatra"} (students), \textit{"aabhyarthi"} (aspirants), \textit{"naujawan"} (youth), and \textit{"satyagrahi / pradarshankari"} (\texttt{[00:30]}, \texttt{[04:15]}, \texttt{[11:20]}) with a \textbf{warm, respectful, level vocal tone and open, inviting hand gestures} holding the microphone toward them \texttt{\{\{gesture: reporter extending mic steadily with empathetic nod\}\}}. When quoting ruling-party or media claims that protesters are \textit{"political elements"} (\texttt{[12:40]}), the studio anchor marks the phrase with \textbf{skeptical rising intonation \texttt{$\uparrow$} and raised eyebrows}, visually rejecting the criminalisation of student dissent.\par\smallskip
\end{itemize}

\noindent\textbf{2.} \textbf{Delivery of students' stated demands: When the protesters' core demands (Education Minister's resignation, compensation, release of Sonam Wangchuk, exam reform) are mentioned, are they delivered slowly with deliberate pauses and finger-counting, rushed through in a flat monotone, mocked with air-quotes/scoffs, or trailed off (...)?}\par\smallskip
\begin{itemize}
  \item \textbf{Notes:}
The protesters' core demands---(1) immediate resignation of Education Minister Dharmendra Pradhan, (2) abolition/complete restructuring of the NTA, (3) re-NEET / fair justice for affected candidates, and (4) stopping police repression of peaceful activists (\texttt{[03:10]}--\texttt{[05:45]}, \texttt{[13:15]})---are delivered \textbf{both directly by the students in their own voices and summarised slowly by the host with deliberate \texttt{(0.8)} pauses and finger-counting} \texttt{\{\{gesture: counting demands on fingers with firm nods\}\}}. Far from mocking the demand for the Education Minister's resignation, the host delivers it with \textbf{grave falling intonation \texttt{$\downarrow$}} (\textit{"Chhatron ka seedha sawaal hai---Dharmendra Pradhan istifa kyon nahin dete$\downarrow$"}, \texttt{[13:28]}), treating ministerial resignation as standard democratic accountability.\par\smallskip
\end{itemize}

\noindent\textbf{3.} \textbf{Where does vocal \& kinetic energy peak---grievance, street disruption, or electoral horse-race?: Track where the speaker's volume, pitch range, pace, and gesture intensity reach their climax: during accounts of the paper leak/student suicides, during street clashes and barricade breaches, or during partisan BJP-vs-Opposition sparring.}\par\smallskip
\begin{itemize}
  \item \textbf{Notes:}
Vocal and kinetic energy peaks during \textbf{the students' own raw, emotional testimonies on the street (\texttt{[02:15]}--\texttt{[06:30]}) and during ground footage of police dragging and detaining peaceful protesters (\texttt{[07:40]}--\texttt{[10:15]})}:\par\smallskip
  \item When a female NEET aspirant at the protest site speaks with a trembling voice, rising pitch \texttt{$\uparrow$}, and tears in her eyes (\textit{"Humne \texttt{3} saal apni zindagi ke diye hain, aur badle mein humein lathiyan mil rahi hain$\uparrow$!"}, \texttt{[03:42]}), the camera holds a steady, unflinching eye-level close-up while the reporter listens in complete silence with a solemn, empathetic nod.
\end{itemize}

\noindent\textbf{4.} \textbf{Non-verbal attribution of responsibility (sighs, stops, head shakes, pointing): When examination authorities (NTA), the Education Minister, police, opposition politicians, or students are named as responsible, what non-verbal cues accompany the attribution (exasperated sighs, mocking laughs, abrupt self-stops ---, or direct finger-pointing at the lens)?}\par\smallskip
\begin{itemize}
  \item \textbf{Notes:}
Responsibility is attributed directly to the \textbf{Union Education Minister (Dharmendra Pradhan), the NTA leadership (Subodh Kumar Singh), and the ruling establishment's pattern of institutional capture}:\par\smallskip
  \item When discussing how the NTA chief and Education Minister initially denied any paper leak despite the Patna FIR (\texttt{[11:10]}), the studio anchor lets out an audible, weary \texttt{[sigh]}, shakes their head slowly \texttt{\{\{gesture: slow head shake + furrowed brow\}\}}, and taps the desk with a pen on each contradiction (\texttt{[11:25]}).
  \item When referencing the Godhra (Gujarat) and Hazaribagh/Patna exam centre nexus (\texttt{[12:05]}), the host raises an eyebrow and uses pointed downward hand gestures to highlight the political connections of the centre operators.
\end{itemize}

\noindent\textbf{5.} \textbf{Prosody of questions---open inquiry vs. loaded rhetorical challenge: Are questions about the protest delivered with level, open intonation, or with sharp rising pitch $\uparrow$, sarcastic stress, raised eyebrows, and trailing clauses (...) that presuppose protesters are violent or politically manipulated?}\par\smallskip
\begin{itemize}
  \item \textbf{Notes:}
\textit{The Wire}'s prosody of questions is the \textbf{exact mirror image of Aaj Tak}:\par\smallskip
  \item \textbf{Questions to Protesting Students (\texttt{[02:10]}, \texttt{[04:50]}):} Delivered with \textbf{softened, level, open inquiry intonation} (\texttt{"Aap log itni garmi mein yahan baithe hain, aapki mukhya maangein kya hain$\downarrow$?"}, \texttt{[02:14]}), giving the student full agency to frame the issue.
  \item \textbf{Questions Directed at the State / Police / NTA (\texttt{[08:50]}, \texttt{[13:25]}):} Delivered with \textbf{pointed rising pitch \texttt{$\uparrow$}, emphatic stress, and prosecutorial gravity} (\texttt{"Kya apna haq maangne wale chhatron par lathi chalana aur unhe hirasat mein lena hi sarkar ka jawab hai$\uparrow$?"}, \texttt{[08:55]}).
\end{itemize}

\noindent\textbf{6.} \textbf{Visual \& paralinguistic performance of injury and violence: How are claims of violence substantiated non-verbally on screen? (Compare camera close-ups and pained vocal tone on police bandages/fractured hands vs. visual/acoustic treatment of student head injuries, lathi-charge footage, or defensive stammers/pauses when an officer or guest is asked who used force first.)}\par\smallskip
\begin{itemize}
  \item \textbf{Notes:}
Whereas Aaj Tak substantiates violence by zooming in on a police officer's bandaged hand, \textit{The Wire} substantiates violence through \textbf{eye-level camera close-ups of bruised students, torn clothes, and police personnel grabbing protesters by the limbs and shoving them into buses} (\texttt{[07:45]}--\texttt{[09:30]}). A student protester holds up his scraped arm/wrist to \textit{The Wire}'s camera \texttt{\{\{gesture: student holds bruised wrist up to camera lens\}\}} (\texttt{[08:12]}) while breathing heavily \texttt{[gasp]} and saying: \textit{"Dekhiye police ne kaise ghasita hai$\downarrow$"} (\texttt{[08:15]}). The reporter lowers their voice in solidarity (\texttt{$^\circ$aapko chot aayi hai$^\circ$}, \texttt{[08:18]}), visually and acoustically validating student injury at the hands of state force.\par\smallskip
\end{itemize}

\noindent\textbf{7.} \textbf{Asymmetric delivery of police actions vs. protester actions: Compare the vocal cadence, pauses, and hand gestures used when describing police actions (e.g., calm, matter-of-fact administrative tone on "tear gas chhode gae" / "preventive arrangement") versus protester actions (alarmist pitch, striking gestures, or ironic softening as "retaliation").}\par\smallskip
\begin{itemize}
  \item \textbf{Notes:}
  \item \textbf{Police Actions:} Described in an \textbf{indignant, critical cadence} that strips away bureaucratic euphemisms: when quoting the police claim of \textit{"preventive detention"} or \textit{"law and order arrangement"} (\texttt{[09:10]}), the presenter uses a \textbf{satirical, dry inflection and physical air-quotes} \texttt{\{\{gesture: air-quotes on 'preventive detention'\}\}}, contrasting the euphemism with the visual footage of students being dragged across the asphalt.
  \item \textbf{Protester Actions:} Described in a \textbf{dignified, empathetic cadence} (\texttt{"shantipurna tarike se Sansad gherao karne ja rahe chhatra"}, \texttt{[07:50]}), framing their march as a constitutional right to peaceful assembly.
\end{itemize}

\noindent\textbf{8.} \textbf{Acoustic \& visual space granted to students' own voices: Are protesting students heard in their own voice with clear audio, steady close-ups, and empathetic listening nods, or are they reduced to noisy background chaos B-roll while studio anchors/politicians speak over them?}\par\smallskip
\begin{itemize}
  \item \textbf{Notes:}
Protesting students receive \textbf{over 60\% of the total speaking time in the segment} with \textbf{clear, unedited close-up audio and zero anchor \texttt{[overlap]} interruptions}. During each student interview (\texttt{[02:15]}--\texttt{[06:30]}, \texttt{[08:05]}--\texttt{[10:00]}), the reporter holds the microphone steadily, maintains eye contact, nods empathetically \texttt{\{\{gesture: continuous empathetic listening nods\}\}}, and waits a full \texttt{(0.8)}--\texttt{(1.0)} second after the student finishes speaking before asking a follow-up question---the exact opposite of television panel cross-talk.\par\smallskip
\end{itemize}

\noindent\textbf{9.} \textbf{Non-verbal shift between the NEET paper-leak issue and the street protest: Does the speaker's non-verbal demeanour shift from sympathetic concern (sombre tone, sighs, solidarity gestures) when discussing the exam leak itself to combative, authoritarian, or dismissive delivery once students march toward Parliament?}\par\smallskip
\begin{itemize}
  \item \textbf{Notes:}
\textbf{No authoritarian shift occurs toward the students.} Instead, as the video moves from the exam paper leak (\texttt{[00:00]}--\texttt{[02:00]}) to the street march and police crackdown (\texttt{[07:40]}--\texttt{[10:30]}), the presenter's non-verbal demeanour shifts from \textbf{sombre sympathy for exam victims to grave constitutional outrage against police high-handedness}, treating the street protest as the natural democratic extension of the exam grievance.\par\smallskip
\end{itemize}

\subsubsection*{Apply the five non-verbal features (Event 2)}

\begin{longtable}{|p{0.31\textwidth}|p{0.31\textwidth}|p{0.31\textwidth}|}
\hline
Non-verbal feature & Timestamp + spoken anchor words + non-verbal notation & What it does in context (stance, framing, implication) \\ \hline \endhead
\textbf{1. Pauses, silences \& stops} & \texttt{[03:48]} \texttt{[(1.0) respectful silence by reporter after crying student's testimony]} / \texttt{[11:20]} \textit{"Pehle kaha koi leak nahin hua--- (1.2) aur ab CBI jaanch$\downarrow$"} & The \texttt{(1.0)} silence honours the student's raw emotion without studio intrusion; the \texttt{(1.2)} pause highlights the government's U-turn on the paper leak. \\ \hline
\textbf{2. Trailing sentences (\texttt{...})} & \texttt{[14:15]} \textit{"Jab har badi pariksha ka paper leak ho raha ho, aur sawaal poochhne par chhatron ko lathiyan milein... toh is desh ke yuvaon ka bhavishya..."} & Trails off with a slow, sorrowful head shake, inviting the viewer to reflect on authoritarian intolerance toward youth dissent. \\ \hline
\textbf{3. Intonation, pitch \& volume} & \texttt{[02:14]} \textit{"$^\circ$Aapki mukhya maangein kya hain$^\circ$$\downarrow$?"} (gentle, open inquiry to student) vs. \texttt{[08:55]} \textit{"Kya chhatron par \_LATHI\_ chalana hi sarkar ka jawab hai$\uparrow$?!"} & Gentle, low-volume open prosody empowers student speakers, whereas sharp rising pitch \texttt{$\uparrow$} and stress on \texttt{\_LATHI\_} indict state repression. \\ \hline
\textbf{4. Coughs, sighs \& vocalisations} & \texttt{[08:12]} Student: \texttt{[[gasp]]} \textit{"Dekhiye police ne kaise ghasita hai$\downarrow$"} / \texttt{[11:10]} Anchor: \texttt{[[sigh]]} (on Education Minister's initial denial) & Including the student's breathless \texttt{[gasp]} on camera conveys authentic physical distress; the anchor's \texttt{[sigh]} performs civic exhaustion with official denialism. \\ \hline
\textbf{5. Hand gestures, face \& posture} & \texttt{[08:12]} \texttt{\{\{gesture: student holds bruised wrist up to camera lens\}\}} / \texttt{[09:10]} \texttt{\{\{gesture: anchor air-quotes on 'preventive detention'\}\}} / \texttt{[02:20]} \texttt{\{\{gesture: reporter's steady mic hold + empathetic nods\}\}} & The student's bruised wrist directly counters TV visuals of police bandages; air-quotes debunk police euphemisms; empathetic nods validate student testimony. \\ \hline
\end{longtable}

\subsubsection*{Short summary (Event 2 Non-Verbal Synthesis)}

\begin{itemize}
  \item \textbf{What overall interpretation or emotional stance does the non-verbal delivery encourage?}
The non-verbal delivery encourages \textbf{solidarity with protesting students and moral indignation at both the NTA's exam failure and the police crackdown on peaceful marchers}. Through eye-level close-ups, unmuted student audio, and grave studio commentary, the viewer is guided to see the protesters as courageous citizens defending public education against an unaccountable state.\par\smallskip
  \item \textbf{Whose voice or physical presence is acoustically/visually foregrounded, undermined, or silenced?}
\textbf{Protesting NEET aspirants and youth activists} are acoustically and visually foregrounded, receiving uninterrupted close-ups, empathetic nods, and visual validation of their bruises (\texttt{[08:12]}). \textbf{Police claims of "preventive arrangement" and the Education Minister's denials} are undermined via air-quotes, weary \texttt{[sigh]} vocalisations, and direct juxtaposition with crackdown footage.\par\smallskip
  \item \textbf{What evidence supports your non-verbal reading? What complicates it (e.g., speaker baseline, format)?}
The zero-interruption turn-taking structure during street interviews (\texttt{[02:15]}--\texttt{[06:30]}) and the visual close-up on a student's scraped wrist (\texttt{[08:12]}) provide concrete multimodal evidence of a citizen-centred frame. What complicates a purely neutral reading is that \textit{The Wire}'s ground-report format selectively foregrounds articulate, injured protesters while granting minimal screen time to police spokespersons, reflecting its counter-establishment editorial baseline.\par\smallskip
  \item \textbf{How do the non-verbal cues interact with the spoken words? (Do they reinforce, dramatise, soften/hedge, invert meaning via sarcasm/irony, or leave an accusation unspoken via trailing off?)}
The non-verbal cues \textbf{reinforce student testimonies} through silent, respectful listening nods and \texttt{(1.0)} pauses, \textbf{invert state policing euphemisms via irony} (air-quotes and skeptical pitch on \textit{"'preventive detention'"} at \texttt{[09:10]}), and \textbf{leave a broader indictment of democratic backsliding unspoken via trailing off} (\texttt{...} at \texttt{[14:15]}).\par\smallskip
  \item \textbf{What context, speaker baseline, or format factors support or complicate this interpretation? (e.g., Is this the speaker's normal style across all videos, or a marked shift on this specific actor/claim?)}
Across both Event 1 and Event 2, \textit{The Wire} maintains a consistent \textbf{low-decibel, grave, civil-liberties baseline} (\textasciitilde{}120 WPM, minimal shouting, heavy reliance on pauses \texttt{(1.5)}, weary \texttt{[sigh]} vocalisations, and skeptical eyebrow raises). However, Event 2 incorporates high-arousal ground footage of breathless \texttt{[gasp]}, crying students (\texttt{[03:42]}, \texttt{[08:12]}), creating a powerful multimodal contrast between raw street distress and calm studio indictment.\par\smallskip
  \item \textbf{How does another outlet perform the same moment, actor, or claim non-verbally?}
Compare \textbf{The Wire} (\texttt{Event 2}, \texttt{[08:12]}) with \textbf{Aaj Tak} (\texttt{Event 2}, \texttt{[08:40]}): at the exact same protest, Aaj Tak puts an Additional DCP on screen holding up his \textbf{bandaged hand} while the anchor lowers her voice in sympathy (\texttt{$^\circ$DCP sahab ke haath mein fracture hai$^\circ$}) and shouts over students \texttt{[overlap]} as \textit{"upadravi"}. \textit{The Wire} shows the exact inverse visual prop---a \textbf{student holding up his bruised wrist} after being dragged into a police bus (\texttt{[08:12]})---while the anchor uses air-quotes to mock police claims of \textit{"preventive arrangement"}.\par\smallskip
  \item \textbf{My interpretation, with reasons:}
\textit{The Wire}'s multimodal discourse across Events 1 and 2 exemplifies \textbf{counter-hegemonic journalistic framing}. In Event 1, it uses secular newsroom posture, document-tapping, and pregnant \texttt{(1.5)} silences to strip the Ram Mandir Trust of its religious immunity and demand a public audit. In Event 2, it uses acoustic generosity (uninterrupted mic time, empathetic nods) and visual counter-evidence (student bruises vs. police bandages) to legitimise street protest as an essential democratic right.\par\smallskip
\end{itemize}

\end{document}
